\documentclass[11pt]{article}

\usepackage[a4paper,margin=1in]{geometry}
\usepackage{amsmath,amssymb,amsthm,mathtools}
\usepackage[hidelinks]{hyperref}
\hypersetup{
  pdfauthor={Bangyen Pham},
  pdftitle={Sharpness of the Coefficient Order-Statistic Bound},
  pdfsubject={Sharpness and necessity results for the coefficient order-statistic bound on polynomial multiples},
  pdfkeywords={polynomial multiples, coefficient mass, reduced height, Chebyshev systems, exponential sums},
  pdflang={en-US}
}

\newtheorem{theorem}{Theorem}[section]
\newtheorem{lemma}[theorem]{Lemma}
\newtheorem{corollary}[theorem]{Corollary}
\newtheorem{proposition}[theorem]{Proposition}
\theoremstyle{definition}
\newtheorem{conjecture}[theorem]{Conjecture}
\theoremstyle{plain}
\newcommand{\R}{\mathbb R}
\newcommand{\N}{\mathbb N}
\newcommand{\E}{\mathbb E}
\DeclareMathOperator{\Tail}{Tail}
\DeclareMathOperator{\sgn}{sgn}

\title{Sharpness of the Coefficient Order-Statistic Bound}
\author{\shortstack{Bangyen Pham\\
\small Independent Researcher\\
\small\href{mailto:bangyenp@gmail.com}{\texttt{bangyenp@gmail.com}}\\
\small\href{https://orcid.org/0009-0006-9928-2088}{ORCID: 0009-0006-9928-2088}}}
\date{September 2026}

\begin{document}
\maketitle

\begin{abstract}
For real roots $2\le r_1\le\cdots\le r_L$ and a nonzero multiple $F$ of
$P=\prod_{i=1}^L(x-r_i)$ with leading coefficient $f_D$, the companion
paper~\cite{coefficient-mass} proves that the $(u+1)$-st largest nonleading
coefficient magnitude satisfies $b_{u+1}\ge|f_D|\prod_{i>u}(r_i-1)$.  We show
that this bound is the infimum over monic multiples exactly when every
partial sum of the leading coefficients of $P_u=\prod_{i>u}(x-r_i)$ is at
most $|P_u(1)|$ in absolute value, repeated roots included.  For all real
roots above $1$ we express the infimum as $1/T$, where $T$ is the value of a
game over the positions of the $u$ coefficients that $b_{u+1}$ ignores,
played with $\ell^1$ tails of exponential sums on the reciprocal roots.
When these positions escape to the bottom the value is $\tau^*$, with
$1/\tau^*=\min_M\max_{j\ge1}|S_j(P_uM)|$ over monic $M$, where $S_j$ is the
sum of the $j+1$ leading coefficients.  For roots at least $2$ we prove
$T=\tau^*$, so the infimum is $1/\tau^*$.  Below $2$ this fails: at
$(\tfrac{11}{10},3,5,7)$ with $u=1$ a reduction to finitely many positions
and an exact computation give the infimum $\tfrac{3398808}{96935}$, below
$1/\tau^*$.  Along $(r,3,5,7)$ with $u=1$ the infimum tends to $24$ as
$r\to1^+$, exceeding it by a quantity of exact order
$(r-1)^{\alpha}$, $\alpha=\log(5/3)/\log3$, whose ratio to $(r-1)^\alpha$
has no limit: it oscillates log-periodically between $80\cdot18^{-\alpha}$
and $\tfrac{400}3\cdot18^{-\alpha}$, and we determine it uniformly,
across the points where the worst position jumps.  The infimum is never
attained.
\end{abstract}

\noindent\textbf{Keywords.} Polynomial multiples; coefficient mass; reduced
height; Chebyshev systems; exponential sums.

\noindent\textbf{2020 Mathematics Subject Classification.} Primary 11C08; Secondary 26C05, 41A50.

\section{Introduction}

The companion paper~\cite{coefficient-mass} proves that, for a multiple of a
polynomial with real roots $r_i\ge2$, the $(u+1)$-st largest nonleading
coefficient magnitude is at least the leading coefficient magnitude times
$\prod_{i>u}(r_i-1)$, whatever the degree.
Precisely, for real roots $2\le r_1\le\cdots\le r_L$, repetitions allowed, and $F$ a
nonzero multiple of $\prod_{i=1}^L(x-r_i)$ of degree $D\ge L$ with leading
coefficient $f_D$, its main bound~\cite[Theorem~3.2]{coefficient-mass} reads,
for every $0\le u\le L-1$,
\begin{equation}\label{eq:order}
  b_{u+1}\ge|f_D|\prod_{i=u+1}^L(r_i-1),
\end{equation}
where $b_k$ is the $k$th largest of the nonleading coefficient magnitudes
$|f_j|$, $j<D$.  The engine behind \eqref{eq:order}
is~\cite[Theorem~2.2]{coefficient-mass}, a
comparison between the infinite $\ell^1$ tail of a normalized exponential
sum with prescribed zeros and the tail of the sum with the same zero count
but a consecutive zero set: for nodes $0<y_1<\cdots<y_c\le1/2$ and a zero
set $Z$ of $c-1$ positive integers with leading run
$\ell=\max\{0\le k\le c-1:\{1,\ldots,k\}\subseteq Z\}$, the unique exponential
sum $\xi$ on those nodes with $\xi_0=1$ and $\xi_z=0$ ($z\in Z$) satisfies
\[
  \Tail(\xi)=\sum_{d\ge1}|\xi_d|\le\prod_{i=1}^{\ell+1}\frac{y_i}{1-y_i},
\]
the product over the $\ell+1$ smallest nodes, with equality exactly at $Z=\{1,\ldots,c-1\}$.
The consecutive case is evaluated in~\cite[Lemma~2.1]{coefficient-mass},
and~\cite[Lemma~3.1]{coefficient-mass} converts the tail bound into a
coefficient bound.

\paragraph{Main results.}  Write $P=\prod_{i=1}^L(x-r_i)$, let
$P_u=\prod_{i>u}(x-r_i)$ be the factor of the $m=L-u$ surviving roots, and let
$S_j(P_u)$ be the sum of its $j+1$ leading coefficients, those of
$x^m,\ldots,x^{m-j}$.  Thus $S_0(P_u)=1$, and
$S_j(P_u)=P_u(1)=(-1)^m\prod_{i>u}(r_i-1)$ for $j\ge m$.  All infima below
are of $b_{u+1}(QP)$ over monic $Q\in\R[x]$.  The \emph{exempt}
coefficients are the $u$ largest nonleading ones, which $b_{u+1}$ ignores.

\emph{The criterion.}  For $2\le r_1\le\cdots\le r_L$, repetitions allowed,
\eqref{eq:order} is the infimum if and only if
\[
  |S_j(P_u)|\le|P_u(1)|\qquad(1\le j<m),
\]
which is condition \eqref{eq:partial} below (Corollary~\ref{cor:charrep}).
For distinct roots the multiples approaching the bound carry a short prefix
$S_1(P_u),\ldots,S_{m-1}(P_u)$, a long run of the value $P_u(1)$, and $u$
free coefficients at the bottom, which $b_{u+1}$ discards.

\emph{The value.}  For $1<r_1\le\cdots\le r_L$ the infimum is $1/T$, where
$T=\sup_{|X|=u}\tau(X)$ is the largest cost $\tau(X)$ of a \emph{placement}
$X\subset\N_{>0}$, $|X|=u$, the distances of the exempt coefficients below
the leading one (Theorem~\ref{thm:valuerep}).  Here a \emph{certificate} for
$X$ is an exponential sum $v_d=\sum_i\lambda_ir_i^{-d}$ with $v_0=1$ and
$v|_X=0$ (for repeated roots, an element of the Hermite space of the roots),
and $\tau(X)$ is the least tail $\sum_{d\ge1}|v_d|$ of a certificate for $X$.
As the exempt
coefficients escape to the bottom the cost tends to $\tau^*$, the least tail
of a sum with $v_0=1$ on the surviving roots alone, so $T\ge\tau^*$.  This
$\tau^*$ is a minimum, and
\[
  \frac1{\tau^*}=\min_M\max_{j\ge1}\bigl|S_j(P_uM)\bigr|
\]
over monic $M\in\R[x]$ (Proposition~\ref{prop:attainrep}).  No monic
multiple attains the infimum (Theorem~\ref{thm:noattain}).

\emph{Roots at least $2$.}  For $2\le r_1\le\cdots\le r_L$ the escaping
placement is the worst one, that is, $T=\tau^*$, and the multiplier formula
computes the infimum $1/\tau^*$ (Corollary~\ref{cor:extremalrep}).

\emph{Roots below $2$.}  Here the escaping placement need not be the worst,
but for one exempt coefficient $T$ is still the larger of $\tau^*$ and the
costs of finitely many placements.
For distinct roots above $1$ the supremum $T$ is a maximum, achieved by a
placement or at a level where some exempt coefficients have escaped
(Proposition~\ref{prop:attainT}).  For $u=1$ an exact identity for the cost
of a zero placed above those of an escaping minimiser shows that every
placement beyond an explicit $N$ costs less than $\tau^*$; hence $T$ is the
larger of $\tau^*$ and the finitely many $\tau(\{x\})$ with $x<N$
(Theorem~\ref{thm:halfmass}).  At $(\tfrac{11}{10},3,5,7)$ with $u=1$ this
gives $N=11$, the worst placement $\{5\}$ and the infimum
$\tfrac{3398808}{96935}=35.06\ldots$, against $1/\tau^*=\tfrac{1704}{31}$
(Theorem~\ref{thm:elevenvalue}).  Along $(r,3,5,7)$ with $1<r<3$ and
$u=1$, the escaping placement is the worst one exactly when
$930r^2-7r-1988\ge0$, that is, $r\ge r_c=1.4658\ldots$; for
$\tfrac{13}{10}\le r\le r_c$ the infimum is $48(r-1)(15r+71)/(49r-34)$
(Theorem~\ref{thm:family}).  Below $\tfrac{13}{10}$ the zero set minimising the
cost of the worst placement keeps moving, and the worst placement changes
from $\{4\}$ to $\{5\}$; on $[1.0745\ldots,\tfrac{13}{10}]$ the infimum is
a piecewise rational function of $r$ with seven pieces, each proved by a
multiplier and explicit certificates (Theorem~\ref{thm:familylow}).  As
$r\to1^+$ the pieces accumulate: the infimum tends to $24$, the value at
$(3,5,7)$, and exceeds it by a quantity of exact order $(r-1)^\alpha$,
$\alpha=\log(5/3)/\log3$, while the worst placement is
$\log_3\frac1{r-1}+O(1)$ (Theorem~\ref{thm:familyone}); the costs that
matter are those at the roots $3,5,7$, in closed form
(Lemma~\ref{lem:threenode}), and a witness with one sign flip far out
transfers them to $r$ near $1$.  The ratio of the excess to $(r-1)^\alpha$
has no limit: off neighbourhoods of the points where $3^x(r-1)=18$ for an
integer $x$, the worst placement is the largest $x^*$ with
$3^{x^*}(r-1)<18$, its minimising zero set is $\{1,x^*,z\}$ for some
$z\ge x^*+3$, and the excess is $(80+o(1))(\tfrac35)^{x^*}$, so the ratio
oscillates log-periodically between $80\cdot18^{-\alpha}$ and
$\tfrac{400}3\cdot18^{-\alpha}$ (Theorem~\ref{thm:familyosc}), on an
explicit range of $r$ (Proposition~\ref{prop:familyrdelta}).  In fact every
cost along the family is explicit: for each placement $\{x\}$, $x\ge2$,
and each $1<r<3$ the minimising zero set is $\{1,x,z\}$ or $\{1,z,x\}$,
with $z$ a half-mass point of an explicit exponential sum
(Theorem~\ref{thm:familyzero}).  As $3^x(r-1)\to\kappa$ the cost of
$\{x\}$ tends to an explicit function of $\kappa$, equal to
$\tfrac1{24}$ up to $\kappa=18$ and beyond it the least of the limits
$t+\ell/\kappa$ of finitely many certificates at the roots $3,5,7$ lifted
to $r$, with breakpoints accumulating at $18$
(Theorem~\ref{thm:familylimit}).  Across a jump this profile decides
between $x^*$ and $x^*+1$: as $r$ decreases to a jump the ratio falls by
the factor $\tfrac35$ within a window of width of order
$(r-1)^{\alpha/(1+\alpha)}$ ending at the jump, and off that window it
converges uniformly to the log-periodic profile
(Theorem~\ref{thm:familyjump}).  These
proofs end in a finite exact computation, over $\mathbb Q$ and over
$\mathbb Q(r)$, described in the paragraph on computer assistance before the
appendix.  Along $(r,5,7)$ the
escaping placement is not the worst for any $1<r<2$, so the cutoff $2$
cannot be lowered uniformly (Proposition~\ref{prop:cutoffsharp}).

\emph{The tail comparison.}  \cite[Theorem~2.2]{coefficient-mass} extends
to confluent nodes, strictly unless the zero set is consecutive
(Proposition~\ref{prop:confluent}); the criterion needs this at repeated
roots.

\paragraph{Organisation.}  Section~\ref{sec:extend} proves the confluent
tail comparison.  Section~\ref{sec:criterion} proves the partial-sum
criterion: sufficiency by an explicit construction, necessity by a
certificate argument built on a Descartes rule at infinity.
Section~\ref{sec:value} expresses the infimum as a placement game, proves
that $\tau^*$ is a minimum, with a multiplier witness, completes the criterion
at repeated roots, and shows that the infimum is never attained.
Section~\ref{sec:extremal} proves that at roots at least $2$ the escaping
placement is extremal, by lifting minimisers one node at a time, and shows
that below $2$ it need not be.  Section~\ref{sec:belowtwo} drops the
cutoff: that $T$ is a maximum and the reduction for one exempt coefficient hold
for all distinct roots above $1$, followed by exact examples below $2$.  Section~\ref{sec:scope} closes with a remark
on what is left open and the relation to the literature.

\section{The confluent tail comparison}\label{sec:extend}

The necessity half of the criterion at repeated roots needs~\cite[Theorem~2.2]{coefficient-mass} at
confluent nodes, where the exponential sums are replaced by elements of a
Hermite space, together with its strict form.

\begin{proposition}[Confluent analogue]\label{prop:confluent}
Let $0<y_1<\cdots<y_c\le1/2$ and $m_1,\ldots,m_c\ge1$, put $m=m_1+\cdots+m_c$,
and let $y_{(1)}\le\cdots\le y_{(m)}$ be the multiset with each $y_i$ repeated
$m_i$ times.  Let
\[
  \xi_d=\sum_{i=1}^c p_i(d)y_i^d,\qquad \deg p_i<m_i,
\]
satisfy $\xi_0=1$ and $\xi|_Z=0$ for a set $Z$ of $m-1$ positive integers with
leading run $\ell=\max\{0\le k\le m-1:\{1,\ldots,k\}\subseteq Z\}$.  Then
\[
  \Tail(\xi)=\sum_{d\ge1}|\xi_d|\le\prod_{j=1}^{\ell+1}\frac{y_{(j)}}{1-y_{(j)}}.
\]
\end{proposition}

\begin{proof}
The interpolation system is indexed by the Hermite basis, not by the
ordinary exponentials, and the two must be kept apart.  For
$\varepsilon>0$, $i=1,\ldots,c$ and $k=0,\ldots,m_i-1$ put
\[
  \phi_{i,k,\varepsilon}(d)
  =y_i^d\Bigl(\frac{1-\exp(-\varepsilon d)}{\varepsilon}\Bigr)^k .
\]
For fixed $\varepsilon>0$, the binomial expansion gives
\[
  \phi_{i,k,\varepsilon}(d)
  =\frac1{\varepsilon^k}\sum_{j=0}^{k}(-1)^j\binom{k}{j}
    \bigl(y_i\exp(-j\varepsilon)\bigr)^d,
\]
so each $\phi_{i,k,\varepsilon}$ is a finite linear combination of the
ordinary exponentials $d\mapsto z^d$ at the distinct nodes
$\zeta_{i,j}=y_i\exp(-j\varepsilon)$.  For all sufficiently small $\varepsilon>0$ the $m$
nodes $\zeta_{i,j}$ are distinct, lie in $(0,1/2]$, and stay separated between
distinct $i$.  Within one group $i$ the change of basis from
$\{\zeta_{i,j}^{d}\}_{j<m_i}$ to
$\{\phi_{i,k,\varepsilon}\}_{k<m_i}$ is triangular with diagonal entries
$(-1)^k\varepsilon^{-k}$, hence invertible for $\varepsilon>0$, and the same
function space is spanned either way.

For such an $\varepsilon$ let $\xi^{(\varepsilon)}$ be the unique exponential
sum on the nodes $\zeta_{i,j}$ with $\xi^{(\varepsilon)}_0=1$ and
$\xi^{(\varepsilon)}|_Z=0$; existence and uniqueness are the
Vandermonde argument opening~\cite[Section~2]{coefficient-mass}.
The comparison~\cite[Theorem~2.2]{coefficient-mass} applies with the same $\ell$ and gives
\[
  \Tail(\xi^{(\varepsilon)})
  \le\prod_{j=1}^{\ell+1}\frac{y_{(j)}(\varepsilon)}{1-y_{(j)}(\varepsilon)},
\]
where $y_{(1)}(\varepsilon)\le\cdots\le y_{(m)}(\varepsilon)$ is the sorted
perturbed multiset.

It remains to pass to the limit.  Write
$a^{(\varepsilon)}=(a^{(\varepsilon)}_{i,k})$ for the coefficients of
$\xi^{(\varepsilon)}$ in the $\phi$-basis; they solve
$M(\varepsilon)a^{(\varepsilon)}=\mathbf e_0$, where $M(\varepsilon)$ is the
$m\times m$ matrix with rows indexed by $d\in\{0\}\cup Z$ and columns by
$(i,k)$, with entry $\phi_{i,k,\varepsilon}(d)$, and $\mathbf e_0$ is the first
standard basis vector.  Each entry
$\phi_{i,k,\varepsilon}(d)=y_i^d\bigl((1-\exp(-\varepsilon d))/\varepsilon\bigr)^k$
has a removable singularity at $\varepsilon=0$ and extends analytically there
with limit $d^ky_i^d$, so $M(\varepsilon)\to M(0)$ as $\varepsilon\to0$, where
$M(0)$ is the confluent matrix.  The functions
$d\mapsto d^ky_i^d$ form an extended Chebyshev system on $\R$: a nonzero
$\sum_ip_i(d)y_i^d$ with $\deg p_i<m_i$ has at most $m-1$ real
zeros~\cite[Part~Five, Ch.~1]{polya-szego}.  The
interpolation here is at the $m$ \emph{distinct} points $\{0\}\cup Z$---the
confluence sits in the basis, not in the nodes---so $M(0)$ is invertible by
ordinary unisolvence of a Chebyshev system at distinct
points~\cite[Chapter~I, Theorem~4.1]{karlin-studden}.  Matrix inversion is
continuous on the open set of invertible matrices, so
$a^{(\varepsilon)}\to M(0)^{-1}\mathbf e_0$, and for every fixed $d$,
\[
  \xi^{(\varepsilon)}_d
  =\sum_{i,k}a^{(\varepsilon)}_{i,k}\phi_{i,k,\varepsilon}(d)
  \longrightarrow\sum_{i,k}a_{i,k}d^ky_i^d=\xi_d .
\]
Fatou's lemma for series, applied to the nonnegative terms
$|\xi^{(\varepsilon)}_d|$, gives
\[
  \Tail(\xi)=\sum_{d\ge1}|\xi_d|
  \le\liminf_{\varepsilon\to0}\Tail(\xi^{(\varepsilon)}).
\]
Finally the order statistics of a converging multiset converge, so the
product $\prod_{j=1}^{\ell+1}y_{(j)}(\varepsilon)/(1-y_{(j)}(\varepsilon))$
tends to $\prod_{j=1}^{\ell+1}y_{(j)}/(1-y_{(j)})$.
\end{proof}

\paragraph{The displaced step.}  Lemmas~\ref{lem:confstrict}
and~\ref{lem:delete} below use one step of the proof
of~\cite[Theorem~2.2]{coefficient-mass}, which we restate.  Let $c\ge2$ and
$0<y_1<\cdots<y_c\le1/2$, let $Z\ne\{1,\ldots,c-1\}$ be a set of $c-1$
positive integers, so that its leading run is below $c-1$ as that step
assumes, $z=\max Z$ and $Z'=Z\setminus\{z\}$, and let $w$ be the sum on all $c$
nodes with $w_0=1$ and $w|_Z=0$.  Let $V$ be the sum on
$y_1,\ldots,y_{c-1}$, all nodes but the largest, with $V_0=1$ and
$V|_{Z'}=0$, and $G$ the sum on all $c$ nodes with
$G|_{\{0\}\cup Z'}=0$ and $G_z=1$; then $w=V-V_zG$, and $V_z\ne0$ by the
zero bound.  Put $\mu=|V_z|$ and, for $X\in\{V,G\}$,
\[
  \Gamma_X(z)=\frac{\sum_{d\ge z}|X_d|}{|X_z|},
  \qquad \Sigma_G=\sum_{1\le d<z}|G_d| .
\]
The displaced step of the proof of~\cite[Theorem~2.2]{coefficient-mass}
proves the splitting identity and the correction-ratio bound
\[
  \Tail(V)-\Tail(w)=\mu\bigl(2\Gamma_V(z)-\Gamma_G(z)+\Sigma_G\bigr),
  \qquad
  \Gamma_G(z)\le\frac{\Gamma_V(z)}{1-y_c}\le2\Gamma_V(z),
\]
the identity from the zero bound and the distinctness of the nodes, and the
bound from~\cite[Lemma~2.3]{coefficient-mass}, proved by weak alternation,
and, in its last inequality, $y_c\le1/2$; neither uses the induction
hypothesis of that proof.
Together they give
\[
  \Tail(w)\le\Tail(V)-\mu\Sigma_G ,
\]
and the same step shows $\Sigma_G>0$.

\begin{lemma}[Strictness at confluent nodes]\label{lem:confstrict}
In Proposition~\ref{prop:confluent} the inequality is strict unless
$Z=\{1,\ldots,m-1\}$.
\end{lemma}

\begin{proof}
Suppose $Z\ne\{1,\ldots,m-1\}$.  Then $\ell\le m-2$, the largest element
$z=\max Z$ satisfies $z\ge m$, and $Z'=Z\setminus\{z\}$ has the same leading
run $\ell$.  Keep the perturbed nodes $\zeta_{i,j}=y_i\exp(-j\varepsilon)$ of the
proof of Proposition~\ref{prop:confluent}; the largest of them is
$y_c$ itself ($j=0$), and it is strictly larger than all the others.  Run the
displaced step above once at these $m$ distinct nodes, all at most $1/2$:
$V^{(\varepsilon)}$ is the sum on the perturbed nodes other than $y_c$ with
$V^{(\varepsilon)}_0=1$ and $V^{(\varepsilon)}|_{Z'}=0$, $G^{(\varepsilon)}$
the sum on all of them with $G^{(\varepsilon)}|_{\{0\}\cup Z'}=0$ and
$G^{(\varepsilon)}_z=1$, and $\mu_\varepsilon=|V^{(\varepsilon)}_z|$, so
\[
  \Tail(\xi^{(\varepsilon)})
  \le\Tail(V^{(\varepsilon)})-\mu_\varepsilon \Sigma_{G^{(\varepsilon)}},
  \qquad
  \Sigma_{G^{(\varepsilon)}}=\sum_{1\le d<z}|G^{(\varepsilon)}_d| .
\]
The comparison~\cite[Theorem~2.2]{coefficient-mass} on the $m-1$ nodes of $V^{(\varepsilon)}$ bounds
$\Tail(V^{(\varepsilon)})$ by the product over their $\ell+1\le m-1$ smallest
members, which are the $\ell+1$ smallest of all $m$ perturbed nodes because the
one removed is the largest.  So
$\Tail(\xi^{(\varepsilon)})\le\prod_{j\le \ell+1}y_{(j)}(\varepsilon)/(1-y_{(j)}(\varepsilon))
-\mu_\varepsilon \Sigma_{G^{(\varepsilon)}}$.

Now let $\varepsilon\to0$.  The sums $V^{(\varepsilon)}$ and
$G^{(\varepsilon)}$ are interpolants at the distinct points $\{0\}\cup Z'$
and $\{0\}\cup Z$ in spaces spanned by functions of the form
$\phi_{i,k,\varepsilon}$ (for the group of $y_c$, which has lost its top
member, with $y_c$ replaced by $y_c\exp(-\varepsilon)$), so exactly as in the
proof of Proposition~\ref{prop:confluent} they converge pointwise to the
confluent interpolants: $V^{(0)}$ in the Hermite space of the multiset with one
copy of $y_c$ removed, with $V^{(0)}_0=1$ and $V^{(0)}|_{Z'}=0$, and $G^{(0)}$
in the full Hermite space, with $G^{(0)}|_{\{0\}\cup Z'}=0$ and $G^{(0)}_z=1$.
By the P\'olya--Szeg\H o zero bound $V^{(0)}$, a nonzero element of an
$(m-1)$-dimensional Hermite space with the $m-2$ zeros $Z'$, has no other real
zero, so $V^{(0)}_z\ne0$; likewise $G^{(0)}$ vanishes only on
$\{0\}\cup Z'$, and since $|Z'|=m-2<z-1$ some $1\le d<z$ avoids $Z'$, so
$\Sigma_{G^{(0)}}>0$.  By Fatou's lemma, as in Proposition~\ref{prop:confluent},
\[
  \Tail(\xi)\le\liminf_{\varepsilon\to0}\Tail(\xi^{(\varepsilon)})
  \le\prod_{j=1}^{\ell+1}\frac{y_{(j)}}{1-y_{(j)}}-|V^{(0)}_z|\,\Sigma_{G^{(0)}}
  <\prod_{j=1}^{\ell+1}\frac{y_{(j)}}{1-y_{(j)}} .
\]
\end{proof}

\section{The partial-sum criterion}\label{sec:criterion}

This section proves that \eqref{eq:order} is the infimum exactly when the
partial-sum condition \eqref{eq:partial} below holds.  At roots $(2,3)$ the
companion paper already exhibits monic multiples of $(x-2)(x-3)$ whose
$b_2$ decreases to $2$~\cite[Proposition~4.4]{coefficient-mass}, so the $u=1$
bound is an infimum there.  The general construction writes the multiple
from the top: one leading $1$, a short free prefix, a long run of a single
value $t$, and $u$ free coefficients at the bottom, which the bound on
$b_{u+1}$ discards.  After division by
$r_i^D$ these matter only at the $u$ smallest roots,
which they absorb; the run must therefore satisfy the remaining $L-u$ roots
alone.

Write $P_u(x)=\prod_{i>u}(x-r_i)$ for the
surviving roots, $m=L-u$, and let $\widetilde P_u(y)=y^mP_u(1/y)$ be its
reverse.  In the reversed variables the limit construction is
\begin{equation}\label{eq:wgen}
  W(y)=\sum_{d\ge0}w_dy^d=\frac{\widetilde P_u(y)}{1-y},
\end{equation}
which vanishes at $y=1/r_i$ for every surviving root because
$\widetilde P_u$ does.  Dividing by $1-y$, its coefficients are the partial sums
$S_j=c_m+c_{m-1}+\cdots+c_{m-j}$ of the coefficients $c_k$ of $P_u$ read down
from the leading $c_m=1$, and they stabilize at
$S_m=P_u(1)=(-1)^m\prod_{i>u}(r_i-1)$ once $j\ge m$.  So the free prefix is
$S_1,\ldots,S_{m-1}$ and the run value is $S_m=P_u(1)$, whose magnitude
$|S_m|$ is exactly the right side of \eqref{eq:order} with $f_D=1$: the construction reproduces the bound because the sum of the
coefficients of $P_u$ is $P_u(1)$.

At $u=0$ this criterion and this series are due to Dubickas and
Jankauskas~\cite[Corollary~7]{dj}: their Corollary~7 gives
$H(P)=|P(1)|$ for the reduced height $H(P)$, the $u=0$ infimum, under the
same partial-sum condition and their same-sign hypothesis (which~\cite[Lemma~2.1]{coefficient-mass}
removes), and exhibits $\widetilde P/(1-y)$ as the extremal
series.  What is added here is the passage to $u\ge1$,
where the $u$ free bottom coefficients absorb the $u$ smallest roots.  The run is
the $(u+1)$-st largest magnitude only if no partial sum exceeds the total,
\begin{equation}\label{eq:partial}
  |S_j|\le|S_m|\qquad(1\le j<m),
\end{equation}
and \eqref{eq:partial} can fail.

\begin{proposition}[Partial-sum criterion]\label{prop:partial}
Let $2\le r_1<\cdots<r_L$ be distinct, $0\le u\le L-1$ and $m=L-u$.  If
$|S_j(P_u)|\le|P_u(1)|$ for $1\le j<m$ (condition \eqref{eq:partial}), then
$\inf_Qb_{u+1}\bigl(Q\prod_{i=1}^L(x-r_i)\bigr)=\prod_{i=u+1}^L(r_i-1)$ over
monic $Q\in\R[x]$.
\end{proposition}

\begin{proof}
The bound~\cite[Theorem~3.2]{coefficient-mass} gives $\ge$.  For the other direction put $m=L-u$ and write, for $D\ge L$,
$F_D=x^D+\sum_{k=1}^{m-1}a_kx^{D-k}+t\sum_{j=u}^{D-m}x^j+C$ with $\deg C<u$,
the $L$ unknowns fixed by $F_D(r_i)=0$.  This linear system is
nonsingular: if $H$ is a kernel element, the polynomial $(x-1)H$ has at most
$L+1$ nonzero terms, hence at most $L$ positive roots by Descartes' rule,
yet it vanishes at the $L+1$ positive numbers $1,r_1,\ldots,r_L$, so $H=0$.
Divide by $r_i^D$, sum the run, and trade the coefficients of $C$ for their
values $\sigma_i=r_i^{-D}C(r_i)$ ($i\le u$), a Vandermonde change of
variables: with $\ell_1,\ldots,\ell_u$ the Lagrange basis at $r_1,\ldots,r_u$,
$C=\sum_{q\le u}\sigma_qr_q^D\ell_q$.  In the unknowns $a_k$, $t$, $\sigma_q$
and with $y_i=1/r_i$, the equation at $r_i$ reads
\[
  1+\sum_{k=1}^{m-1}a_ky_i^k+t\,\frac{y_i^m-y_i^{D-u+1}}{1-y_i}
  +\sum_{q\le u}\sigma_q\Bigl(\frac{r_q}{r_i}\Bigr)^D\ell_q(r_i)=0 ,
\]
where the last sum is $\sigma_i$ for $i\le u$.  As $D\to\infty$ the
coefficient matrix converges entrywise: $y_i^{D-u+1}\to0$, and for $i>u$ the
cross coefficients are $O((r_u/r_{u+1})^D)\to0$ (for $u\ge1$; at $u=0$ there
is no $C$).  In the limiting matrix
each row $i\le u$ has coefficient $1$ on $\sigma_i$ and $0$ on the other
$\sigma_q$, and each row $i>u$ has $0$ on every $\sigma_q$, so it is block
triangular.  Multiplied by $r_i^{m-1}(r_i-1)$, a limiting surviving row says
that the monic polynomial $(x-1)A+t$ of degree $m$, with
$A=x^{m-1}+a_1x^{m-2}+\cdots+a_{m-1}$, vanishes at $r_i$.  Every monic $R$ of
degree $m$ is $(x-1)A+t$ for exactly one such pair, namely the quotient and
remainder of $R$ by $x-1$, and exactly one monic $R$ of degree $m$ vanishes
at the $m$ distinct surviving roots, namely $P_u$.  So the surviving block
is nonsingular, being square with a unique solution, hence the
limiting matrix is, and since matrix inversion is
continuous on the invertible matrices the solutions converge: as
$D\to\infty$,
\[
  t\to P_u(1),\qquad
  a_k\to S_k\quad(1\le k<m),
\]
since $(P_u-P_u(1))/(x-1)$ has as coefficients the prefix
$S_1,\ldots,S_{m-1}$ of \eqref{eq:wgen}.  The nonleading coefficients of
$F_D$ are $a_1,\ldots,a_{m-1}$, copies of $t$, and the $u$ coefficients of
$C$; discarding the latter, $b_{u+1}(F_D)\le\max(|a_1|,\ldots,|a_{m-1}|,|t|)$,
so
\[
  \limsup_{D\to\infty}b_{u+1}(F_D)\le\max_{1\le j\le m}|S_j|,
\]
which is $|S_m|$ when \eqref{eq:partial} holds.  Distinctness enters three
times---the Descartes step, the Vandermonde change, and the surviving
block---and $r_u<r_{u+1}$ (for $u\ge1$) is what kills the cross terms.
\end{proof}

The conclusion does not need distinct roots (Corollary~\ref{cor:charrep}).

\subsection*{Necessity}

With $y_i=1/r_i$, an exponential sum $v_d=\sum_i\lambda_iy_i^d$ on some of
these nodes is a \emph{certificate for} $X\subset\N_{>0}$ if $v_0=1$ and
$v_w=0$ for $w\in X$; certificates show that the criterion is also
necessary, repeated roots included.  The argument runs
on the dual side of~\cite[Lemma~3.1]{coefficient-mass}; we set it up for distinct roots, and the proof
of Theorem~\ref{thm:converse} carries it to repeated roots.  Here
$1/2\ge y_1>\cdots>y_L>0$, the reverse of the order in~\cite[Theorem~2.2]{coefficient-mass}.  If $F$ is a monic
multiple of $\prod_{i=1}^L(x-r_i)$ of degree $D$ and $X$ is the set of
distances $D-j$ of $u$ positions $j$ carrying its $u$ largest nonleading
magnitudes, the identity in the proof of~\cite[Lemma~3.1]{coefficient-mass} reads
$0=1+\sum_{j<D}f_jv_{D-j}$, the terms with $D-j\in X$ vanish, and every other
$|f_j|$ is at most $b_{u+1}(F)$, so
\begin{equation}\label{eq:weak}
  1\le b_{u+1}(F)\,\Tail(v)
\end{equation}
for every certificate $v$ for $X$ (this uses only $r_i>1$).  Write
$B=\prod_{i=u+1}^L(r_i-1)^{-1}$.  So, to show that the infimum exceeds $1/B$ when \eqref{eq:partial}
fails, it suffices to find a $\delta>0$ such that every $X$ of size $u$ has a certificate with
$\Tail(v)\le B-\delta$.  Three ingredients supply the pieces: the partial
sums alternate in sign (Lemma~\ref{lem:signs}), a failing partial sum
improves the consecutive certificate (Lemma~\ref{lem:improve}), and zeros far
from the origin cost only the largest nodes
(Proposition~\ref{prop:farrep}).

\begin{lemma}[Signs of the partial sums]\label{lem:signs}
Let $r_1,\ldots,r_m>1$, let $S_0,\ldots,S_m$ be the partial sums of the
coefficients of $\prod_{i=1}^m(x-r_i)$ read down from the leading one, and put
$s_i=r_i-1$.  Then
\[
  (-1)^jS_j=\sum_{k=0}^{j}e_k(s)\binom{m-1-k}{j-k}>0\quad(0\le j\le m-1),
  \qquad (-1)^mS_m=e_m(s),
\]
where $e_k$ is the elementary symmetric polynomial.
\end{lemma}

\begin{proof}
By \eqref{eq:wgen}, $S_j$ is the coefficient of $y^j$ in
$\prod_i(1-r_iy)/(1-y)$, so $(-1)^jS_j$ is the coefficient of $y^j$ in
$\prod_i(1+r_iy)/(1+y)$.  Writing $1+r_iy=(1+y)+s_iy$ and grouping the
expansion by the number of factors $s_iy$ chosen,
\[
  \frac{\prod_i(1+r_iy)}{1+y}
  =\sum_{k=0}^{m-1}e_k(s)\,y^k(1+y)^{m-1-k}+e_m(s)\frac{y^m}{1+y}.
\]
The terms with $k<m$ are polynomials of degree $m-1$ and contribute
$e_k(s)\binom{m-1-k}{j-k}$ to $y^j$; the last contributes nothing below $y^m$
and $e_m(s)$ at $y^m$.  The $k=0$ term alone gives $\binom{m-1}{j}\ge1$.
\end{proof}

In particular \eqref{eq:partial} reads
$\sum_{k\le j}e_k(s)\binom{m-1-k}{j-k}\le e_m(s)$ for $1\le j<m$, with
$s_i=r_i-1$ over the surviving roots.

\begin{lemma}[Improving the consecutive certificate]\label{lem:improve}
Let $\rho_1,\ldots,\rho_n>1$ be distinct and $R=\prod_{i=1}^n(x-\rho_i)$ with
partial sums $S_j(R)$.  Call an exponential sum $v_d=\sum_i\lambda_iy_i^d$ on
the nodes $y_i=1/\rho_i$ a certificate for $X\subset\N_{>0}$ if $v_0=1$ and
$v|_X=0$, and let $v^*$ be the certificate with
zero set $\{1,\ldots,n-1\}$, so that $\Tail(v^*)=1/|R(1)|$ by~\cite[Lemma~2.1]{coefficient-mass}.  If
$X_0\subseteq\{1,\ldots,n-1\}$ and some $j\in\{1,\ldots,n-1\}\setminus X_0$
has $|S_j(R)|>|R(1)|$, then some certificate for $X_0$ on the same nodes has
tail less than $1/|R(1)|$.
\end{lemma}

\begin{proof}
Put $y_{\max}=\max_iy_i$.  Let $\eta$ be the exponential sum on the nodes with
$\eta_d=[d=j]$ for $0\le d\le n-1$; it exists by the Vandermonde argument.
With $\Pi(x)=\prod_i(1-x/\rho_i)=\sum_s\pi_sx^s$, its generating function is
$N(x)/\Pi(x)$ with $\deg N\le n-1$, and $N$ is the truncation of
$\Pi\sum_d\eta_dx^d$ below degree $n$, namely
$N(x)=\sum_{k=j}^{n-1}\pi_{k-j}x^k$.  Since $\Pi$ is $R$ times the constant
$\prod_i(-1/\rho_i)$, evaluating at $x=1$ gives
\[
  \sum_{d\ge n}\eta_d=\frac{N(1)}{\Pi(1)}-1
  =-\frac{\pi_{n-j}+\cdots+\pi_n}{\Pi(1)}
  =-\frac{S_j(R)}{R(1)} .
\]
By~\cite[Lemma~2.1]{coefficient-mass}, $v^*_d$ has the constant sign $(-1)^{n+1}$ for $d\ge n$
and $|v^*_d|\ge\bigl(\prod_iy_i\bigr)y_{\max}^{\,d-n}$, while
$|\eta_d|\le C_\eta y_{\max}^{\,d}$ for a constant $C_\eta$; so for
$\kappa>0$ small enough $v=v^*+t\eta$ with
$t=(-1)^{n+1}\kappa\sgn\bigl(S_j(R)/R(1)\bigr)$ keeps the sign
$(-1)^{n+1}$ at every $d\ge n$.  It is a certificate for $X_0$, since it vanishes
on $\{1,\ldots,n-1\}\setminus\{j\}$, and
\[
  \Tail(v)=\kappa+(-1)^{n+1}\sum_{d\ge n}\bigl(v^*_d+t\eta_d\bigr)
  =\Tail(v^*)-\kappa\Bigl(\Bigl|\frac{S_j(R)}{R(1)}\Bigr|-1\Bigr)
  <\Tail(v^*).
\]
\end{proof}

\paragraph{Lemma~\ref{lem:improve} with repetitions.}  Let
$\rho_1,\ldots,\rho_n>1$, repetitions allowed, $R=\prod_{i=1}^n(x-\rho_i)$
and $y_i=1/\rho_i$, and let $z_1>\cdots>z_p$ be the distinct values among the
$y_i$, $z_l$ occurring $m_l$ times, so that $y_{\max}=z_1$ has multiplicity
$\mu=m_1$.  The \emph{Hermite space} $\mathcal H(R)$ is the span of the $n$
sequences $d\mapsto d^az_l^d$ ($d\ge0$, $1\le l\le p$, $0\le a<m_l$); for
$R=\prod_{i>k}(x-r_i)$ it is the space $\mathcal H_k$ of
Proposition~\ref{prop:farrep}, and for distinct $\rho_i$ it is the span of
the $y_i^d$.  A certificate for $X$ on $\mathcal H(R)$ is a $v\in\mathcal H(R)$
with $v_0=1$ and $v|_X=0$.  Then Lemma~\ref{lem:improve} holds with the
nodes replaced by $\mathcal H(R)$: there is exactly one certificate $v^*$ on
$\mathcal H(R)$ with zero set $\{1,\ldots,n-1\}$, it has
$\Tail(v^*)=1/|R(1)|$, and if $X_0\subseteq\{1,\ldots,n-1\}$ and some
$j\in\{1,\ldots,n-1\}\setminus X_0$ has $|S_j(R)|>|R(1)|$, then some
certificate for $X_0$ on $\mathcal H(R)$ has tail less than $1/|R(1)|$.

\begin{proof}
Three facts about $\mathcal H(R)$ replace the Vandermonde argument and the
explicit exponential sums of the distinct case.  Put
$\Pi(x)=\prod_{i=1}^n(1-y_ix)=\sum_{s=0}^n\pi_sx^s$.

(i) \emph{Unisolvence.}  A sum $\sum_lq_l(d)z_l^d$ with $\deg q_l<m_l$ and
not all $q_l$ zero has at most $n-1$ real
zeros~\cite[Part~Five, Ch.~1]{polya-szego}.  So the $n$ spanning sequences
are linearly independent, $\dim\mathcal H(R)=n$, and for distinct integers
$d_1,\ldots,d_n\ge0$ the linear map $v\mapsto(v_{d_1},\ldots,v_{d_n})$ from
$\mathcal H(R)$ to $\R^n$ is injective, hence bijective: any $n$ values at
$d_1,\ldots,d_n$ are taken by exactly one element of $\mathcal H(R)$.

(ii) \emph{Generating functions.}  $\mathcal H(R)$ is the set of sequences
whose generating function is $N(x)/\Pi(x)$ with $\deg N\le n-1$.  By
induction on $a$,
$\sum_{d\ge0}d^az^dx^d=(x\,d/dx)^a(1-zx)^{-1}=A_a(x)/(1-zx)^{a+1}$ with
$\deg A_a\le a$, since $x\,d/dx$ sends $A/(1-zx)^{a+1}$ to
$\bigl(xA'(1-zx)+(a+1)zxA\bigr)/(1-zx)^{a+2}$.  For $z=z_l$ and $a<m_l$ the
factor $(1-z_lx)^{a+1}$ divides $\Pi$, so this is $N/\Pi$ with
$\deg N\le a+n-(a+1)=n-1$.  So $\mathcal H(R)$ lies in the space of such
sequences, which has dimension $n$, and by (i) the two coincide.  For
$v\in\mathcal H(R)$ with generating function $V=N/\Pi$, $N=\Pi V$ as formal
power series, and its coefficients below $x^n$ involve only
$v_0,\ldots,v_{n-1}$:
\[
  N(x)=\sum_{k=0}^{n-1}\Bigl(\sum_{s=0}^k\pi_sv_{k-s}\Bigr)x^k .
\]

(iii) \emph{Growth.}  For $v\in\mathcal H(R)$ some $C_v$ has
$|v_d|\le C_v(1+d)^{\mu-1}y_{\max}^{\,d}$ for all $d\ge0$: the spanning
sequences at $z_1$ have $a\le\mu-1$, and for $l\ge2$ each
$d^a(z_l/z_1)^d$ is bounded.  As $y_{\max}<1$, $\sum_d|v_d|<\infty$, so the
identity $\Pi V=N$ may be evaluated at $x=1$ (a Cauchy product of an
absolutely convergent series with a polynomial):
$\Pi(1)\sum_{d\ge0}v_d=N(1)$, where $\Pi(1)=\prod_i(1-y_i)>0$.

\emph{The certificate $v^*$.}  By (i) it exists and is unique.  Its values
at $0,\ldots,n-1$ are $1,0,\ldots,0$, so by (ii)
$N=\sum_{k<n}\pi_kx^k=\Pi-\pi_nx^n$ and its generating function is
$1-\pi_nx^n/\Pi(x)$.  Expanding
$1/\Pi=\prod_i\sum_{e\ge0}y_i^ex^e=\sum_{e\ge0}h_e(y)x^e$, with $h_e$ the
complete homogeneous symmetric polynomial of the multiset
$y_1,\ldots,y_n$, and using $\pi_n=(-1)^n\prod_iy_i$,
\[
  v^*_d=(-1)^{n+1}\Bigl(\prod_iy_i\Bigr)h_{d-n}(y)\qquad(d\ge n).
\]
All $y_i$ are positive, so $h_e(y)>0$, $v^*_d$ has the sign $(-1)^{n+1}$
for $d\ge n$, and
\[
  \Tail(v^*)=\prod_iy_i\sum_{e\ge0}h_e(y)
  =\prod_i\frac{y_i}{1-y_i}=\prod_i\frac1{\rho_i-1}=\frac1{|R(1)|}.
\]
Keeping in $h_e(y)$ only the monomials in the $\mu$ copies of $y_{\max}$, of
which there are $\binom{e+\mu-1}{\mu-1}\ge(e+1)^{\mu-1}/(\mu-1)!$, gives
\[
  |v^*_d|\ge\frac{\prod_iy_i}{(\mu-1)!}\,(d-n+1)^{\mu-1}y_{\max}^{\,d-n}
  \qquad(d\ge n).
\]

\emph{The sum $\eta$.}  By (i) there is exactly one $\eta\in\mathcal H(R)$
with $\eta_d=[d=j]$ for $0\le d\le n-1$, and by (ii) its generating function
is $N/\Pi$ with $N(x)=\sum_{k=j}^{n-1}\pi_{k-j}x^k$.  Since
$\Pi(x)=\prod_i(-1/\rho_i)\cdot R(x)$, each $\pi_s$ is that constant times
the coefficient of $x^s$ in $R$, so (iii) gives
\[
  \sum_{d\ge n}\eta_d=\frac{N(1)}{\Pi(1)}-1
  =-\frac{\pi_{n-j}+\cdots+\pi_n}{\Pi(1)}
  =-\frac{S_j(R)}{R(1)} ,
\]
exactly as for distinct nodes.

\emph{Domination.}  By (iii), $|\eta_d|\le C_\eta(1+d)^{\mu-1}y_{\max}^{\,d}$;
since $(n+1)(d-n+1)-(1+d)=n(d-n)\ge0$ for $d\ge n$, the lower bound on
$|v^*_d|$ gives
\[
  |\eta_d|\le C|v^*_d|\quad(d\ge n),\qquad
  C=\frac{C_\eta\,(\mu-1)!\,(n+1)^{\mu-1}y_{\max}^{\,n}}{\prod_iy_i}.
\]

\emph{The improved certificate.}  Here $R(1)\ne0$ and $|S_j(R)|>|R(1)|$.
Fix $0<\kappa<1/C$ and put $t=(-1)^{n+1}\kappa\sgn\bigl(S_j(R)/R(1)\bigr)$
and $v=v^*+t\eta\in\mathcal H(R)$.  Then $v_0=1$ because $j\ge1$, $v$
vanishes on $\{1,\ldots,n-1\}\setminus\{j\}\supseteq X_0$, and $v_j=t$, so
$v$ is a certificate for $X_0$ on $\mathcal H(R)$; and for $d\ge n$,
$|t\eta_d|<|v^*_d|$, so $v_d$ has the sign $(-1)^{n+1}$.  All the series
converge absolutely, and
\[
  \Tail(v)=\kappa+(-1)^{n+1}\sum_{d\ge n}\bigl(v^*_d+t\eta_d\bigr)
  =\Tail(v^*)-\kappa\Bigl(\Bigl|\frac{S_j(R)}{R(1)}\Bigr|-1\Bigr)
  <\Tail(v^*).
\]
\end{proof}

Zeros far from the origin cost only the largest nodes because the leading
functions of the basis cannot produce them, and that is a Descartes rule at
infinity.  Let $\mathcal C$ be the
set of functions $x\mapsto R(x)w^x$ with $R\in\R(x)$ nonzero and $w>0$, each
considered on a half-line beyond the real poles of $R$.  It is a group under
multiplication; each member has finitely many zeros, so it is eventually of
one sign; the quotient of two members,
$(R_1/R_2)(w_1/w_2)^x$, has a limit in $[-\infty,\infty]$ as $x\to\infty$;
and the derivative of a nonconstant member lies in $\mathcal C$, because
$(Rw^x)'=(R'+R\log w)w^x$ and, writing $R=p/q$, the numerator
$p'q-pq'+pq\log w$ of $R'+R\log w$ has a term of degree $\deg p+\deg q$ that
the other two cannot cancel when $w\ne1$, while $R'\ne0$ for nonconstant $R$
when $w=1$.

\begin{lemma}[Descartes at infinity]\label{lem:infinity}
Let $\varphi_1,\ldots,\varphi_M\in\mathcal C$ satisfy
$\varphi_{j+1}(x)/\varphi_j(x)\to0$ as $x\to\infty$.  Let
$c^{(\nu)}\in\R^M$ converge to $c\ne0$, let $j_0$ be the least index with
$c_{j_0}\ne0$, and let $n_\nu\to\infty$.  Then for all large $\nu$ the
function $f_\nu=\sum_jc^{(\nu)}_j\varphi_j$ has at most $j_0-1$ zeros in
$[n_\nu,\infty)$.
\end{lemma}

\begin{proof}
Induct on $j_0$, for all such families at once.  Every ratio
$\varphi_j/\varphi_1$, $j\ge2$, is a product of consecutive ratios and tends
to $0$.  If $j_0=1$ then on $[n_\nu,\infty)$
\[
  \Bigl|\frac{f_\nu}{\varphi_1}-c^{(\nu)}_1\Bigr|
  \le\max_j|c^{(\nu)}_j|\sum_{j\ge2}\sup_{x\ge n_\nu}
  \Bigl|\frac{\varphi_j(x)}{\varphi_1(x)}\Bigr|\longrightarrow0,
\]
while $c^{(\nu)}_1\to c_1\ne0$, so $f_\nu$ has no zero there for large $\nu$.

If $j_0\ge2$, then $c_1=0$.  Put $\chi_j=(\varphi_j/\varphi_1)'$ for
$2\le j\le M$.  Each $\varphi_j/\varphi_1$ tends to $0$ and is not constant,
so $\chi_j\in\mathcal C$.  For $2\le j<M$ the functions
$\varphi_{j+1}/\varphi_1$ and $\varphi_j/\varphi_1$ tend to $0$, their
quotient $\varphi_{j+1}/\varphi_j$ tends to $0$, and the quotient
$\chi_{j+1}/\chi_j$ of their derivatives lies in $\mathcal C$ and so has a
limit in $[-\infty,\infty]$, while $\chi_j$ is eventually nonzero.
By L'H\^opital's rule the limit of $\chi_{j+1}/\chi_j$ is that of
$\varphi_{j+1}/\varphi_j$, which is $0$, so the family $\chi_2,\ldots,\chi_M$ satisfies the hypothesis.  Its coefficient
vectors $(c^{(\nu)}_2,\ldots,c^{(\nu)}_M)$ converge to $(c_2,\ldots,c_M)$,
whose first nonzero entry has index $j_0-1$ in the new family, so by induction
$(f_\nu/\varphi_1)'=\sum_{j\ge2}c^{(\nu)}_j\chi_j$ has at most $j_0-2$ zeros
in $[n_\nu,\infty)$ for large $\nu$.  If $f_\nu$ had $j_0$ zeros there, so
would $f_\nu/\varphi_1$, and Rolle's theorem would give $j_0-1$ zeros of its
derivative between them.
\end{proof}

\begin{proposition}[Far zeros at repeated roots]\label{prop:farrep}
Let $r_1\le\cdots\le r_L$ with $r_1>1$, repetitions allowed, let
$z_1>\cdots>z_p$ be the distinct values of $1/r_i$ with multiplicities $m_i$,
and let $\psi_1,\ldots,\psi_L$ be the Hermite basis $d^az_i^d$,
$0\le a<m_i$, ordered by growth: $z$ descending, and $a$ descending within a
node.  For $0\le k\le L$ let
$\mathcal H_k=\operatorname{span}\{\psi_{k+1},\ldots,\psi_L\}$, and call
$v\in\mathcal H_k$ a certificate for $X$ on $\mathcal H_k$ if $v_0=1$ and
$v_x=0$ for $x\in X$.  Fix $0\le k<L$, and let $E_\nu$ be $k$-sets of positive
integers with $\min E_\nu\to\infty$.
\begin{enumerate}
\item[(a)] If $f_\nu\in\mathcal H_0$ vanishes on $E_\nu$ and its coefficient
vector has sup-norm $1$, its coefficients on $\psi_1,\ldots,\psi_k$ tend to
$0$.
\item[(b)] With $J^*=\{1,\ldots,k\}$ and
$\alpha_J(E)=\det(\psi_j(e))_{e\in E,\,j\in J}$, for large $\nu$ one has
$\alpha_{J^*}(E_\nu)\ne0$, and
$\eta_\nu=\max_{J\ne J^*}|\alpha_J(E_\nu)/\alpha_{J^*}(E_\nu)|\to0$.
\item[(c)] If $v^\infty$ is a certificate for a finite set $X_0$ on
$\mathcal H_k$, then for large $\nu$ there are certificates $v^{(\nu)}$ for
$X_0\cup E_\nu$ on $\mathcal H_0$ with $\Tail(v^{(\nu)}-v^\infty)\to0$.
\end{enumerate}
\end{proposition}

Removing the first $k$ functions removes whole blocks at the largest nodes
and, at the node where the cut falls, the top exponents; so $\mathcal H_k$ is
the Hermite space of the multiset $r_{k+1},\ldots,r_L$, and by the zero bound
for exponential polynomials~\cite[Part~Five, Ch.~1]{polya-szego} a nonzero
element of $\mathcal H_k$ has at most $L-k-1$ real zeros.  For distinct roots
$\mathcal H_k$ is the span of
$y_{k+1}^d,\ldots,y_L^d$.

\begin{proof}
(a) Otherwise, along a subsequence, the coefficient vectors converge to some
$c$ with a nonzero coordinate among the first $k$, so the least index $j_0$
with $c_{j_0}\ne0$ is at most $k$.  Each $\psi_j(x)=x^az_i^x$ lies in
$\mathcal C$, and consecutive ratios are $1/x$ within a node and
$x^{m_{i+1}-1}(z_{i+1}/z_i)^x$ across nodes, both tending to $0$.  By
Lemma~\ref{lem:infinity} with $n_\nu=\min E_\nu$, $f_\nu$ has at most
$j_0-1\le k-1$ zeros in $[\min E_\nu,\infty)$ for large $\nu$, but it vanishes
on the $k$ points of $E_\nu$.

(b) If $\alpha_{J^*}(E_\nu)=0$ along a subsequence, a combination of
$\psi_1,\ldots,\psi_k$ of sup-norm $1$ vanishes on $E_\nu$, against (a).
Otherwise let $\Psi(E)=(\psi_j(e))_{e\in E,\,j\le L}$ and $\Psi_{J^*}(E)$ its
first $k$ columns, so $\alpha_{J^*}=\det\Psi_{J^*}$, and put
$\Lambda_\nu=\Psi_{J^*}(E_\nu)^{-1}\Psi(E_\nu)$, a $k\times L$ matrix with rows
indexed by $j\le k$ whose first $k$ columns form the identity; its minor on
the columns $J$ is $\alpha_J/\alpha_{J^*}$.  For
$l>k$ the function $\psi_l-\sum_{j\le k}(\Lambda_\nu)_{jl}\psi_j$ vanishes on
$E_\nu$; its coefficient vector has sup-norm $\max(1,t)$ with
$t=\max_j|(\Lambda_\nu)_{jl}|$, so by (a) $t/\max(1,t)\to0$, that is $t\to0$.
Every $J\ne J^*$ contains some $l>k$, so its minor is a determinant with a
column tending to $0$ and the others bounded, and $\eta_\nu\to0$.

(c) Put $n=L-k$; the zero bound in $\mathcal H_k$ gives $|X_0|\le n-1$, so
we may fix a set $K$ of $n-1-|X_0|$ positive integers outside $X_0$.  For
large $\nu$ the $L$ points of $Y_\nu=\{0\}\cup X_0\cup K\cup E_\nu$ are
distinct, so there is a unique $v^{(\nu)}\in\mathcal H_0$ with
$v^{(\nu)}_0=1$, $v^{(\nu)}|_{X_0}=0$, $v^{(\nu)}_q=v^\infty_q$ for $q\in K$
and $v^{(\nu)}|_{E_\nu}=0$; it is a certificate for $X_0\cup E_\nu$, and
$v^\infty$ is the unique element of $\mathcal H_k$ meeting the same
conditions at the $n$ points $\{0\}\cup X_0\cup K$.  Let $c^{(\nu)}$ be the
coefficient vector of $v^{(\nu)}$.  If $\|c^{(\nu)}\|_\infty\to\infty$ along
a subsequence, then $v^{(\nu)}/\|c^{(\nu)}\|_\infty$ has, by (a) and
compactness, a limit $g\in\mathcal H_k$ of sup-norm $1$ vanishing at the $n$
points $\{0\}\cup X_0\cup K$, against the zero bound.  So the $c^{(\nu)}$ are
bounded, and also $\|c^{(\nu)}\|_\infty\ge1$ because $v^{(\nu)}_0=1$; by (a)
their first $k$ coordinates tend to $0$, so every limit point is the
coefficient vector of an element of $\mathcal H_k$ meeting the $n$ conditions
of $v^\infty$, that is, of $v^\infty$.  Hence $c^{(\nu)}\to c^\infty$, and
\[
  \Tail(v^{(\nu)}-v^\infty)\le\sum_{j=1}^L|c^{(\nu)}_j-c^\infty_j|
  \sum_{d\ge1}|\psi_j(d)|\longrightarrow0 .
\]
\end{proof}

\begin{theorem}[Necessity of the partial-sum criterion]\label{thm:converse}
Let $2\le r_1\le\cdots\le r_L$, repetitions allowed, $0\le u\le L-1$ and
$m=L-u$.  If $|S_j(P_u)|>|P_u(1)|$ for some $1\le j<m$, that is, if
\eqref{eq:partial} fails, then
\[
  \inf_Qb_{u+1}\Bigl(Q\prod_{i=1}^L(x-r_i)\Bigr)>\prod_{i=u+1}^L(r_i-1)
\]
over monic $Q\in\R[x]$.  With Corollary~\ref{cor:charrep},
\eqref{eq:order} is the infimum exactly when \eqref{eq:partial} holds.
\end{theorem}

\begin{proof}
Use the notation of Proposition~\ref{prop:farrep}.  If $F$ is a monic
multiple of $P=\prod_i(x-r_i)$ of degree $D$ and $v\in\mathcal H_0$, then
$\sum_{j\le D}f_jv_{D-j}=0$: for $\psi=d^az^d$ with $z=1/r$, expanding
$(D-j)^a$ in powers of $j$ writes $\sum_jf_j(D-j)^ar^{j-D}$ as a combination
of the values $\bigl((x\,d/dx)^bF\bigr)(r)$, $b\le a$, which vanish because
$r$ is a root of multiplicity greater than $a$.  So \eqref{eq:weak} holds
for every certificate on $\mathcal H_0$, and, as in the distinct case, it
suffices to find $\delta>0$ such that every $X\subset\N_{>0}$ of size $u$ has
a certificate on $\mathcal H_0$ with $\Tail\le B-\delta$.  Suppose not, and
take $X_\nu$ whose certificates on $\mathcal H_0$ all have
$\Tail\ge B-1/\nu$.  Passing to a subsequence, $X_\nu=X_0\cup E_\nu$ with
$X_0$ fixed, $|E_\nu|=k$ and, if $k\ge1$, $\min E_\nu\to\infty$.  If
$k\ge1$, by Proposition~\ref{prop:farrep}(c) a certificate $v^\infty$ for $X_0$ on
$\mathcal H_k$ with $\Tail(v^\infty)<B$ yields certificates for $X_\nu$
with tails tending to $\Tail(v^\infty)$, a contradiction; if $k=0$, then
$E_\nu=\varnothing$, $X_\nu=X_0$, and $v^\infty$ itself is a certificate for
$X_\nu$ on $\mathcal H_0$ with tail below $B-1/\nu$ for large $\nu$, again a
contradiction.  So it suffices to produce one.  Put
$n=L-k$ and $m=L-u$, so $|X_0|=u-k$; recall that $\mathcal H_k$ is the
Hermite space of $r_{k+1},\ldots,r_L$.

If $X_0=\varnothing$, then $k=u$, so $\mathcal H_k=\mathcal H_u$ is the Hermite space of the roots of $P_u$, and
since \eqref{eq:partial} fails, Lemma~\ref{lem:improve} (with repetitions)
for $R=P_u$ and $X_0=\varnothing$ gives $\Tail(v^\infty)<1/|P_u(1)|=B$.

If $X_0\ne\varnothing$, take the certificate of the proof of~\cite[Lemma~3.1]{coefficient-mass}
on $\mathcal H_k$:
zero set $Z$ consisting of $X_0$ and the first $m-1$ positive integers
outside it, so $|Z|=n-1$ and its leading run is $\ell\ge m-1$.
Proposition~\ref{prop:confluent}, with nodes at most $1/2$, bounds $\Tail$ by
the product of $1/(r_i-1)$ over the $\ell+1$ largest of $r_{k+1},\ldots,r_L$,
strictly unless $Z=\{1,\ldots,n-1\}$ by Lemma~\ref{lem:confstrict}.  Every
factor is at most $1$ and $\ell+1\ge m$, so $\Tail\le B$, strictly unless
$Z=\{1,\ldots,n-1\}$ and $r_{k+1}=\cdots=r_u=2$.  In that case
$X_0\subseteq\{1,\ldots,n-1\}$ and $P_k=\prod_{i>k}(x-r_i)=(x-2)^{u-k}P_u$
has $|P_k(1)|=|P_u(1)|=1/B$, so Lemma~\ref{lem:improve} (with repetitions) for
$R=P_k$ and $X_0$ finishes the proof once some $j\in\{1,\ldots,n-1\}\setminus X_0$
has $|S_j(P_k)|>|P_k(1)|$.

For a monic $R$ with all roots above $1$ let
$\mathcal A(R)=\{1\le j\le\deg R-1:|S_j(R)|>|R(1)|\}$.  If $R$ has
coefficient $\zeta_l$ at $x^{\deg R-l}$ ($\zeta_0=1$, $\zeta_{-1}=0$), then
$(x-2)R$ has $\zeta_l-2\zeta_{l-1}$ at $x^{\deg R+1-l}$, so
$S_j((x-2)R)=S_j(R)-2S_{j-1}(R)$, and by Lemma~\ref{lem:signs}
$(-1)^jS_j(R)>0$ for $0\le j\le\deg R$; hence
$|S_j((x-2)R)|=|S_j(R)|+2|S_{j-1}(R)|$ for $1\le j\le\deg R$, while
$|((x-2)R)(1)|=|R(1)|$.  Thus $j\in\mathcal A(R)$ puts both $j$ and $j+1$ in
$\mathcal A((x-2)R)$.  Since \eqref{eq:partial} fails, $\mathcal A(P_u)$
contains some $i\le m-1$, and $u-k$ applications give
$\{i,i+1,\ldots,i+u-k\}\subseteq\mathcal A(P_k)\subseteq\{1,\ldots,n-1\}$.
These are $u-k+1$ integers and $X_0$ has only $u-k$ elements.
\end{proof}

\section{The placement game and attainment}\label{sec:value}

Keep the notation of the necessity argument in
Section~\ref{sec:criterion}, except that only $r_1>1$ is assumed: $1<r_1<\cdots<r_L$
are distinct, $y_i=1/r_i\in(0,1)$, and $m=L-u$.  For finite $X\subset\N_{>0}$ and
$0\le k<L$ put
\[
  \tau_k(X)=\inf\bigl\{\Tail(v):v\text{ a certificate for }X\text{ on }
  y_{k+1},\ldots,y_L\bigr\},
  \quad \tau(X)=\tau_0(X),\quad \tau^*=\tau_u(\varnothing).
\]
A certificate on $c$ nodes is \emph{full} if it vanishes at $c-1$ positive
integers, its \emph{zero set}.  By the zero bound it has no other real zeros,
so it is the unique certificate with that zero set, and it keeps one sign
past its largest zero; it is a certificate for $X$ exactly when its zero set
contains $X$.  By \eqref{eq:weak} every monic multiple $F$ has
$b_{u+1}(F)\ge1/\tau(X_F)$, where $X_F$ holds the distances from the top of
its $u$ largest nonleading coefficients.  So the infimum is set by the \emph{worst}
placement, the one maximising $\tau(X)$: worst for the certificate, and so
best for the multiple.  We call $\tau(X)$ the \emph{cost} of $X$.  The
first result says that the infimum is exactly $1/T$, where $T$ is the largest
cost of a placement, and that exempt coefficients escaping to the bottom
cost, in the limit, $\tau^*$, the best tail on the surviving roots alone.

\begin{theorem}[The infimum as a placement game]\label{thm:value}
Let $1<r_1<\cdots<r_L$ and $0\le u\le L-1$.  Then
\[
  \inf_Qb_{u+1}\Bigl(Q\prod_{i=1}^L(x-r_i)\Bigr)=\frac1T,
  \qquad T=\sup_{|X|=u}\tau(X),\qquad \tau^*\le T<\infty,
\]
over monic $Q\in\R[x]$, where $\tau(X)$ is the least tail of a certificate
for $X$ on $y_1,\ldots,y_L$ and $\tau^*$ the least tail of a sum $v$ on the
surviving nodes $y_{u+1},\ldots,y_L$ with $v_0=1$.
\end{theorem}

\begin{proof}
First, $T<\infty$.  Take $u$-sets $X_\nu$ with $\tau(X_\nu)\to T$.
Passing to a subsequence, each coordinate of $X_\nu$, listed in increasing
order, is either constant or tends to $\infty$, and the constant ones come
first.  Let $X_0$ be the set of the constant ones and $E'_\nu$ the rest, the
escaping part, and put $p=|E'_\nu|$, so $X_\nu=X_0\cup E'_\nu$, $|X_0|=u-p$,
and $\min E'_\nu\to\infty$ if $p\ge1$.  Since $u-p\le L-p-1$, enlarge $X_0$
to a set $Z$ of $L-p-1$ positive integers.  The interpolation system on
$y_{p+1},\ldots,y_L$ at the $L-p$ distinct points $\{0\}\cup Z$ is
nonsingular by the zero bound, so there is a full certificate $v^\infty$ on
these nodes with zero set $Z\supseteq X_0$.  If $p=0$, then
$T=\tau(X_0)\le\Tail(v^\infty)$.  Otherwise
Proposition~\ref{prop:farrep}(c), with $k=p$ and $E'_\nu$ in place of
$E_\nu$, gives certificates for $X_\nu$ with tails tending to
$\Tail(v^\infty)$, so $T\le\Tail(v^\infty)$.  Either way $T<\infty$.

The lower bound $b_{u+1}(F)\ge1/\tau(X_F)\ge1/T$ is \eqref{eq:weak}.

For the upper bound fix $X$ with $|X|=u$ and $D\ge\max(L,\max X)$, and let
$\tau_D(X)$ be the least $\sum_{d=1}^D|v_d|$ over certificates $v$ for $X$.
It is positive, since a certificate vanishing at $1,\ldots,D$ would vanish
identically.  Minimizing the largest $|f_j|$ over the positions $j<D$ with
$D-j\notin X$, subject to $F=x^D+\sum_{j<D}f_jx^j$ vanishing at every $r_i$,
is a feasible linear program.  Write $\mu_i$ for the multiplier of
$F(r_i)=0$, and put $\lambda_i=-\mu_ir_i^D$ and $v_d=\sum_i\lambda_iy_i^d$.
The coefficient of $f_j$ in the Lagrangian is then $-v_{D-j}$, and the dual
program is to maximize $v_0$ subject to $v_d=0$ for $d\in X$ and
$\sum_{1\le d\le D,\,d\notin X}|v_d|\le1$.  Its value is $1/\tau_D(X)$.  By
linear-programming duality some monic multiple of degree $D$ has every
coefficient off the $u$ positions of $X$ at most $1/\tau_D(X)$ in magnitude,
so its $b_{u+1}$ is at most $1/\tau_D(X)$.

Two limits finish the proof.  First, $\tau_D(X)$ is nondecreasing in $D$ and
at most $\tau(X)$.  Take minimizers $v^{(D)}=\sum_i\lambda^{(D)}_iy_i^d$.
Their values at $L$ fixed points of $\{0,1,\ldots\}\setminus X$, one of them
$0$, are bounded by $\max(1,\tau(X))$.  The matrix $(y_i^d)$ at $L$ distinct
points is nonsingular by the zero bound, so the $\lambda^{(D)}$ are bounded.
A subsequence converges to a certificate $v$ for $X$ with
$\sum_{d\le N}|v_d|\le\lim_D\tau_D(X)$ for every $N$, whence
$\tau(X)\le\Tail(v)\le\lim_D\tau_D(X)$.  So the infimum is at most $1/\tau(X)$
for every $X$, that is, at most $1/T$.

Second, put $E_\nu=\{\nu+1,\ldots,\nu+u\}$ and suppose
$\tau(E_\nu)\le\tau^*-\varepsilon$ along a subsequence.  Take certificates
$v^\nu=\sum_i\lambda^\nu_iy_i^d$ for $E_\nu$ with
$\Tail(v^\nu)\le\tau^*-\varepsilon/2$.  For $\nu\ge L$ their values at
$0,\ldots,L-1$ are bounded, so we may assume $\lambda^\nu\to\lambda$.  The $u$
equations $v^\nu_e=0$, $e\in E_\nu$, read
$A_\nu\lambda^\nu_{\le u}=-R_\nu\lambda^\nu_{>u}$ with
$A_\nu=(y_j^e)_{e\in E_\nu,\,j\le u}$ and
$R_\nu=(y_j^e)_{e\in E_\nu,\,j>u}$.  By Cramer's rule each entry of
$A_\nu^{-1}R_\nu$ is, up to sign, a ratio $\alpha_J/\alpha_{J^*}$ of
Proposition~\ref{prop:farrep}(b), whose basis is $y_1^d,\ldots,y_L^d$ for
distinct roots, with $k=u$ and $J\ne J^*$.  So it is at most $\eta_\nu\to0$, and $\lambda_{\le u}=0$.  Then
$v=\sum_{i>u}\lambda_iy_i^d$ is a certificate on the surviving nodes, and as
before $\Tail(v)\le\tau^*-\varepsilon/2$, against the definition of $\tau^*$.
Hence $\liminf_\nu\tau(E_\nu)\ge\tau^*$, and $T\ge\tau^*$.
\end{proof}

The proof gives $\liminf_\nu\tau(E_\nu)\ge\tau^*$ for the escaping
placements $E_\nu=\{\nu+1,\ldots,\nu+u\}$, and the reverse inequality
$\limsup_\nu\tau(E_\nu)\le\tau^*$ also holds: Proposition~\ref{prop:farrep}(c)
with $k=u$ and $X_0=\varnothing$ turns any sum $v^\infty$ on the surviving
nodes with $v^\infty_0=1$ into certificates for $E_\nu$ with tails tending
to $\Tail(v^\infty)$.  So $\tau(E_\nu)\to\tau^*$:
the escaping placement costs exactly $\tau^*$.  The escaping placement is
this sequence of placements, not a placement, and its cost is the limit
$\tau^*$; we say that \emph{the escaping placement is the worst one} when
$T=\tau^*$, that is, when no placement costs more than $\tau^*$.

On the dual side, pairing a sum on the surviving nodes with a sequence
whose series vanishes there bounds $\tau^*$ below, and each monic $M$ gives
such a sequence, so $\tau^*\ge1/\max_{j\ge1}|S_j(P_uM)|$.  Precisely, if
$v=\sum_{i>u}\lambda_iy_i^d$ is a certificate on the surviving nodes and $w$
is a bounded sequence with $w_0=1$ whose series $W(y)=\sum_dw_dy^d$ vanishes
at $y_{u+1},\ldots,y_L$, then $0=\sum_i\lambda_iW(y_i)=\sum_dw_dv_d$, so
\begin{equation}\label{eq:dualpair}
  1\le\sup_{d\ge1}|w_d|\cdot\Tail(v).
\end{equation}
Such $W$ are the $\widetilde P_uH$ with $H(0)=1$.  Taking
$H=\widetilde M/(1-y)$ for a monic $M\in\R[x]$ with reverse $\widetilde M$
makes $w_d$ the partial sums $S_d(P_uM)$, as in \eqref{eq:wgen}.  So
\begin{equation}\label{eq:multiplier}
  \inf_Qb_{u+1}\Bigl(Q\prod_{i=1}^L(x-r_i)\Bigr)\le\frac1{\tau^*}
  \le\max_{j\ge1}\bigl|S_j(P_uM)\bigr|
  \qquad\text{for every monic }M\in\R[x].
\end{equation}
At $M=1$ this is the construction of Proposition~\ref{prop:partial}.  A
witness $W$ with $\sup_{d\ge1}|w_d|=1/\Tail(v)$ proves that $v$ achieves
$\tau^*$.  Such a witness exists at every full minimiser.

\begin{proposition}[The minimum $\tau^*$ and a multiplier witness]\label{prop:attain}
Write $m=L-u$.  Then $\tau^*$ is a minimum, achieved by a full certificate $v$ on
$y_{u+1},\ldots,y_L$.  At any full certificate $v$ achieving $\tau^*$ there are reals $s_d$ with
$|s_d|\le1$, $s_d=\sgn(v_d)$ whenever $v_d\ne0$, and
$\sum_{d\ge1}s_dy_i^d=\tau^*$ for $u<i\le L$; then $w_0=1$,
$w_d=-s_d/\tau^*$ is a witness as in \eqref{eq:dualpair} achieving equality,
and it is $\widetilde P_u\widetilde M/(1-y)$ for a monic $M$, so
\eqref{eq:multiplier} is an equality at that $M$.
\end{proposition}

\begin{proof}
A certificate for $\varnothing$ on the surviving nodes is a
$\lambda\in\R^m$ with $\sum_i\lambda_i=1$, writing
$v_d(\lambda)=\sum_{i>u}\lambda_iy_i^d$; put $f(\lambda)=\Tail(v(\lambda))$.

\emph{The infimum is a minimum.}  Since $|v_d|\le\sum_i|\lambda_i|y_i^d$ and
the nodes lie in $(0,1)$, $f$ is finite, and it is homogeneous and
subadditive.  If $f(\lambda)=0$ then $v_d=0$ for $1\le d\le m$, a Vandermonde
system in the $\lambda_i$ with distinct nonzero nodes, so $\lambda=0$.  Thus
$f$ is a norm, hence continuous and coercive, and it has a minimum on
the closed nonempty affine set $\{\sum_i\lambda_i=1\}$; that minimum is
$\tau^*$.

\emph{Some minimiser is full.}  Let $\lambda$ minimise, let
$Z=\{d\ge1:v_d=0\}$, and suppose $|Z|\le m-2$.  The conditions
$\sum_ih_i=0$ and $v_d(h)=0$ for $d\in Z$ are at most $m-1$ linear
conditions on $h\in\R^m$, so some $h\ne0$ meets them.  As $v_d(h)=0$ on $Z$,
\[
  \varphi(t)=f(\lambda+th)=\sum_{d\notin Z}\bigl|v_d+tv_d(h)\bigr|,
\]
and each term is differentiable at $t=0$ with difference quotients dominated
by the summable $|v_d(h)|$, so $\varphi$ is differentiable there and
minimality gives $\varphi'(0)=\sum_{d\notin Z}\sgn(v_d)v_d(h)=0$.
Subtracting that vanishing linear term,
\[
  \varphi(t)-\varphi(0)=\sum_{d\notin Z}
  \bigl(|v_d+tv_d(h)|-|v_d|-t\sgn(v_d)v_d(h)\bigr),
\]
all of whose terms are nonnegative, the $d$-th vanishing exactly when
$v_d+tv_d(h)$ is zero or has the sign of $v_d$, that is when $t/t_d\le1$ for
$t_d=-v_d/v_d(h)$ (always, when $v_d(h)=0$).  Neither $v(\lambda)$ nor
$v(h)$ vanishes identically, so by the zero bound each has constant nonzero
sign past its largest zero, and hence $t_d$ has one sign for all large $d$.
Only finitely many $t_d$ have the other sign, all nonzero; let $t^*$ be the
one of them nearest $0$.  It exists, for otherwise every term would vanish
for all $t$ of that other sign and $\varphi$ would be constant on a
half-line, against coercivity.  For $t$ between $0$ and $t^*$ every $d$ has
$t/t_d\le1$, so every term above vanishes and
$\varphi(t^*)=\varphi(0)$: the certificate $\lambda+t^*h$ still minimises,
keeps every zero of $Z$, and has gained one.  The zero bound caps the zeros
at $m-1$, so iterating reaches a full minimiser.

\emph{The multipliers.}  Let $\lambda$ be a full minimiser with zero set
$Z$, $|Z|=m-1$, and set $s_d=\sgn(v_d)$ for $d\notin Z$.  A nontrivial real
combination of the $m$ distinct powers $y^d$, $d\in\{0\}\cup Z$, has at most
$m-1$ positive zeros, so the matrix $(y_i^d)_{i>u,\,d\in\{0\}\cup Z}$ is
nonsingular.  Hence there are unique reals $(s_d)_{d\in Z}$ and $\theta$
with $\sum_{d\ge1}s_dy_i^d=\theta$ for every surviving $i$, the series
converging because the $s_d$ are bounded.  Pairing with any $h$ with
$\sum_ih_i=0$ gives $\sum_{d\ge1}s_dv_d(h)=\theta\sum_ih_i=0$, while
minimality gives
\[
  0\le f'(\lambda;h)=\sum_{d\notin Z}\sgn(v_d)v_d(h)+\sum_{d\in Z}|v_d(h)|,
\]
so subtracting the pairing leaves
$0\le\sum_{d\in Z}\bigl(|v_d(h)|-s_dv_d(h)\bigr)$.  By the same
nonsingularity $h\mapsto(v_d(h))_{d\in Z}$ carries $\{\sum_ih_i=0\}$ onto
$\R^{m-1}$, so taking the preimages of $\pm\mathbf e_d$ gives $|s_d|\le1$ on $Z$.
Finally $\sum_{d\ge1}s_dv_d=\sum_i\lambda_i\sum_{d\ge1}s_dy_i^d=\theta$,
while $v_d=0$ on $Z$ makes the left side $\sum_{d\notin Z}|v_d|=\Tail(v)$,
so $\theta=\tau^*$.

\emph{The witness is a monic multiplier.}  Put $w_0=1$ and
$w_d=-s_d/\tau^*$.  Then $W(y_i)=1-\tau^*/\tau^*=0$ at every surviving node,
and $\sup_{d\ge1}|w_d|=1/\tau^*$ since $|s_d|=1$ off $Z$ and $|s_d|\le1$ on
it, so \eqref{eq:dualpair} is an equality.  Past $g=\max Z$ the sign of
$v_d$ is constant by the zero bound, so $w_d$ is constant there and
$(1-y)W(y)$ is a polynomial of degree at most $g+1$, vanishing at the $m$
surviving nodes and equal to $1$ at $0$.  Dividing out those nodes writes it
as $\widetilde P_uN$ with $N(0)=1$, and the reverse of $N$ is a monic $M$ of
degree at most $g+1-m$ with $N=\widetilde M$.
\end{proof}

\paragraph{Hermite spaces.}  For real roots
$1<r_1\le\cdots\le r_L$, repetitions allowed, $\mathcal H_k$ is the Hermite
space of the multiset $r_{k+1},\ldots,r_L$ as in
Proposition~\ref{prop:farrep}: the span of $d\mapsto d^az^d$, $z=1/r$ a
distinct value among them and $a$ below its multiplicity there, ordered by
growth.  Its dimension is $L-k$, and by the P\'olya--Szeg\H o zero bound a
nonzero element has at most $L-k-1$ real zeros counted with multiplicity.
A \emph{certificate for $X$ on $\mathcal H_k$} is a $v\in\mathcal H_k$ with
$v_0=1$ and $v|_X=0$; it is \emph{full} if it vanishes at $L-k-1$ positive
integers, its zero set.  A full certificate has no other real zero, is the
unique certificate with its zero set (unisolvence at distinct points),
changes sign at each point of it and keeps one sign past its largest zero; it
is a certificate for $X$ exactly when its zero set contains $X$.  For finite $X\subset\N_{>0}$ put
\[
  \tau_k(X)=\inf\{\Tail(v):v\text{ a certificate for }X\text{ on }\mathcal H_k\},
  \qquad \tau(X)=\tau_0(X),\qquad \tau^*=\tau_u(\varnothing).
\]
For distinct roots these are the quantities defined at the start of this section.  The top basis function $\psi_{k+1}$ of $\mathcal H_k$ is
$d^ay^d$, where $y=1/r_{k+1}$ is the largest node of $\mathcal H_k$ and
$a+1$ its multiplicity there; $\mathcal H_{k+1}$ is the span of the others.

\begin{theorem}[The placement game at repeated roots]\label{thm:valuerep}
Let $1<r_1\le\cdots\le r_L$, repetitions allowed, and $0\le u\le L-1$.  Then
\[
  \inf_Qb_{u+1}\Bigl(Q\prod_{i=1}^L(x-r_i)\Bigr)=\frac1T,
  \qquad T=\sup_{|X|=u}\tau(X),\qquad \tau^*\le T<\infty,
\]
over monic $Q\in\R[x]$, where $\tau(X)$ is the least tail of a certificate
for $X$ on $\mathcal H_0$ and $\tau^*$ the least tail of a $v\in\mathcal H_u$
with $v_0=1$.
\end{theorem}

\begin{proof}
$T<\infty$ by the first paragraph of the proof of Theorem~\ref{thm:value},
with $v^\infty$ a full certificate on $\mathcal H_p$ whose zero set
$Z$ contains $X_0$, which exists by unisolvence of $\mathcal H_p$ at the
$L-p$ distinct points $\{0\}\cup Z$, and Proposition~\ref{prop:farrep}(c),
which is stated for repeated roots.

The lower bound is \eqref{eq:weak}, which holds for every certificate on
$\mathcal H_0$ by the first paragraph of the proof of
Theorem~\ref{thm:converse}.

For the upper bound fix $X$ with $|X|=u$ and $D\ge\max(L,\max X)$, and let
$\tau_D(X)$ be the least $\sum_{d=1}^D|v_d|$ over certificates $v$ for $X$
on $\mathcal H_0$; it is positive because a nonzero element of
$\mathcal H_0$ has at most $L-1$ zeros.  Write $m_\rho$ for the multiplicity
of a distinct value $\rho$ among the $r_i$.  Since $\rho\ne0$,
$F=x^D+\sum_{j<D}f_jx^j$ is a multiple of $P$ exactly when
$\bigl((x\,d/dx)^aF\bigr)(\rho)=\sum_{j\le D}f_jj^a\rho^j=0$ for every
distinct $\rho$ and $a<m_\rho$, which are $L$ constraints.  Minimising the
largest $|f_j|$ over the positions $j<D$ with $D-j\notin X$ subject to these
constraints is a feasible linear program ($x^{D-L}P$ is feasible).  With
multipliers $\mu_{\rho,a}$ put
$v_d=-\sum_{\rho,a}\mu_{\rho,a}\rho^D(D-d)^a\rho^{-d}$; then the coefficient
of $f_j$ in the Lagrangian is $-v_{D-j}$, and as $\{(D-d)^a\}_{a<m_\rho}$ and
$\{d^a\}_{a<m_\rho}$ span the same polynomials, $\mu\mapsto v$ is a linear
bijection from $\R^L$ onto $\mathcal H_0$.  The dual program is to maximise $v_0$
subject to $v|_X=0$ and $\sum_{1\le d\le D,\,d\notin X}|v_d|\le1$, with value
$1/\tau_D(X)$, so some monic multiple of degree $D$ has
$b_{u+1}\le1/\tau_D(X)$.  The limit $\tau_D(X)\uparrow\tau(X)$ is the
compactness step of the proof of Theorem~\ref{thm:value}, using unisolvence
of $\mathcal H_0$ at $L$ distinct points; so the infimum is at most $1/T$.

For $T\ge\tau^*$ put $E_\nu=\{\nu+1,\ldots,\nu+u\}$ and suppose
$\tau(E_\nu)\le\tau^*-\varepsilon$ along a subsequence; take certificates
$v^\nu$ for $E_\nu$ on $\mathcal H_0$ with
$\Tail(v^\nu)\le\tau^*-\varepsilon/2$.  For $\nu\ge L$ their values at
$0,\ldots,L-1$ are bounded, so their coefficient vectors $c^\nu$ are bounded,
and $\|c^\nu\|_\infty$ is bounded below because $v^\nu_0=1$.
Proposition~\ref{prop:farrep}(a) with $k=u$, applied to
$v^\nu/\|c^\nu\|_\infty$, sends the coordinates of $c^\nu$ on
$\psi_1,\ldots,\psi_u$ to $0$.  A convergent subsequence therefore has a
limit $v\in\mathcal H_u$ with $v_0=1$, and by Fatou's lemma
$\Tail(v)\le\tau^*-\varepsilon/2$, against the definition of $\tau^*$.
\end{proof}

\begin{proposition}[The minimum $\tau^*$ at repeated roots]\label{prop:attainrep}
In the setting of Theorem~\ref{thm:valuerep}, $\tau^*$ is a minimum, achieved
by a full certificate on $\mathcal H_u$, and some monic $M\in\R[x]$ has
$\max_{j\ge1}|S_j(P_uM)|=1/\tau^*$; moreover
$\max_{j\ge1}|S_j(P_uM)|\ge1/\tau^*$ for every monic $M$.
\end{proposition}

\begin{proof}
Run the proof of Proposition~\ref{prop:attain} on $\mathcal H_u$, with the
coefficient vector $\lambda$ in the Hermite basis, the constraint $v_0=1$ in
place of $\sum_i\lambda_i=1$, and directions $h\in\mathcal H_u$ with $h_0=0$
(as sequences) in place of $\sum_ih_i=0$.  $\Tail$ is a
norm on $\mathcal H_u$ (an element vanishing at $1,\ldots,m$ vanishes), a
nonzero element has finitely many real zeros and so one sign for large $d$,
and the zero bound caps the zeros at $m-1$; so the first two steps hold
verbatim.  At a full minimiser with zero set $Z$, unisolvence at
$\{0\}\cup Z$ gives unique $(s_d)_{d\in Z}$ and $\theta$ with
$\sum_{d\ge1}s_dh_d=\theta h_0$ for all $h\in\mathcal H_u$ (with
$s_d=\sgn v_d$ off $Z$), and makes $h\mapsto(h_d)_{d\in Z}$ map
$\{h_0=0\}$ onto $\R^{m-1}$; so $|s_d|\le1$ and $\theta=\tau^*$ as there.
Put $w_0=1$, $w_d=-s_d/\tau^*$.  Then $\sum_dw_dh_d=0$ on $\mathcal H_u$;
with $h=d^ay_i^d$ this reads $((y\,d/dy)^aW)(y_i)=0$ for $a$ below the
multiplicity of $y_i$, so $W$ vanishes there to that order.  Past $\max Z$
the sign of $v$ is constant, $(1-y)W$ is a polynomial divisible by
$\widetilde P_u=\prod_{i>u}(1-r_iy)$ with value $1$ at $0$, and the quotient
is $\widetilde M$ for a monic $M$, whose partial sums $S_j(P_uM)$ are the
$w_j$.  For the last claim, the dual pairing \eqref{eq:dualpair} holds on
$\mathcal H_u$: if $\widetilde P_u$ divides $W$, then
$\sum_dw_dd^ay^d=((y\,d/dy)^aW)(y)$ vanishes at $y=1/r_i$ for $a$ below the
multiplicity of $r_i$ among $r_{u+1},\ldots,r_L$, so $\sum_dw_dv_d=0$ for
every $v\in\mathcal H_u$; with $W=\widetilde P_u\widetilde M/(1-y)$ this
gives $1\le\max_{j\ge1}|S_j(P_uM)|\,\tau^*$.
\end{proof}

The dual side settles sufficiency of the criterion at repeated roots, where
the construction of Proposition~\ref{prop:partial} is not available.

\begin{corollary}[The criterion at repeated roots]\label{cor:charrep}
Let $2\le r_1\le\cdots\le r_L$, repetitions allowed, and $0\le u\le L-1$.
Then \eqref{eq:order} is the infimum of $b_{u+1}$ over monic multiples of
$\prod_i(x-r_i)$ exactly when $|S_j(P_u)|\le|P_u(1)|$ for $1\le j<L-u$
(condition \eqref{eq:partial}); no alignment
$r_u<r_{u+1}$ is needed.
\end{corollary}

\begin{proof}
Necessity is Theorem~\ref{thm:converse}.  For sufficiency, put
$B=\prod_{i>u}(r_i-1)^{-1}=1/|P_u(1)|$.  By Theorem~\ref{thm:valuerep} the
infimum is $1/T$ with $T\ge\tau^*$, and by Proposition~\ref{prop:attainrep}
with $M=1$, $1/\tau^*\le\max_{j\ge1}|S_j(P_u)|$, which is $|P_u(1)|=1/B$
under \eqref{eq:partial}.  Hence $T\ge B$, while~\cite[Theorem~3.2]{coefficient-mass} gives
$1/T\ge1/B$.
\end{proof}

\subsection*{The infimum is never attained}

No monic multiple attains the infimum: a multiple of degree $D$ pairs only
with the first $D$ terms of a certificate, while each cost $\tau(X)$ is
achieved, by the argument of Proposition~\ref{prop:attain}, by a full
certificate, which has nonzero terms past every $D$.

\begin{lemma}[Full minimisers]\label{lem:attainX}
Let $0\le k<L$ and let $X$ be a set of at most $L-k-1$ positive integers.  Then
$\tau_k(X)$ is a minimum, achieved by a full certificate on $y_{k+1},\ldots,y_L$ whose
zero set contains $X$.
\end{lemma}

\begin{proof}
Run the first two steps of the proof of Proposition~\ref{prop:attain} on the
affine set $\{\lambda:\sum_i\lambda_i=1,\ v_x(\lambda)=0\ (x\in X)\}$ in place
of $\{\sum_i\lambda_i=1\}$.  It is closed, and nonempty because it contains the
full certificate of any zero set of size $L-k-1$ containing $X$, so the norm
$f$ has a minimum on it.  If a minimiser has a zero set $Z\supseteq X$
with $|Z|\le L-k-2$, the conditions $\sum_ih_i=0$ and $v_d(h)=0$ for $d\in Z$
are at most $L-k-1$ linear conditions on $h\in\R^{L-k}$, so some $h\ne0$ meets
them, and the move $\lambda+th$ of that proof keeps the tail and every zero of
$Z$, and so every point of $X$, until it gains a zero.  The zero bound caps
the zeros at $L-k-1$.
\end{proof}

\begin{lemma}[Full minimisers at repeated roots]\label{lem:attainXrep}
Let $1<r_1\le\cdots\le r_L$, $0\le k<L$, and let $X$ be a set of at most
$L-k-1$ positive integers.  Then $\tau_k(X)$ is a minimum, achieved by a full
certificate on $\mathcal H_k$ whose zero set contains $X$.
\end{lemma}

\begin{proof}
The proof of Lemma~\ref{lem:attainX}, with the changes listed in the proof
of Proposition~\ref{prop:attainrep}.
\end{proof}

\begin{theorem}[The infimum is never attained]\label{thm:noattain}
Let $1<r_1\le\cdots\le r_L$, repetitions allowed, and $0\le u\le L-1$.  Every
monic multiple $F$ of $P=\prod_i(x-r_i)$ has
\[
  b_{u+1}(F)>\frac1T=\inf_Qb_{u+1}(QP),
\]
with $T$ as in Theorem~\ref{thm:valuerep}.  In particular the
bound \eqref{eq:order}, whenever it is the infimum, is not attained.
\end{theorem}

\begin{proof}
Let $F$ have degree $D$ and let $X=X_F$ hold the distances $D-j$ of $u$
positions $j<D$ carrying its $u$ largest nonleading magnitudes.  The identity
behind \eqref{eq:weak} involves only $v_1,\ldots,v_D$: for a certificate $v$
for $X$ on $\mathcal H_0$, $0=\sum_{j\le D}f_jv_{D-j}$ (first paragraph of
the proof of Theorem~\ref{thm:converse}) gives
\[
  1\le b_{u+1}(F)\sum_{d=1}^{D}|v_d| .
\]
By Lemma~\ref{lem:attainXrep} with $k=0$, $\tau(X)$ is achieved by a
full certificate $v$.  It is a nonzero element of $\mathcal H_0$, so it has
at most $L-1$ real zeros and $v_d\ne0$ for some $d>D$.  Hence
$\sum_{d\le D}|v_d|<\Tail(v)=\tau(X)\le T$, and $b_{u+1}(F)>1/T$.  The last
sentence follows because the infimum is $1/T$.
\end{proof}

\section{The escaping placement is extremal}\label{sec:extremal}

This section answers the question left by Theorem~\ref{thm:value}.  For roots
at least $2$ the escaping placement is the worst one, $T=\tau^*$, repeated
roots included (Theorem~\ref{thm:extremal} and
Corollary~\ref{cor:extremalrep}); below $2$ it need not be
(Proposition~\ref{prop:belowtwo}).  Keep the notation of
Section~\ref{sec:value}: $1<r_1<\cdots<r_L$ are distinct, $y_i=1/r_i$, and
$\tau_k(X)$ is the least tail of a certificate for $X$ on
$y_{k+1},\ldots,y_L$.  The
proof inserts the exempt coefficients one at a time, each paid for by one
node, and the payment is the deletion step of~\cite[Theorem~2.2]{coefficient-mass} read backwards.

\begin{lemma}[The deletion step]\label{lem:delete}
Let $c\ge2$, let $\tfrac12\ge y_1>y_2>\cdots>y_c>0$ be any nodes, let $Z$ be a set of $c-1$
positive integers with $z=\max Z$, let $w$ be the full certificate on
$y_1,\ldots,y_c$ with zero set $Z$, and let $V$ be the full certificate on
$y_2,\ldots,y_c$ with zero set $Z\setminus\{z\}$.  Then
$\Tail(w)\le\Tail(V)$.
\end{lemma}

\begin{proof}
If $Z=\{1,\ldots,c-1\}$, the evaluation~\cite[Lemma~2.1]{coefficient-mass} gives both tails as products, which
differ by the factor $y_1/(1-y_1)\le1$.  Otherwise
$Z\ne\{1,\ldots,c-1\}$, and the displaced step of
Section~\ref{sec:extend}, with the nodes listed in increasing order, deletes
the largest prescribed zero $z$ together with the largest node $y_1$: its $V$
is the present one, and it gives
$\Tail(w)\le\Tail(V)-|V_z|\Sigma_G<\Tail(V)$.
\end{proof}

Inserting a zero $g$ below the largest zero of the certificate being lifted
is not covered by Lemma~\ref{lem:delete}, and it can fail: the full
certificate on $(\tfrac13,\tfrac15,\tfrac17)$ with zero set $\{1,3\}$ has tail
$\tfrac{31}{1704}$, while adding the node $\tfrac12$ and the zero $2$ gives the
consecutive zero set $\{1,2,3\}$ and, by~\cite[Lemma~2.1]{coefficient-mass}, the tail
$\tfrac1{48}>\tfrac{31}{1704}$.  The remedy is to insert $g$ and shift every
zero above it up by one, and to use that the certificate being lifted is a
minimiser, comparing it with the \emph{neighbours} of its zero set, in which
one run of zeros above $g$ moves up by one.  For an exponential sum $V$ we
also write $V(d)$ for $V_d$.

\begin{proposition}[The run chord]\label{prop:runchord}
Let $0\le k\le L-3$ and $y_{k+1}\le\tfrac12$, let $g\ge1$, let
$B_0\subseteq\{1,\ldots,g-1\}$, and let $A$ be a nonempty set of integers
above $g$ with $|B_0|+|A|=L-k-2$.  Write $A$ as the disjoint union of its
maximal runs of consecutive integers $[s_1,e_1],\ldots,[s_p,e_p]$, put
$A+1=\{a+1:a\in A\}$, and let the neighbour of $B_0\cup A$ at the $j$th run
be
\[
  Z_j=\bigl((B_0\cup A)\setminus\{s_j\}\bigr)\cup\{e_j+1\}.
\]
Let $V$ and $V_j$ be the full certificates on
$y_{k+2},\ldots,y_L$ with zero sets $Z=B_0\cup A$ and $Z_j$, and let $w$ be
the full certificate on $y_{k+1},\ldots,y_L$ with zero set
$\widehat Z=B_0\cup\{g\}\cup(A+1)$.  Put $\rho_j=|w(s_j)|/|V_j(s_j)|$ (the
denominator is nonzero), $\mu=1/(1+\sum_j\rho_j)$ and $\lambda_j=\mu\rho_j$,
so that $\lambda_j,\mu>0$ and $\sum_j\lambda_j+\mu=1$.  Then
\begin{equation}\label{eq:runchord}
  \Tail(V)\ge\sum_{j=1}^p\lambda_j\Tail(V_j)+\mu\,\Tail(w),
\end{equation}
with equality if and only if $y_{k+1}=\tfrac12$ and $B_0=\{1,\ldots,g-1\}$.
\end{proposition}

\begin{proof}
\emph{Signs.}  For an integer $d\ge0$ put $\varphi(d)=(-1)^{|Z\cap(0,d)|}$.
Each of $V$, $V_j$, $w$ is full, so by the zero bound its only real zeros
are the points of its zero set, all simple, and it is positive at $0$; its sign at an
integer $d$ outside its zero set is $-1$ to the number of its zeros below $d$.
Hence $V(d)=\varphi(d)|V(d)|$.  For an integer $d$,
$|Z_j\cap(0,d)|-|Z\cap(0,d)|=[e_j+1<d]-[s_j<d]$, which vanishes unless
$s_j<d\le e_j+1$, and those $d$ lie in $Z_j$; so $V_j(d)=\varphi(d)|V_j(d)|$
for every integer $d\ge0$.  Since $|(A+1)\cap(0,d)|=|A\cap(0,d-1)|$,
\[
  |\widehat Z\cap(0,d)|=|Z\cap(0,d)|+[g<d]-[d-1\in A],
\]
and $d-1\in A$ puts $d$ in $\widehat Z$; so
\begin{equation}\label{eq:runsigns}
  w(d)=\varphi(d)|w(d)|\quad(0\le d\le g),\qquad
  w(d)=-\varphi(d)|w(d)|\quad(d>g).
\end{equation}
In words: the $V_j$ keep the signs of $V$ at every integer, and $w$ keeps
them up to $g$ and reverses them above.

\emph{The combination.}  At a run start $s_i$ we have $V_j(s_i)=0$ for
$j\ne i$, as $s_i\in Z_j$; $V_i(s_i)\ne0$, as $s_i\notin Z_i$; and
$w(s_i)\ne0$, as $s_i>g$ and $s_i-1\notin A$ keep $s_i$ out of $\widehat Z$.  With the weights of the statement put
\[
  q=\sum_{j=1}^p\lambda_jV_j+\mu w .
\]
Then $q(0)=1$, and by \eqref{eq:runsigns}
$q(s_i)=\varphi(s_i)\bigl(\lambda_i|V_i(s_i)|-\mu|w(s_i)|\bigr)=0$.  Every
other point of $Z$ lies in $B_0$ or in $A\cap(A+1)$, hence in every $Z_j$ and
in $\widehat Z$, so $q$ vanishes on all of $Z$.  Let $\omega$ be the weight of $w$ on
$y_{k+1}$.  It is nonzero, for otherwise $w$ would be a sum on $L-k-1$ nodes
with the $L-k-1$ zeros $\widehat Z$; and since $y_{k+1}$ is the largest node and $w$
has the sign $(-1)^{|\widehat Z|}=(-1)^{L-k-1}$ past $\max\widehat Z$,
$\sgn\omega=(-1)^{L-k-1}$.  The sum $G=V-q$ on $y_{k+1},\ldots,y_L$ has weight
$-\mu\omega\ne0$ on $y_{k+1}$ and vanishes on the $L-k-1$ points
$\{0\}\cup Z$, so by the zero bound these are its only zeros and all are
simple.  For large $d$ its sign is $-\sgn\omega=(-1)^{L-k}=(-1)^{|Z|}$, which
is $\varphi(d)$ there; as $G$ and $\varphi$ change sign exactly at the points
of $Z$,
\begin{equation}\label{eq:runG}
  \varphi(d)G(d)=|G(d)|\qquad(d\ge0).
\end{equation}

\emph{Decomposition.}  Put $\delta_d=|V(d)|-|q(d)|$ and
$\gamma_d=\sum_j\lambda_j|V_j(d)|+\mu|w(d)|-|q(d)|$, so that, all the series
converging absolutely, the left side of \eqref{eq:runchord} minus the right is
$\sum_{d\ge1}(\delta_d-\gamma_d)$.  For $1\le d\le g$ every term of $q(d)$ has
the sign $\varphi(d)$ or vanishes, so $\gamma_d=0$ and
$\delta_d=\varphi(d)(V(d)-q(d))=\varphi(d)G(d)$.  For $d>g$ write
$\alpha=\sum_j\lambda_j|V_j(d)|$ and $\beta=\mu|w(d)|$; by \eqref{eq:runsigns}
$q(d)=\varphi(d)(\alpha-\beta)$, so $\gamma_d=\alpha+\beta-|\alpha-\beta|$ and
\[
  \delta_d-\gamma_d=|V(d)|-\alpha-\beta=\varphi(d)G(d)-2\mu|w(d)| ,
\]
whatever the sign of $\alpha-\beta$.  Hence the difference is
$\sum_{d\ge1}\varphi(d)G(d)-2\mu\sum_{d>g}|w(d)|$.  Since $G(0)=0$, and
$\varphi(d)\ne\varphi(d-1)$ only when $d-1\in Z$, where $G$ vanishes,
$\sum_{d\ge1}\varphi(d)G(d)=\sum_{d\ge1}\varphi(d)G(d-1)$.  With
\eqref{eq:runsigns} and \eqref{eq:runG} this gives
\begin{equation}\label{eq:runsplit}
  \Tail(V)-\sum_j\lambda_j\Tail(V_j)-\mu\Tail(w)
  =\sum_{d=1}^{g}|G(d-1)|+\sum_{d>g}\varphi(d)E(d),
  \qquad E(d)=G(d-1)+2\mu\,w(d),
\end{equation}
and it remains to show $\varphi(d)E(d)\ge0$ for every integer $d>g$.

\emph{Zeros of $E$.}  $E$ is an exponential sum on $y_{k+1},\ldots,y_L$, and
its weight on $y_{k+1}$ is $\mu\omega(2-1/y_{k+1})$.  It vanishes at every $d$
with $d-1\in\{0\}\cup Z$ and $d\in\widehat Z$.  With $\Sigma_0=\{0\}\cup B_0\cup\{g\}$ that
covers every $d$ with $d-1,d\in\Sigma_0$, and every $d\in A+1$.  Let $\Sigma_0$ consist of
the maximal runs $R_1,\ldots,R_t$ of consecutive integers, $0\in R_1$ and
$g\in R_t$.  At a gap start $h\notin\Sigma_0$ with $h-1\in\Sigma_0$ and $h-1<g$, we
have $G(h-1)=0$ and $h\notin\widehat Z$, so $E(h)=2\mu w(h)\ne0$, of sign
$\varphi(h)$ by \eqref{eq:runsigns}.

\emph{Counting.}  Suppose first $t\ge2$, with gap starts
$h_1<\cdots<h_{t-1}$, and $E\not\equiv0$.  The run $R_1$ supplies $|R_1|-1$
zeros below $h_1$; for $1\le j\le t-2$ the run $R_{j+1}\subseteq B_0$
supplies $|R_{j+1}|-1$ zeros in $(h_j,h_{j+1})$, and since
$\varphi(h_{j+1})=\varphi(h_j)(-1)^{|R_{j+1}|}$ that interval holds at least
$|R_{j+1}|$ zeros counted with multiplicity; and beyond $h_{t-1}$ lie the
known zeros $a+1,\ldots,g$, with $a=\min R_t$, and the $|A|$ points of $A+1$.
In all $E$ has at least $|\Sigma_0|-2+|A|=|B_0|+|A|=L-k-2$ zeros, of which at least
$|R_t|-1+|A|$ lie in $(h_{t-1},\infty)$.  At $h_{t-1}$ the sign of $E$ is
$\varphi(h_{t-1})=(-1)^{|B_0|-|R_t|+1}$.  If $y_{k+1}<\tfrac12$ the weight on
$y_{k+1}$ is nonzero, $E$ has at most $L-k-1$ zeros, hence at most
$|R_t|+|A|$ in $(h_{t-1},\infty)$, and its sign at infinity is
$-\sgn\omega=(-1)^{|B_0|+|A|}$; the number of its zeros in
$(h_{t-1},\infty)$ is then congruent to $|R_t|-1+|A|$ modulo $2$, so it is
exactly $|R_t|-1+|A|$.  If $y_{k+1}=\tfrac12$ the weight vanishes, $E$ is a
sum on $L-k-1$ nodes with at most $L-k-2$ zeros, and again only the known
zeros occur.  Either way the zeros of $E$ beyond $h_{t-1}$ are exactly
$a+1,\ldots,g$ and the points of $A+1$, all simple.  So for an integer $d>g$
outside $A+1$, where $d-1\notin A$ and hence $|(A+1)\cap(0,d)|=|A\cap(0,d)|$,
\[
  \sgn E(d)=(-1)^{|B_0|-|R_t|+1}(-1)^{|R_t|-1}(-1)^{|A\cap(0,d)|}
  =(-1)^{|Z\cap(0,d)|}=\varphi(d),
\]
while $E$ vanishes on $A+1$.  If $t=1$, that is $B_0=\{1,\ldots,g-1\}$, the
known zeros $1,\ldots,g$ and $A+1$ number $L-k-1$.  For $y_{k+1}<\tfrac12$
they are all the zeros, simple, and counting down from the sign
$(-1)^{|B_0|+|A|}$ at infinity gives $\sgn E(d)=\varphi(d)$ for $d>g$ outside
$A+1$ as before; for $y_{k+1}=\tfrac12$ the bound $L-k-2$ forces
$E\equiv0$.

So $\varphi(d)E(d)\ge0$ for every $d>g$, and \eqref{eq:runsplit} proves
\eqref{eq:runchord}.  Its first term vanishes exactly when $G$ vanishes on
$\{1,\ldots,g-1\}$, that is when $B_0=\{1,\ldots,g-1\}$; then $E\equiv0$ if
$y_{k+1}=\tfrac12$, and otherwise $\varphi(g+1)E(g+1)>0$, as $g+1\notin A+1$.
\end{proof}

Only the new node is held to the cutoff, and only through the sign of
$2-1/y_{k+1}$; the nodes $y_{k+2},\ldots,y_L$ may be any distinct reals in
$(0,y_{k+1})$.  For $y_{k+1}>\tfrac12$ that sign flips and
\eqref{eq:runchord} with these weights can fail: with the nodes
$(\tfrac34,\tfrac13,\tfrac14,\tfrac1{10},\tfrac1{11})$, $B_0=\{1,2\}$, $g=3$
and $A=\{7\}$,
\begin{align*}
  \Tail(V)&=\tfrac{39343118507}{13623090657480}=0.002887\ldots,\\
  \sum_j\lambda_j\Tail(V_j)+\mu\Tail(w)
  &=\tfrac{1246797277705214239223}{416394790925719545961974}=0.002994\ldots,
\end{align*}
computed by exact rational arithmetic (the finite part up to the largest
zero, then geometric series, the sign past it fixed by the zero bound, as in
the proof of Theorem~\ref{thm:elevenvalue}) in
\texttt{test\_run\_chord\_control\_above\_cutoff}
(\texttt{tests/test\_attainment.py}).

\begin{theorem}[The escaping placement is extremal]\label{thm:extremal}
Let $2\le r_1<\cdots<r_L$ and $0\le u\le L-1$.  For every $0\le k<u$ and every
set $X$ of $u-k$ positive integers,
\begin{equation}\label{eq:step}
  \tau_k(X)\le\tau_{k+1}\bigl(X\setminus\{\max X\}\bigr).
\end{equation}
Consequently $T=\tau^*$, and
\[
  \inf_Qb_{u+1}\Bigl(Q\prod_{i=1}^L(x-r_i)\Bigr)=\frac1{\tau^*}
  =\min_M\max_{j\ge1}\bigl|S_j(P_uM)\bigr|
\]
over monic $Q,M\in\R[x]$.
\end{theorem}

\begin{proof}
Fix $k$ and $X$, and put $g=\max X$ and $X'=X\setminus\{g\}$, a set of
$u-k-1\le L-k-2$ integers.  By Lemma~\ref{lem:attainX} on
$y_{k+2},\ldots,y_L$, $\tau_{k+1}(X')$ is achieved by a full certificate $V$
whose zero set $Z^-\supseteq X'$ has $L-k-2$ elements.  Every node is at most
$\tfrac12$.

If $g\in Z^-$, put $Z^+=Z^-\cup\{N\}$ with any $N>\max Z^-$; if $Z^-$ is
empty or $g>\max Z^-$, put $Z^+=Z^-\cup\{g\}$.  Either way
$Z^+\supseteq X$ has $L-k-1$ elements, its largest element is the one added,
and the full certificate $v^+$ on $y_{k+1},\ldots,y_L$ with zero set $Z^+$ is a
certificate for $X$.  Lemma~\ref{lem:delete}, with the largest node $y_{k+1}$
and the largest zero deleted, gives
$\tau_k(X)\le\Tail(v^+)\le\Tail(V)=\tau_{k+1}(X')$.

Otherwise $g\notin Z^-$ and $g<\max Z^-$.  Put $B_0=Z^-\cap(0,g)$, which
contains $X'$ because every element of $X'$ is below $g$, and
$A=Z^-\cap(g,\infty)\ne\varnothing$; then $L-k-2=|Z^-|\ge1$, so $k\le L-3$.
For a run $[s,e]$ of $A$ the neighbour $(Z^-\setminus\{s\})\cup\{e+1\}$ has
$L-k-2$ elements and contains $X'$, as $s>g$, so its full certificate is a
certificate for $X'$ on $y_{k+2},\ldots,y_L$ and its tail is at least
$\tau_{k+1}(X')=\Tail(V)$.  In \eqref{eq:runchord} every $\Tail(V_j)$ is
therefore at least $\Tail(V)$, so
$\Tail(V)\ge(1-\mu)\Tail(V)+\mu\Tail(w)$ with $\mu>0$, that is
$\Tail(w)\le\Tail(V)$.  The zero set $B_0\cup\{g\}\cup(A+1)$ of $w$ contains
$X$, so $\tau_k(X)\le\Tail(w)\le\tau_{k+1}(X')$.

Chaining \eqref{eq:step} from $k=0$ upwards, each step deleting the largest
remaining element, gives $\tau(X)=\tau_0(X)\le\tau_u(\varnothing)=\tau^*$ for
every $X$ of size $u$, so $T\le\tau^*$, and Theorem~\ref{thm:value} gives the
reverse inequality and the value $1/\tau^*$.  Proposition~\ref{prop:attain}
makes \eqref{eq:multiplier} an equality at some monic $M$.
\end{proof}

For example, at $(2,3,5,7)$ with $u=1$ the surviving factor
$(x-3)(x-5)(x-7)$ has partial sums $1,-14,57,-48$, so \eqref{eq:partial}
fails and \eqref{eq:order} is not the infimum.  The full certificate on
$(\tfrac13,\tfrac15,\tfrac17)$ with zero set $\{1,3\}$ is
\[
  v_d=\tfrac1{142}\bigl(81\cdot3^{-d}-625\cdot5^{-d}+686\cdot7^{-d}\bigr):
\]
$v_0=1$, $v_1=v_3=0$ and $v_2=-\tfrac1{71}$, and $v_d>0$ for $d\ge4$ by the
zero bound, so its tail is
$\tfrac1{71}+\tfrac1{142}\bigl(\tfrac32-\tfrac54+\tfrac13\bigr)=\tfrac{31}{1704}$,
the geometric sums $\sum_{d\ge4}\rho^{-d}$ at $\rho=3,5,7$ giving the
bracket.  The multiplier $M=x+\tfrac9{62}$ gives the partial sums
$-\tfrac{859}{62},\tfrac{1704}{31},-\tfrac{2463}{62},-\tfrac{1704}{31}$ of
$P_1M$, constant from $j=4$ on, so
$\max_{j\ge1}|S_j(P_1M)|=\tfrac{1704}{31}$, and by \eqref{eq:dualpair} both
achieve $\tau^*=\tfrac{31}{1704}$ and Theorem~\ref{thm:extremal} gives
\[
  \inf_Qb_2\bigl(Q(x-2)(x-3)(x-5)(x-7)\bigr)=\frac{1704}{31}=54.967\ldots,
\]
strictly between the bound $48$ of \eqref{eq:order} and the upper bound
$57=\max_{j\ge1}|S_j(P_1)|$ that the construction of
Proposition~\ref{prop:partial} gives.

\begin{corollary}[Repeated roots]\label{cor:extremalrep}
Let $2\le r_1\le\cdots\le r_L$, repetitions allowed, and $0\le u\le L-1$.
Then $T=\tau^*$ in Theorem~\ref{thm:valuerep}, and
\[
  \inf_Qb_{u+1}\Bigl(Q\prod_{i=1}^L(x-r_i)\Bigr)=\frac1{\tau^*}
  =\min_M\max_{j\ge1}\bigl|S_j(P_uM)\bigr|
\]
over monic $Q,M\in\R[x]$, with $\tau^*$ taken on the Hermite space
$\mathcal H_u$.
\end{corollary}

\begin{proof}
By Theorem~\ref{thm:valuerep} the infimum is $1/T\le1/\tau^*$, and
Proposition~\ref{prop:attainrep} gives the multiplier value, so it suffices
to show $b_{u+1}(F)\ge1/\tau^*$ for every monic multiple $F=QP$, where
$P=\prod_i(x-r_i)$.  Perturb the roots: if $r_i$ is the $(j+1)$st copy, in the
order of the indices, of a value $\rho$ of multiplicity $m_\rho$, put
$r_{i,\varepsilon}=\rho\exp(j\varepsilon)$ ($0\le j<m_\rho$).  For small
$\varepsilon>0$ these are distinct, at least $2$, and increasing in $i$, and
$P_\varepsilon=\prod_i(x-r_{i,\varepsilon})$ and $F_\varepsilon=QP_\varepsilon$
are monic with coefficients tending to those of $P$ and $F$; so
$b_{u+1}(F_\varepsilon)\to b_{u+1}(F)$.  Let $\tau^*_\varepsilon$ be $\tau^*$
at the roots $r_{i,\varepsilon}$; by Theorem~\ref{thm:extremal},
$b_{u+1}(F_\varepsilon)\ge1/\tau^*_\varepsilon$.

It remains to show $\limsup_{\varepsilon\to0}\tau^*_\varepsilon\le\tau^*$.  By
Proposition~\ref{prop:attainrep} some $v\in\mathcal H_u$ with $v_0=1$ has
$\Tail(v)=\tau^*$; write
$v_d=\sum_\rho\sum_{a<m'_\rho}c_{\rho,a}d^a\rho^{-d}$, where $m'_\rho$ is the
number of indices $i>u$ with $r_i=\rho$.  Those indices carry the $m'_\rho$
largest copies of $\rho$, so the surviving perturbed nodes at $\rho$ are
$z_{\rho,\varepsilon}\exp(-j\varepsilon)$, $0\le j<m'_\rho$, with
$z_{\rho,\varepsilon}=\rho^{-1}\exp\bigl(-(m_\rho-m'_\rho)\varepsilon\bigr)$.  Put
\[
  v^\varepsilon_d=\sum_\rho\sum_{a<m'_\rho}c_{\rho,a}\,
  z_{\rho,\varepsilon}^{\,d}\Bigl(\frac{1-\exp(-\varepsilon d)}{\varepsilon}\Bigr)^a .
\]
By the binomial expansion in the proof of Proposition~\ref{prop:confluent},
$v^\varepsilon$ is an exponential sum on the surviving perturbed nodes, and
$v^\varepsilon_0=\sum_\rho c_{\rho,0}=v_0=1$, so
$\tau^*_\varepsilon\le\Tail(v^\varepsilon)$.  As $\varepsilon\to0$,
$v^\varepsilon_d\to v_d$ for each $d$, and since
$z_{\rho,\varepsilon}\le\rho^{-1}$ and
$0\le(1-\exp(-\varepsilon d))/\varepsilon\le d$, every $|v^\varepsilon_d|$ is
at most $\sum_{\rho,a}|c_{\rho,a}|d^a\rho^{-d}$, which is summable; by
dominated convergence $\Tail(v^\varepsilon)\to\Tail(v)=\tau^*$.  Hence
$b_{u+1}(F)=\lim_\varepsilon b_{u+1}(F_\varepsilon)
\ge\limsup_\varepsilon1/\Tail(v^\varepsilon)=1/\tau^*$.
\end{proof}

Below $2$ the deletion step fails, and so can the conclusion.

\begin{proposition}[Below the cutoff]\label{prop:belowtwo}
Let $P=(x-\tfrac{11}{10})(x-3)(x-5)(x-7)$ and $u=1$.  Then
$\tau^*=\tfrac{31}{1704}$, while the monic multiple $F=QP$ with
$Q=x^4+\tfrac52x^3+\tfrac{11}6x^2+\tfrac{13}{12}x+\tfrac5{12}$,
\[
  F=x^8-\tfrac{68}5x^7+\tfrac{589}{12}x^6+\tfrac{433}{60}x^5
  -\tfrac{23863}{120}x^4+\tfrac{823}{20}x^3+\tfrac{997}{20}x^2
  +\tfrac{293}6x+\tfrac{385}8,
\]
has $b_2(F)=\tfrac{997}{20}<\tfrac{1704}{31}=1/\tau^*$.  Hence
$T>\tau^*$: the conclusion of Theorem~\ref{thm:extremal} fails at
$r_1=\tfrac{11}{10}$.
\end{proposition}

\begin{proof}
The surviving roots are those of the example after
Theorem~\ref{thm:extremal}, whose certificate and multiplier give
$\tau^*=\tfrac{31}{1704}$; that argument uses only \eqref{eq:dualpair},
which holds for any roots above $1$.  Expanding $QP$ gives the displayed
$F$; its nonleading coefficient magnitudes in decreasing order begin
$\tfrac{23863}{120}>\tfrac{997}{20}$, so $b_2(F)=\tfrac{997}{20}=49.85$.  By
Theorem~\ref{thm:value}, $1/T\le b_2(F)<1/\tau^*$.
\end{proof}

The conclusion of Lemma~\ref{lem:delete} fails at these roots, where
$y_1=\tfrac{10}{11}>\tfrac12$:
lifting the escaping minimiser, with zero set $\{1,3\}$, by the zero $5$
gives the full certificate on all four nodes with zero set $\{1,3,5\}$,
\[
  w_d=-\tfrac{5153632}{2063784541}\bigl(\tfrac{10}{11}\bigr)^d
  +\tfrac{3601989}{5381446}\,3^{-d}
  -\tfrac{17296875}{3682042}\,5^{-d}
  +\tfrac{6004901}{1193629}\,7^{-d},
\]
and tail $|w_2|+|w_4|+\bigl|\sum_{d\ge6}w_d\bigr|
=\tfrac{96935}{3398808}=0.02852\ldots>\tfrac{31}{1704}$ (by the zero bound
$w_d$ has one sign for $d\ge6$, so the last term is a sum of four
geometric series).  These weights are checked, with the tail, by exact
rational arithmetic in \texttt{test\_below\_two}
(\texttt{tests/test\_attainment.py}).
Section~\ref{sec:belowtwo} shows that $\{5\}$ is the worst placement there,
so the infimum is $\tfrac{3398808}{96935}$ (Theorem~\ref{thm:elevenvalue}).

\section{Below the cutoff}\label{sec:belowtwo}

Keep the notation of Section~\ref{sec:value}: $1<r_1<\cdots<r_L$ are
distinct, $y_i=1/r_i$, and $\tau_k(X)$ is the least tail of a certificate
for $X$ on $y_{k+1},\ldots,y_L$.  Nothing in this section assumes
$r_1\ge2$.  The game value $T$ is always a maximum, achieved possibly at a level
where some exempt coefficients have escaped
(Proposition~\ref{prop:attainT}); for one exempt coefficient it reduces to
$\tau^*$ and finitely many placements (Theorem~\ref{thm:halfmass}).  This
gives the exact infimum at $(\tfrac{11}{10},3,5,7)$
(Theorem~\ref{thm:elevenvalue}) and, along $(r,3,5,7)$, the exact range of
$r$ at which the escaping placement is extremal
(Theorem~\ref{thm:family}), and the infimum there down to $r=1.0745\ldots$
(Theorem~\ref{thm:familylow}) and its order as $r\to1^+$
(Theorem~\ref{thm:familyone}), with the constant oscillating
(Theorem~\ref{thm:familyosc} and Proposition~\ref{prop:familyrdelta});
every cost along the family is explicit (Theorem~\ref{thm:familyzero}),
with a limit profile (Theorem~\ref{thm:familylimit}) that settles the
transition at the jumps (Theorem~\ref{thm:familyjump}); at $(r,5,7)$ it
fails for every $1<r<2$
(Proposition~\ref{prop:cutoffsharp}).  Lower bounds on $T$ come from
multipliers with an exempt position.

\begin{lemma}[Pairing with exempt positions]\label{lem:exemptpair}
Let $1<r_1<\cdots<r_L$, $P=\prod_{i=1}^L(x-r_i)$, $X\subset\N_{>0}$ finite,
and $M\in\R[x]$ monic.  Every certificate $v$ for $X$ on $y_1,\ldots,y_L$
satisfies
\[
  1\le\max_{j\ge1,\,j\notin X}\bigl|S_j(PM)\bigr|\cdot\Tail(v),
\]
so $\tau(X)\ge1/\max_{j\ge1,\,j\notin X}|S_j(PM)|$.  If $|X|=u\le L-1$, then
$\inf_Qb_{u+1}(QP)\le\max_{j\ge1,\,j\notin X}|S_j(PM)|$ over monic $Q$.
\end{lemma}

\begin{proof}
Put $w_d=S_d(PM)$.  As in \eqref{eq:wgen}, $\sum_dw_dy^d$ is the reverse of
$PM$ divided by $1-y$, so it vanishes at every $y_i$, and $w$ is bounded,
equal to $(PM)(1)$ from $d=\deg PM$ on.  For $v_d=\sum_i\lambda_iy_i^d$,
absolute convergence gives
$\sum_{d\ge0}w_dv_d=\sum_i\lambda_i\sum_dw_dy_i^d=0$.  Since
$w_0=v_0=1$ and $v|_X=0$, $1=-\sum_{d\ge1,\,d\notin X}w_dv_d$, which gives
the inequality.  The last claim is Theorem~\ref{thm:value}:
the infimum is $1/T\le1/\tau(X)$.
\end{proof}

With $X=\varnothing$ and $P_u$ in place of $P$ this is \eqref{eq:dualpair}.

\begin{proposition}[The game value is a maximum]\label{prop:attainT}
Let $1<r_1<\cdots<r_L$ and $0\le u\le L-1$, and for $0\le k\le u$ put
$T_k=\sup\{\tau_k(X):X\subset\N_{>0},\ |X|=u-k\}$, so that $T_0=T$ and
$T_u=\tau^*$.  Then $T_0\ge T_1\ge\cdots\ge T_u$, and each $T_k$ is finite
and equals $\tau_l(X)$ for some $k\le l\le u$ and some $X$ with $|X|=u-l$.
In particular, if $u\ge1$ and $T>T_1$, then $T=\tau(X)$ for a placement $X$.
\end{proposition}

\begin{proof}
First, let $0\le k<L$, let $X_0$ be finite and $j\ge1$ with
$|X_0|+j\le L-k-1$, and let $E_\nu$ be $j$-sets disjoint from $X_0$ with
$\min E_\nu\to\infty$; we show $\tau_k(X_0\cup E_\nu)\to\tau_{k+j}(X_0)$.
Apply Proposition~\ref{prop:farrep} to the roots $r_{k+1},\ldots,r_L$, whose
spaces $\mathcal H_0$ and $\mathcal H_j$ there are the spans of
$y_{k+1}^d,\ldots,y_L^d$ and of $y_{k+j+1}^d,\ldots,y_L^d$.  By
Lemma~\ref{lem:attainX} some certificate $v^\infty$ for $X_0$ on
$y_{k+j+1},\ldots,y_L$ has tail $\tau_{k+j}(X_0)$, and part~(c) gives
certificates for $X_0\cup E_\nu$ on $y_{k+1},\ldots,y_L$ with tails tending
to it; so $\limsup_\nu\tau_k(X_0\cup E_\nu)\le\tau_{k+j}(X_0)$.
Conversely take certificates $v^\nu$ for $X_0\cup E_\nu$ on
$y_{k+1},\ldots,y_L$ with $\Tail(v^\nu)\le\tau_k(X_0\cup E_\nu)+1/\nu$,
along a subsequence realising the $\liminf$, so that the tails are bounded.
Fix $L-k$ points of $\{0,1,\ldots\}\setminus X_0$, one of them $0$; for
large $\nu$ they avoid $E_\nu$, the values of $v^\nu$ there are bounded, and
by the zero bound the coefficient vectors $c^\nu$ are bounded, while
$\|c^\nu\|_\infty\ge1/(L-k)$ because their entries sum to $v^\nu_0=1$.  Part~(a),
applied to $v^\nu/\|c^\nu\|_\infty$, sends the coefficients on
$y_{k+1},\ldots,y_{k+j}$ to $0$.  A convergent subsequence has a limit
$v$ on $y_{k+j+1},\ldots,y_L$ with $v_0=1$ and $v|_{X_0}=0$, and by Fatou's
lemma $\tau_{k+j}(X_0)\le\Tail(v)\le\liminf_\nu\tau_k(X_0\cup E_\nu)$.

For $k<u$ and $|X_0|=u-k-1$ this with $E_\nu=\{\nu\}$ gives
$\tau_{k+1}(X_0)=\lim_\nu\tau_k(X_0\cup\{\nu\})\le T_k$, so $T_{k+1}\le T_k$.
Now fix $k$ and take $(u-k)$-sets $X_\nu$ with $\tau_k(X_\nu)\to T_k$.
Passing to a subsequence, each coordinate of $X_\nu$, listed in increasing
order, is either constant or tends to $\infty$, and the constant ones come
first.  Let $X_0$ be the set of the constant ones and $E_\nu$ the rest, the
escaping part, so $X_\nu=X_0\cup E_\nu$ and either $E_\nu=\varnothing$ or
$|E_\nu|=j\ge1$ and $\min E_\nu\to\infty$.  In the first case $T_k=\tau_k(X_0)$; in the second
$|X_0|+j=u-k\le L-k-1$, and $T_k=\tau_{k+j}(X_0)$.  Both are finite, since
full certificates exist.  If $T>T_1$, the second case at $k=0$ would give
$T=\tau_j(X_0)\le T_j\le T_1$.
\end{proof}

For one exempt coefficient the escaped levels collapse to $\tau^*$, and an
exact identity for lifting by a zero above all the others bounds the
placements that can matter.

\begin{theorem}[One exempt coefficient]\label{thm:halfmass}
Let $1<r_1<\cdots<r_L$ with $L\ge2$, and $u=1$.  Let $V$ be a full
certificate on $y_2,\ldots,y_L$ achieving $\tau^*$
(Proposition~\ref{prop:attain}), $Z^*$ its zero set and
$g=\max(Z^*\cup\{0\})$, and let $H$ be the exponential sum on
$y_1,\ldots,y_L$ with weight $1$ on $y_1$, $H_0=0$ and $H|_{Z^*}=0$.
\begin{enumerate}
\item[(a)] $\tau(\{x\})\le\tau^*$ for $x\in Z^*$.
\item[(b)] For $x>g$, the full certificate $w$ on $y_1,\ldots,y_L$ with
zero set $Z^*\cup\{x\}$ (the \emph{aligned lift} of $V$ by $x$) satisfies
\[
  \Tail(w)=\tau^*-2\sum_{d>x}|V_d|
  -\frac{|V_x|}{|H_x|}\Bigl(\sum_{d=1}^{x}|H_d|-\sum_{d>x}|H_d|\Bigr).
\]
\item[(c)] Let $N$ be the least integer $N>g$ with
$2\sum_{d=1}^N|H_d|\ge\Tail(H)$ (it depends on the chosen $V$).  Then $\tau(\{x\})<\tau^*$ for every
$x\ge N$, and
\[
  T=\max\Bigl(\tau^*,\ \max\{\tau(\{x\}):1\le x<N,\ x\notin Z^*\}\Bigr),
\]
the inner maximum omitted when there is no such $x$.
\end{enumerate}
\end{theorem}

\begin{proof}
The matrix $(y_i^d)_{i\ge2,\,d\in\{0\}\cup Z^*}$ is square of size $L-1$ and
nonsingular by the zero bound, so $H$ exists, is unique and nonzero.  It has
the $L-1$ zeros $\{0\}\cup Z^*$, hence by the zero bound no others, all
simple; its weight on the largest node is positive, so $H_d>0$ for $d>g$.
Likewise $V$ has exactly the simple zeros $Z^*$ and $V_0=1$.  So for
integers $d\ge1$ outside $Z^*$ both $V_d$ and $H_d$ are nonzero, and since
both change sign exactly at the points of $Z^*$, $\sgn(V_dH_d)=\eta$ is
constant.

(a) $V$, with weight $0$ on $y_1$, is a certificate for $\{x\}$ on all the
nodes.

(b) Put $c=V_x/H_x$, which has sign $\eta$.  Then $w=V-cH$: it has
$w_0=1$ and vanishes on $Z^*\cup\{x\}$, and the full certificate is unique.
For $d\ge1$ outside $Z^*$, $cH_d$ has the sign of $V_d$, so
$w_d=\sgn(V_d)(|V_d|-|c||H_d|)$.  By the zero bound $w$ has exactly the
simple zeros $Z^*\cup\{x\}$, so, as $w_0>0$, $w_d$ has the sign of $V_d$ for
$d<x$ and the opposite sign for $d>x$.  Hence
\[
  \Tail(w)=\sum_{1\le d<x}\bigl(|V_d|-|c||H_d|\bigr)
  +\sum_{d>x}\bigl(|c||H_d|-|V_d|\bigr),
\]
and $\Tail(V)=\tau^*$ with $|V_x|=|c||H_x|$ gives the identity.

(c) As $x$ grows, $\sum_{d\le x}|H_d|$ increases to $\Tail(H)\in(0,\infty)$,
so $N$ exists and the bracket in (b) is nonnegative for every $x\ge N$.  As
$V$ has no zero beyond $g$, $\sum_{d>x}|V_d|>0$, so (b) gives
$\tau(\{x\})\le\Tail(w)<\tau^*$ for $x\ge N$.  Since $u=1$,
$T=\sup_x\tau(\{x\})\ge\tau^*$ by Theorem~\ref{thm:value}, and (a) leaves
only the $x<N$ outside $Z^*$.
\end{proof}

No bound on the nodes is used.  The identity in (b) is exact, so it also
decides whether the aligned lift costs more than $\tau^*$ at a given $x$;
at roots at least $2$, Lemma~\ref{lem:delete} gives $\Tail(w)\le\tau^*$ for
every $x>g$, and the placements $x<g$ outside $Z^*$ are those of the run
chord (Proposition~\ref{prop:runchord}).  Each remaining $\tau(\{x\})$ is
achieved by a full certificate (Lemma~\ref{lem:attainX}), but we do not
bound its zero set a priori, so the theorem reduces $T$ to finitely many
placements without making each of them a finite computation (along
$(r,3,5,7)$ Theorem~\ref{thm:familyzero} does).  Below, the
worst placement is pinned between an explicit full certificate and
Lemma~\ref{lem:exemptpair}, and the others are only bounded above.  For
$u\ge2$ we have no analogue of the threshold $N$.

\begin{theorem}[The infimum at $(\tfrac{11}{10},3,5,7)$]\label{thm:elevenvalue}
Let $P=(x-\tfrac{11}{10})(x-3)(x-5)(x-7)$ and $u=1$.  Then
$T=\tau(\{5\})=\tfrac{96935}{3398808}$, $\tau(\{x\})<T$ for every $x\ne5$,
and
\[
  \inf_Qb_2(QP)=\frac{3398808}{96935}=35.062\ldots
  =\max_{j\ge1,\,j\ne5}\bigl|S_j(PM)\bigr|,\qquad
  M=x^2+\tfrac{52552}{19387}x+\tfrac{69678}{19387}.
\]
\end{theorem}

\begin{proof}
Put $\theta=\tfrac{96935}{3398808}$.  The partial sums $S_1,\ldots,S_6$ of
the sextic $PM$ are
\[
  -\tfrac{2401917}{193870},\ \tfrac1\theta,\ \tfrac{6064861}{193870},\
  -\tfrac1\theta,\ -\tfrac{36840237}{96935},\ \tfrac1\theta,
\]
and $S_j=S_6$ for $j\ge6$, so Lemma~\ref{lem:exemptpair} with $X=\{5\}$ gives
$\tau(\{5\})\ge\theta$, while the full certificate with zero set
$\{1,3,5\}$ has tail $\theta$ (after Proposition~\ref{prop:belowtwo}).  So
$\tau(\{5\})=\theta$.  For the other placements apply
Theorem~\ref{thm:halfmass} with $V$ the full certificate on
$(\tfrac13,\tfrac15,\tfrac17)$ with zero set $Z^*=\{1,3\}$, which achieves
$\tau^*=\tfrac{31}{1704}<\theta$ (the example after
Theorem~\ref{thm:extremal}).  Then
\[
  H_d=\bigl(\tfrac{10}{11}\bigr)^d-\tfrac{14972607}{378004}\,3^{-d}
  +\tfrac{11210000}{94501}\,5^{-d}-\tfrac{30245397}{378004}\,7^{-d},
  \qquad \Tail(H)=\tfrac{2812589}{378004},
\]
and $2\sum_{d\le10}|H_d|<\Tail(H)\le2\sum_{d\le11}|H_d|$, so $N=11$, and
$\tau(\{x\})\le\tau^*<\theta$ for $x\in\{1,3\}$ and for $x\ge11$.  For
$x=2,4,6,7,8,9,10$ the full certificates with zero sets $\{1,2,10\}$,
$\{1,4,9\}$ and $\{1,3,x\}$ have the tails
\begin{align*}
  \{1,2,10\}&:\ \tfrac{843673514383499998315}{40515728086004858188032}=0.02082\ldots,\\
  \{1,4,9\}&:\ \tfrac{29343173972925245}{1145159451474217216}=0.02562\ldots,\\
  \{1,3,6\}&:\ \tfrac{177859285}{8265843096}=0.02151\ldots,\\
  \{1,3,7\}&:\ \tfrac{719323261}{37569613704}=0.01914\ldots,\\
  \{1,3,8\}&:\ \tfrac{181037572242979}{9826348204397256}=0.01842\ldots,\\
  \{1,3,9\}&:\ \tfrac{94783579348745045}{5199121229605850568}=0.01823\ldots,\\
  \{1,3,10\}&:\ \tfrac{199173967015849717825}{10949935864566580087416}=0.01818\ldots,
\end{align*}
all below $\theta=0.02852\ldots$.  So $T=\theta$, achieved
only at $\{5\}$, and Theorem~\ref{thm:value} gives the infimum.

The facts about $H$ and the seven tails are a computer-assisted
verification by exact rational arithmetic, in \texttt{test\_value\_eleven\_tenths}
(\texttt{tests/test\_attainment.py}); no step uses floating point, and the
decimals above are truncations of exact rationals.  The test solves the
$3\times3$ system for $H$ over $\mathbb Q$ and compares it with the display;
forms $\Tail(H)=\sum_{d\le3}|H_d|+\sum_{d\ge4}H_d$, the last sum a
combination of four geometric series since $H_d>0$ for $d>3$; and compares
$2\sum_{d\le n}|H_d|$ with it for $n=10,11$.  For each of the seven zero
sets $Z$ it solves the $4\times4$ system $w_0=1$, $w|_Z=0$ for the weights
on $(\tfrac{10}{11},\tfrac13,\tfrac15,\tfrac17)$, forms
$\Tail(w)=\sum_{d\le\max Z}|w_d|+\bigl|\sum_{d>\max Z}w_d\bigr|$ (one sign
past $\max Z$ by the zero bound, so the last sum is again geometric),
checks that it is the displayed rational, and checks $\Tail(w)<\theta$.  The partial sums of $PM$ and the tail of
the certificate $\{1,3,5\}$ are checked in the same way, and can be
verified by hand from the data displayed here and after
Proposition~\ref{prop:belowtwo}.
\end{proof}

The minimising zero set $\{1,3,5\}$ is $Z^*\cup\{5\}$, the aligned lift that
Lemma~\ref{lem:delete} would make cheaper than $\tau^*$ at $y_1\le\tfrac12$.
The window $g<x<N$ is needed here.  The aligned lift, with zero set
$\{1,3,x\}$, has tail $\tfrac{199}{4200}$ at $x=4$, $\theta$ at $x=5$, and
the values displayed in the proof at $x=6,\ldots,10$; against
$\tau^*=\tfrac{31}{1704}=0.018192\ldots$, it costs more than $\tau^*$ at
$x=4,\ldots,9$ and less at $x=10$, where its tail is $0.018189\ldots$.
These comparisons are also checked exactly in
\texttt{test\_value\_eleven\_tenths}.  By Theorem~\ref{thm:noattain} no
multiple attains the infimum.

\begin{theorem}[The family $(r,3,5,7)$]\label{thm:family}
Let $1<r<3$, $P_r=(x-r)(x-3)(x-5)(x-7)$ and $u=1$, so $\tau^*=\tfrac{31}{1704}$,
and let $r_c=(7+\sqrt{7395409})/1860=1.46583\ldots$ be the positive root of
$930r^2-7r-1988$.  Then $T>\tau^*$ for $1<r<r_c$, and $T=\tau^*$ for
$r_c\le r<3$.  More precisely, with $\beta(r)=48(r-1)(15r+71)/(49r-34)$,
\[
  \inf_Qb_2(QP_r)\ \begin{cases}
    \le\tfrac{209}{4}, & 1<r\le\tfrac{13}{10},\\[2pt]
    =\beta(r), & \tfrac{13}{10}\le r\le r_c,\\[2pt]
    =\tfrac{1704}{31}, & r_c\le r<3,
  \end{cases}
\]
over monic $Q$, where $\beta(r)<\tfrac{1704}{31}$ for $r<r_c$.  For
$\tfrac{13}{10}\le r<r_c$, moreover, $T=\tau(\{4\})=1/\beta(r)$ and
$\tau(\{x\})<T$ for every positive integer $x\ne4$; for
$\tfrac{29}{20}\le r<r_c$ even $\tau(\{x\})\le\tau^*$.
\end{theorem}

\begin{proof}
The surviving roots are $3,5,7$ for every $r$, so $\tau^*=\tfrac{31}{1704}$,
achieved by the full certificate on $(\tfrac13,\tfrac15,\tfrac17)$ with zero
set $Z^*=\{1,3\}$ (the example after Theorem~\ref{thm:extremal}).  Write
$y=1/r$ for the new node.

\emph{Lower bounds.}  For $M=x+\kappa$ the partial sums of $P_rM$ are
$S_1=\kappa-r-14$, $S_2=14r+57-\kappa(r+14)$, $S_3=\kappa(14r+57)-57r-48$,
$S_j=(P_rM)(1)=48(\kappa+1)(r-1)$ for $j\ge5$, and $S_4$, which we exempt.
Lemma~\ref{lem:exemptpair} with $X=\{4\}$ gives
$T\ge\tau(\{4\})\ge1/\max_{j\ne4}|S_j|$.  At $\kappa=\tfrac32$ each $S_j$ is
affine in $r$, so $|S_j|$ is convex and its maximum on $[1,\tfrac{13}{10}]$
is taken at an end; the values are $-\tfrac{27}2,\tfrac{97}2,\tfrac32,0$ at
$r=1$ and $-\tfrac{69}5,\tfrac{209}4,-\tfrac{93}{10},36$ at
$r=\tfrac{13}{10}$, all at most $\tfrac{209}4<\tfrac{1704}{31}$.  At
$\kappa=(105-34r)/(49r-34)$,
\[
  S_2=S_5=\beta(r),\qquad
  S_1=-\frac{7(7r^2+98r-83)}{49r-34},\qquad
  S_3=-\frac{3269r^2+882r-7617}{49r-34},
\]
where $\beta(r)>0$ for $r>1$.  Since $49r-34>0$ and $7r^2+98r-83>0$ for $r\ge1$ (the
latter is $22$ at $r=1$ and increasing), $S_1<0$, so $|S_1|\le\beta(r)$
amounts to $671r^2+2002r-2827\ge0$, while $|S_3|\le\beta(r)$ amounts to
$3989r^2+3570r-11025\ge0$ and $2549r^2-1806r-4209\le0$; the first two
increase on $r>0$ and the third is convex, and all three hold at
$r=\tfrac{13}{10}$ and $r=\tfrac32$, hence on $[\tfrac{13}{10},\tfrac32]$.
There $\tau(\{4\})\ge1/\beta(r)$, and
\[
  \frac1{\beta(r)}-\tau^*=\frac{1988+7r-930r^2}{3408(r-1)(15r+71)}
\]
has the sign of $1988+7r-930r^2$, positive exactly for $r<r_c$.

\emph{Upper bounds.}  For $1<r<3$, $y>\tfrac13$ is the largest node.
Solving the $4\times4$ systems, with the sign patterns given by the zero
bound (Appendix~\ref{sec:familydata} lists the solutions), the full
certificates with zero sets $\{1,3,4\}$ and $\{1,2,5\}$ have tails
\[
  \frac1{\beta(r)}\quad\text{and}\quad
  t_2(r)=\frac{4267r^2+5070r-5871}{48(r-1)(3466r^2+7455r+11025)},
\]
and $\tau^*-t_2(r)=c_2(r)/\bigl(3408(r-1)(3466r^2+7455r+11025)\bigr)$ with
$c_2(r)=214892r^3-55639r^2-138630r-266709$, so the second is at most
$\tau^*$ exactly when $c_2(r)\ge0$.  For the $H$ of
Theorem~\ref{thm:halfmass}, $H_d>0$ for $d\ge4$ and $H_2<0$ by the zero
bound, and solving the $3\times3$ system for $H$
(Appendix~\ref{sec:familydata}) gives, on $(1,3)$,
\[
  2\sum_{d\le5}|H_d|-\Tail(H)=-H_2+H_4+H_5-\sum_{d\ge6}H_d
  =\frac{c_5(r)(3-r)(5-r)(7-r)}{12524400\,r^5(r-1)},
\]
with $c_5(r)=206070r^3-67123r^2-161312r-238560$, which has the sign of
$c_5(r)$ there.  The two tails and this identity are identities between
rational functions of $r$, verified by exact polynomial arithmetic over
$\mathbb Q$ in \texttt{test\_family\_threshold}
(\texttt{tests/test\_attainment.py}): with $r$ an indeterminate the test
solves the three linear systems over $\mathbb Q(r)$, forms each tail and
the left side above from the sign pattern of the zero bound and the
geometric sums past the largest zero, and compares reduced numerators and
denominators with the displayed forms, so the check covers every $r$ in
$(1,3)$, not sample points; Appendix~\ref{sec:familydata} lists the
solutions and the values entering each tail, from which the identities can
also be checked by hand.  The same test checks the partial sums and the
three quadratics of the lower bounds as identities in $r$.  Both cubics
increase on $r\ge1$, since there $r\le r^2$ and $1\le r^2$ give
$c_2'(r)=644676r^2-111278r-138630\ge394768r^2>0$ and
$c_5'(r)=618210r^2-134246r-161312\ge322652r^2>0$, and both are positive
at $\tfrac{29}{20}$, so for $\tfrac{29}{20}\le r<3$ the certificate
$\{1,2,5\}$ gives $\tau(\{2\})\le\tau^*$ and $N\le5$.  By Theorem~\ref{thm:halfmass},
$T=\max(\tau^*,\tau(\{2\}),\tau(\{4\}))\le\max(\tau^*,1/\beta(r))$.

Now let $\tfrac{13}{10}\le r<\tfrac{29}{20}$, where $c_2$ and $c_5$ are
negative at the left end, and take one more term of $H$.  By the same sign
pattern and Appendix~\ref{sec:familydata},
\[
  2\sum_{d\le6}|H_d|-\Tail(H)=-H_2+H_4+H_5+H_6-\sum_{d\ge7}H_d
  =\frac{c_6(r)(3-r)(5-r)(7-r)}{1315062000\,r^6(r-1)}
\]
on $(1,3)$, with
$c_6(r)=23062470r^4-598283r^3-7874752r^2-16937760r-25048800$.  On $r\ge1$,
$c_6'(r)\ge(92249880-1794849-15749504-16937760)r^3=57767767r^3>0$, and
$c_6(\tfrac{13}{10})=\tfrac{522259242}{125}>0$, so $N\le6$.  The full
certificate with zero set $\{1,3,5\}$ has tail
\[
  t_5(r)=\frac{5(347r^2+250r-501)}{24(r-1)q_5(r)},\qquad
  q_5(r)=960r^2+5041r+7455,
\]
and, with $q(r)=3466r^2+7455r+11025$,
\[
  \frac1{\beta(r)}-t_2(r)=\frac{d_2(r)}{48(r-1)(15r+71)q(r)},\qquad
  \frac1{\beta(r)}-t_5(r)=\frac{d_5(r)}{16(r-1)(15r+71)q_5(r)},
\]
where $d_2(r)=105829r^3-131556r^2+14850r+41991$ and
$d_5(r)=-1670r^3-23167r^2+30517r+34080$.  For $r\ge1$,
$d_2'(r)=317487r^2-263112r+14850\ge54375r^2+14850>0$ and $d_2(1)=31114$, while
$d_5'(r)=-5010r^2-46334r+30517\le-20827$ and
$d_5(\tfrac32)=\tfrac{44187}2$; so $t_2(r)<1/\beta(r)$ and
$t_5(r)<1/\beta(r)$ on $[1,\tfrac32]$.  These three identities are checked
over $\mathbb Q(r)$ in \texttt{test\_family\_window}
(\texttt{tests/test\_attainment.py}), as those above.  By
Theorem~\ref{thm:halfmass}, as $\{1,3\}=Z^*$,
$T=\max(\tau^*,\tau(\{2\}),\tau(\{4\}),\tau(\{5\}))
\le\max(\tau^*,t_2(r),1/\beta(r),t_5(r))=1/\beta(r)$, since $\tau^*<1/\beta(r)$
for $r<r_c$.

For $r_c\le r<3$ the bound $T\le\max(\tau^*,1/\beta(r))$ for
$r\ge\tfrac{29}{20}$ is $\tau^*$, so $T=\tau^*$.  For
$\tfrac{29}{20}\le r<r_c$ the lower bound gives $T\ge\tau(\{4\})\ge1/\beta(r)>\tau^*$,
so $T=\tau(\{4\})=1/\beta(r)$, and every other placement costs at most
$\tau^*$: $\{1\}$ and $\{3\}$ by Theorem~\ref{thm:halfmass}(a), $\{2\}$ by
the certificate $\{1,2,5\}$, and $\{x\}$ with $x\ge5\ge N$ by
Theorem~\ref{thm:halfmass}(c).  For $\tfrac{13}{10}\le r<\tfrac{29}{20}$
likewise $T\ge\tau(\{4\})\ge1/\beta(r)$, so $T=\tau(\{4\})=1/\beta(r)$, and
every other placement costs less: $\{1\}$ and $\{3\}$ at most
$\tau^*<1/\beta(r)$, $\{2\}$ and $\{5\}$ at most $t_2(r)$ and $t_5(r)$, and
$\{x\}$ with $x\ge6\ge N$ less than $\tau^*$.  Theorem~\ref{thm:value} turns
these into the stated infima.
\end{proof}

Below $\tfrac{13}{10}$ the zero set minimising $\tau(\{4\})$ moves, and the
lower bound needs one multiplier for each zero set.  They are the
multipliers of Proposition~\ref{prop:attain}, for one placement.  Let $v$
be a full certificate on $y_1,\ldots,y_L$ whose zero set $Z$ contains $x$.
Put $s_d=\sgn(v_d)$ for $d\ge1$ outside $Z$, and let $(s_d)_{d\in Z}$ and
$\theta$ solve the $L$ linear equations $\sum_{d\ge1}s_dy_i^d=\theta$,
$1\le i\le L$, where the sum past $\max Z$ is geometric; the solution is
unique, the matrix $(y_i^d)_{i\le L,\,d\in\{0\}\cup Z}$ being nonsingular by
the zero bound.  Pairing with $v$
gives $\theta=\sum_i\lambda_i\sum_{d\ge1}s_dy_i^d=\sum_{d\ge1}s_dv_d=\Tail(v)$,
since $v$ vanishes on $Z$.  The sequence $w_0=1$, $w_d=-s_d/\theta$ is
bounded and its series vanishes at every $y_i$, which is all the proof of
Lemma~\ref{lem:exemptpair} uses (and, as at the end of the proof of
Proposition~\ref{prop:attain}, $w_d=S_d(PM)$ for a monic $M$).  So
\begin{equation}\label{eq:witness}
  \tau(\{x\})=\Tail(v)\qquad\text{whenever }|s_d|\le1\text{ for every }
  d\in Z\setminus\{x\}.
\end{equation}
For the placement $4$ with $Z=\{1,3,4\}$ along $(r,3,5,7)$ this is the
linear multiplier of the lower bound in the proof of
Theorem~\ref{thm:family}.

\begin{theorem}[The family below $\tfrac{13}{10}$]\label{thm:familylow}
Let $P_r=(x-r)(x-3)(x-5)(x-7)$ and $u=1$, and let $\beta$ and $t_5$ be as
in Theorem~\ref{thm:family} and its proof.  For $z=5,6,7,8$ let
$\theta_z(r)$ be the tail of the full certificate with zero set $\{1,4,z\}$
on $(\tfrac1r,\tfrac13,\tfrac15,\tfrac17)$, and put $\theta_3=1/\beta$, the
tail for $\{1,3,4\}$ (Appendix~\ref{sec:familydata} lists them).  Let
$\rho_1<\rho_2<\rho'<\rho_3<\cdots<\rho_7$ be the zeros in
$(1,\tfrac32)$ of the polynomials
$\gamma_1,\gamma_2,\gamma',\gamma_3,\ldots,\gamma_7$ of
Appendix~\ref{sec:familydata},
\begin{gather*}
  \rho_1=1.07451\ldots,\quad \rho_2=1.09547\ldots,\quad
  \rho'=1.11142\ldots,\quad \rho_3=1.12444\ldots,\quad
  \rho_4=1.13392\ldots,\\
  \rho_5=1.16696\ldots,\quad \rho_6=1.21576\ldots,\quad
  \rho_7=\tfrac{105\sqrt{4278}-1785}{3989}=1.27417\ldots .
\end{gather*}
Then for $\rho_1\le r\le\tfrac{13}{10}$
\[
  \inf_Qb_2(QP_r)=\frac1{T(r)},\qquad
  T(r)=\begin{cases}
    \theta_5(r), & \rho_1\le r\le\rho_2,\\
    t_5(r), & \rho_2\le r\le\rho_3,\\
    \theta_8(r), & \rho_3\le r\le\rho_4,\\
    \theta_7(r), & \rho_4\le r\le\rho_5,\\
    \theta_6(r), & \rho_5\le r\le\rho_6,\\
    \theta_5(r), & \rho_6\le r\le\rho_7,\\
    1/\beta(r), & \rho_7\le r\le\tfrac{13}{10}.
  \end{cases}
\]
More precisely, on $[\rho_1,\tfrac{13}{10}]$ the placement $\{5\}$ costs
$\min(\theta_5(r),t_5(r))$, which is $\theta_5(r)$ on $[\rho_1,\rho_2]$ and
$t_5(r)$ on $[\rho_2,\tfrac{13}{10}]$; on $[\rho',\tfrac{13}{10}]$ the placement
$\{4\}$ costs $\theta_8$, $\theta_7$, $\theta_6$, $\theta_5$, $1/\beta$ on
$[\rho',\rho_4]$, $[\rho_4,\rho_5]$, $[\rho_5,\rho_6]$, $[\rho_6,\rho_7]$,
$[\rho_7,\tfrac{13}{10}]$; the worst placement is $\{5\}$ for $r<\rho_3$,
$\{4\}$ for $r>\rho_3$, and both at $\rho_3$; and every other placement costs
less than $T(r)$.
\end{theorem}

With Theorem~\ref{thm:family} this gives the infimum on $[\rho_1,3)$; in
particular it is $\beta(r)$ on $[\rho_7,r_c]$, it is
$1/t_5(\tfrac{11}{10})=\tfrac{3398808}{96935}$ at $r=\tfrac{11}{10}$, as in
Theorem~\ref{thm:elevenvalue}, and $1/\theta_5(\tfrac54)=\tfrac{314024}{7627}=41.172\ldots$
at $r=\tfrac54$.  Consecutive pieces agree at the breakpoints: in the
notation of Appendix~\ref{sec:familydata},
$\theta_3-\theta_5=-(r+15)\gamma_7(r)/\bigl(96(r-1)(15r+71)Q_5(r)\bigr)$,
$t_5-\theta_5=-(r+15)\gamma_2(r)/\bigl(96(r-1)q_5(r)Q_5(r)\bigr)$,
$t_5-\theta_8=-\gamma_3(r)/\bigl(96(r-1)q_5(r)Q_8(r)\bigr)$, and $\theta_z-\theta_{z+1}$
for $z=5,6,7$ is a positive rational function times $-\gamma_6$, $-\gamma_5$,
$-\gamma_4$.

\begin{proof}
Write $Z^*=\{1,3\}$, as in the proof of Theorem~\ref{thm:family}.  Every
sign of a rational function of $r$ on an interval below is checked by
counting the real zeros of its numerator and denominator with Sturm
sequences over $\mathbb Q$ between rational endpoints; together with the
closed forms, this is the computer-assisted part of the proof, described
at its end.

\emph{Lower bounds.}  Apply \eqref{eq:witness} to the placement $x=4$ with
$Z=\{1,3,4\}$ and $Z=\{1,4,z\}$, $z=5,\ldots,8$, and to $x=5$ with
$Z=\{1,3,5\}$ and $Z=\{1,4,5\}$.  In all seven cases $|s_1|<1$ on
$[\tfrac{1074}{1000},\tfrac32]$.  At the remaining zero $c\in Z\setminus\{1,x\}$,
solving the $4\times4$ systems over $\mathbb Q(r)$ gives
\[
\begin{array}{cccccc}
  x & Z & D & D(1-s_c) & D(1+s_c) & |s_c|\le1\text{ on}\\[2pt]
  4 & \{1,4,5\} & 32(r-1)Q_5 & \gamma_6 & -35r\gamma_7 & [\rho_6,\rho_7]\\
  4 & \{1,4,6\} & 16(r-1)Q_6 & \gamma_5 & -105r\gamma_6 & [\rho_5,\rho_6]\\
  4 & \{1,4,7\} & 16(r-1)Q_7 & \gamma_4 & -105r\gamma_5 & [\rho_4,\rho_5]\\
  4 & \{1,4,8\} & 32(r-1)Q_8 & \gamma' & -105r\gamma_4 & [\rho',\rho_4]\\
  4 & \{1,3,4\} & 48(r-1)(15r+71) & 4209+1806r-2549r^2 & \gamma_7 & [\rho_7,\tfrac32]\\
  5 & \{1,3,5\} & 48(r-1)q_5 & -\gamma_+ & \gamma_2 & [\rho_2,\rho_+]\\
  5 & \{1,4,5\} & 96(r-1)Q_5 & -\gamma_2 & \gamma_1 & [\rho_1,\rho_2]
\end{array}
\]
with positive denominators $D$ on $r>1$.  Each $\gamma$ has exactly one
zero in $(1,\tfrac32)$ and is negative before it and positive after it, $\rho_+=1.4831\ldots$ being the zero of
$\gamma_+$, and $4209+1806r-2549r^2$ is concave and positive at $1$ and
$\tfrac32$.  So on $(1,\tfrac32]$, $|s_c|\le1$ exactly on the interval in the
last column, and
\eqref{eq:witness} gives $\tau(\{4\})$ on $[\rho',\tfrac{13}{10}]$ and
$\tau(\{5\})=\theta_5$ on $[\rho_1,\rho_2]$, $\tau(\{5\})=t_5$ on
$[\rho_2,\tfrac{13}{10}]$, as stated; as both certificates are certificates
for $\{5\}$, $\tau(\{5\})=\min(\theta_5,t_5)$ on $[\rho_1,\tfrac{13}{10}]$.

\emph{The placements $4$ and $5$.}  By the identity for $t_5-\theta_8$
above, $\theta_8<t_5$ for $1<r<\rho_3$ and $t_5<\theta_8$ for
$\rho_3<r<\tfrac32$; moreover $\theta_8<\theta_5$ on
$[\tfrac{1074}{1000},\tfrac{118}{100}]$, and $t_5<\theta_z$ on
$[\tfrac{11242}{10000},\tfrac{13}{10}]$ for $z=3,5,6,7$, where
$\tfrac{11242}{10000}<\rho_3$.  So $\tau(\{4\})\le\theta_8<\tau(\{5\})$ for
$\rho_1\le r<\rho_3$, and $\tau(\{5\})=t_5<\tau(\{4\})$ for
$\rho_3<r\le\tfrac{13}{10}$, where $\tau(\{4\})$ is one of
$\theta_8,\theta_7,\theta_6,\theta_5,\theta_3$.

\emph{The other placements.}  Apply Theorem~\ref{thm:halfmass} with the
escaping minimiser, $Z^*=\{1,3\}$.  With the signs $H_2<0<H_d$ ($d\ge4$) of
the proof of Theorem~\ref{thm:family},
$2\sum_{d\le n}|H_d|-\Tail(H)=-H_2+\sum_{d=4}^nH_d-\sum_{d>n}H_d$ is a
rational function of $r$, positive on $[\tfrac{1074}{1000},\tfrac{13}{10}]$
for $n=13$ and on $[b_n,\tfrac{13}{10}]$ for $n=6,\ldots,12$, where
$(b_6,\ldots,b_{12})=(1.274,1.197,1.153,1.125,1.106,1.092,1.081)$.  So
$N\le13$, and $N\le x$ once $r\ge b_x$.  The placements $\{1\}$ and $\{3\}$
cost at most $\tau^*$, and $\tau^*<\min(t_5,\theta_5)=\tau(\{5\})$ on
$[\tfrac{1074}{1000},\tfrac{13}{10}]$.  For $6\le x\le12$ and $r<b_x$ the
aligned lift, with zero set $\{1,3,x\}$, has tail below both $t_5$ and
$\theta_5$ on $[\tfrac{1074}{1000},b_x]$, so below $\tau(\{5\})$.  The
placement $\{2\}$ costs at most the tail of the full certificate
$\{1,2,13\}$, below $t_5$ and $\theta_5$ on
$[\tfrac{1074}{1000},\tfrac{1112}{1000}]$, and at most that of $\{1,2,10\}$,
below each of $\theta_3,\theta_5,\ldots,\theta_8$ on
$[\tfrac{1112}{1000},\tfrac{13}{10}]$, hence below $\tau(\{4\})$ there,
since $\rho'<\tfrac{1112}{1000}$.  So every placement other than $\{4\}$ and
$\{5\}$ costs less than $\max(\tau(\{4\}),\tau(\{5\}))$, and by
Theorem~\ref{thm:halfmass}(c) $T$ is that maximum, which the two previous
steps evaluate as displayed.  Theorem~\ref{thm:value} gives the infimum.

The computation is in \texttt{test\_family\_below}
(\texttt{tests/test\_attainment.py}), in the arithmetic of $\mathbb Q(r)$ of
\texttt{test\_family\_threshold}.  It solves the linear systems for the
certificates and for the witnesses with $r$ an indeterminate, checks the
closed forms of the $\theta_z$, the identities displayed above and in
Appendix~\ref{sec:familydata}, and that the series of each witness vanishes
at the four nodes; it brackets each zero $\rho$ between two rationals and
checks that each $\gamma$ has exactly one zero in $(1,\tfrac32)$; and it
checks every sign on an interval quoted in this proof by a Sturm count over
$\mathbb Q$.  The tails of the certificates with zero sets $\{1,3,x\}$ and
$\{1,2,z\}$ enter only through these signs.
\end{proof}

The pieces follow a pattern.  As $r$ decreases from $r_c$ the zero set
minimising $\tau(\{4\})$ moves up, $\{1,3,4\},\{1,4,5\},\ldots,\{1,4,8\}$.
For $\{1,4,z\}$ the entry $s_z$ of the witness passes from $-1$ at the
upper end of its piece to $1$ at the lower end, where, for $z\le7$,
$\{1,4,z+1\}$ takes over with the same tail; at $\rho_7$ the entries $s_3$ of $\{1,3,4\}$ and
$s_5$ of $\{1,4,5\}$ are both $-1$.  The cost of $\{5\}$ overtakes that
of $\{4\}$ at $\rho_3$.
At $r=\tfrac{11}{10}$ the worst placement $\{5\}$ and its zero set
$\{1,3,5\}$ are those of Theorem~\ref{thm:elevenvalue}, whose proof bounds
the other placements with other certificates.  Below $\rho_1$ we have no
closed form in $r$, but at each $r$ every cost is explicit
(Theorem~\ref{thm:familyzero}), and the zero sets above are its two
kinds, $\{1,x,z\}$ and $\{1,z,x\}$.  A floating-point search over zero sets in $\{1,\ldots,40\}$, at the
placements $x\le15$, suggests that the minimiser for $\{5\}$ continues
$\{1,5,6\},\{1,5,7\},\{1,5,8\},\ldots$, that the worst placement becomes $6$
near $r=1.03$ and $7$ near $r=1.01$, and that the half-mass point grows
($N=15$ at $r=1.06$).  The qualitative picture is proved below: infinitely
many pieces accumulate at $r=1$, the worst placement grows like
$\log_3\frac1{r-1}$, and the infimum tends to $24$
(Theorem~\ref{thm:familyone} and Corollary~\ref{cor:familyone}); off
neighbourhoods of the points where $3^x(r-1)=18$, the worst placement is
the largest $x$ with $3^x(r-1)<18$, with zero set $\{1,x,z\}$
(Theorem~\ref{thm:familyosc}), and near them it is that $x$ or $x+1$
(Theorem~\ref{thm:familyjump}).

\paragraph{The family near $r=1$.}  As $r\to1^+$ the node $\tfrac1r$ tends
to $1$, where $\lambda y^d$ is constant and has infinite tail, so a cheap
certificate gives it little weight, and the costs approach those at the
roots $3,5,7$.  No two nodes coalesce, so this is not the confluent
situation of Proposition~\ref{prop:confluent}; the limit $24$ is the
infimum at $(3,5,7)$ with $u=1$, where the surviving roots $5,7$ have
$\tau^*=\tfrac1{24}$ (Theorem~\ref{thm:extremal}, and the proof of
Proposition~\ref{prop:cutoffsharp}).  The placements at $(3,5,7)$ have
closed forms.  Below, $h_k$ is the complete homogeneous symmetric
polynomial of degree $k$ ($h_k=0$ for $k<0$); for distinct
$a_1,\ldots,a_k$ and $d\ge0$ the polynomial $p$ of degree below $k$ with
$p(a_i)=a_i^d$ satisfies
\begin{equation}\label{eq:interr}
  t^d-p(t)=h_{d-k}(a_1,\ldots,a_k,t)\prod_{i=1}^k(t-a_i),
\end{equation}
since the divided difference of $t^d$ at $a_1,\ldots,a_k,t$ is $h_{d-k}$.

\begin{lemma}[Placements at $(3,5,7)$]\label{lem:threenode}
For an integer $x\ge2$ put $A=3^{-x}$, $B=5^{-x}$, $C=7^{-x}$ and
$D=3A-10B+7C$.
\begin{enumerate}
\item[(a)] $D=\tfrac8{105}h_{x-2}(\tfrac13,\tfrac15,\tfrac17)>0$.  The full
certificate $U^{(x)}$ on $(\tfrac13,\tfrac15,\tfrac17)$ with zero set
$\{1,x\}$ has weights
\[
  \lambda_3=\frac{3(5B-7C)}{2D},\qquad
  \lambda_5=-\frac{5(3A-7C)}{2D},\qquad
  \lambda_7=\frac{7(3A-5B)}{2D},
\]
and its tail $t(x)$ satisfies
\[
  \frac1{24}-t(x)=\frac{5B-7C-45AB+84AC-35BC}{12D};
\]
$t(2)=\tfrac1{48}$ and $t(3)=\tfrac{31}{1704}$.
\item[(b)] Put
\[
  s_1=\frac{11A-60B+49C+210AB-280AC+70BC}{8D},\qquad
  s_x=\frac{1-36A+90B-56C}{6D},
\]
and $\sigma_1=s_1$, $\sigma_d=-1$ for $1<d<x$, $\sigma_x=s_x$, $\sigma_d=1$
for $d>x$.  Then $\sum_{d\ge1}\sigma_da^d=t(x)$ for
$a=\tfrac13,\tfrac15,\tfrac17$.
\item[(c)] For $x\ge4$: $0<s_1\le\tfrac12$, $0<s_x\le3^x/8$ and
$\tfrac1{24}-t(x)\ge\tfrac1{16}(\tfrac35)^x$.  For $x\ge6$ also
$\tfrac1{24}-t(x)\le\tfrac16(\tfrac35)^x$.
\item[(d)] The full certificate $W$ on $(\tfrac13,\tfrac15,\tfrac17)$ with
zero set $\{1,2\}$ has $W_d=h_{d-3}(\tfrac13,\tfrac15,\tfrac17)/105$ for
$d\ge3$ and tail $\tfrac1{48}$.  For $1<r<3$ and $x\ge3$ the full
certificate on $(\tfrac1r,\tfrac13,\tfrac15,\tfrac17)$ with zero set
$\{1,2,x\}$ has tail at most
\[
  \frac1{48}+\frac{315}{128}\cdot\frac{(r/3)^x}{r^2(r-1)} .
\]
\end{enumerate}
\end{lemma}

By (b), (c) and \eqref{eq:witness} at the roots $(3,5,7)$, $t(x)$ is the
cost of the placement $\{x\}$ there for $x\ge4$, and by (c) these costs
tend to $\tfrac1{24}$ from below at the rate $(\tfrac35)^x$.

\begin{proof}
Write $\mathbf a=(\tfrac13,\tfrac15,\tfrac17)$.

(a) By \eqref{eq:interr} at the nodes $\tfrac15,\tfrac17$ and
$t=\tfrac13$, the sum
$3^{-d}-\tfrac{10}3\,5^{-d}+\tfrac73\,7^{-d}$ equals
$(\tfrac13-\tfrac15)(\tfrac13-\tfrac17)h_{d-2}(\mathbf a)
=\tfrac8{315}h_{d-2}(\mathbf a)$; at $d=x$ it is $D/3$.  The weights are
checked by substituting them into $\sum_a\lambda_a=1$,
$\sum_a\lambda_aa=0$ and $\lambda_3A+\lambda_5B+\lambda_7C=0$.  By the zero
bound $U^{(x)}_d<0$ for $1<d<x$ and $U^{(x)}_d>0$ for $d>x$, so, with
$\sum_{d=2}^{x-1}a^d=(a^2-a^x)/(1-a)$,
\[
  t(x)=\sum_a\lambda_a\,\frac{a^x+a^{x+1}-a^2}{1-a},
\]
which is the stated closed form; at $x=2,3$ it gives $\tfrac1{48}$ and
$\tfrac{31}{1704}$, the latter as in the example after
Theorem~\ref{thm:extremal}.

(b) For $a\in(0,1)$,
$\sum_{d\ge1}\sigma_da^d=s_1a-(a^2-a^x)/(1-a)+s_xa^x+a^{x+1}/(1-a)$, and
at $a=\tfrac13,\tfrac15,\tfrac17$ the identity is checked by substitution,
$a^x$ being $A$, $B$, $C$.

(c) Put $\beta=(\tfrac35)^x=B/A$ and $\gamma=(\tfrac37)^x=C/A$, so
$\gamma=(\tfrac57)^x\beta$ and $D=A(3-10\beta+7\gamma)$.  For $x\ge4$,
$\beta\le(\tfrac35)^4$, $\gamma\le(\tfrac37)^4$, $(\tfrac57)^x\le(\tfrac57)^4$
and $A\le3^{-4}$, so $3-10\beta+7\gamma\ge3-10(\tfrac35)^4>\tfrac{17}{10}$.
The numerator of $s_x$ lies between $1-36\cdot3^{-4}-56\cdot7^{-4}>0$ and
$1+90\cdot5^{-4}$, and
$(1+90\cdot5^{-4})/(6(3-10(\tfrac35)^4))<\tfrac18$, so $0<s_x\le3^x/8$.
Dividing by $A$,
\begin{align*}
  8(3-10\beta+7\gamma)\,s_1&=11-60\beta+49\gamma+A(210\beta-280\gamma+70\beta\gamma),\\
  8(3-10\beta+7\gamma)(\tfrac12-s_1)&=1+20\beta-21\gamma-A(210\beta-280\gamma+70\beta\gamma),\\
  12(3-10\beta+7\gamma)(\tfrac1{24}-t(x))&=5\beta-7\gamma-A(45\beta-84\gamma+35\beta\gamma).
\end{align*}
The first is at least $11-60(\tfrac35)^4-280\cdot3^{-4}(\tfrac37)^4>0$, the
second at least
$\beta\bigl(20-21(\tfrac57)^4-210\cdot3^{-4}-70\cdot3^{-4}(\tfrac37)^4\bigr)>0$,
and the third at least
$\beta\bigl(5-7(\tfrac57)^4-45\cdot3^{-4}-35\cdot3^{-4}(\tfrac37)^4\bigr)$,
which with $3-10\beta+7\gamma\le3+7(\tfrac37)^4$ gives
$\tfrac1{24}-t(x)\ge\tfrac1{16}\beta$.  For $x\ge6$ the third is at most
$5\beta+84A\gamma\le\beta(5+84\cdot3^{-6}(\tfrac57)^6)$, and
$3-10\beta+7\gamma\ge3-10(\tfrac35)^6$ gives
$\tfrac1{24}-t(x)\le\tfrac16\beta$.

(d) By \eqref{eq:interr} at the nodes $\mathbf a$ and $t=0$, the value
$W_d=p(0)$, a combination of $3^{-d},5^{-d},7^{-d}$, has $W_0=1$,
$W_1=W_2=0$ and $W_d=\tfrac1{105}h_{d-3}(\mathbf a)>0$ for $d\ge3$, so
$\Tail(W)=\tfrac1{105}\sum_kh_k(\mathbf a)=\tfrac1{105}\prod_a(1-a)^{-1}=\tfrac1{48}$.
Let $y=\tfrac1r\in(\tfrac13,1)$.  By \eqref{eq:interr} at $t=y$, the
exponential sum $H'_d=y^d-p(y)$ on the four nodes, with weight $1$ on $y$,
vanishes at $d=0,1,2$ and equals
$h_{d-3}(\mathbf a,y)\prod_a(y-a)>0$ for $d\ge3$.  So
$w=W-(W_x/H'_x)H'$ is the full certificate with zero set $\{1,2,x\}$, and
$\Tail(w)\le\tfrac1{48}+(W_x/H'_x)\Tail(H')$.  Here
$h_k(\mathbf a)=3^{-k}h_k(1,\tfrac35,\tfrac37)\le3^{-k}/((1-\tfrac35)(1-\tfrac37))=\tfrac{35}8\,3^{-k}$,
so $W_x\le\tfrac1{24}3^{3-x}$; $H'_x\ge y^{x-3}\prod_a(y-a)$; and
$\Tail(H')=\prod_a(y-a)\sum_kh_k(\mathbf a,y)
=\prod_a(y-a)/\bigl((1-y)\prod_a(1-a)\bigr)$.  Hence
$(W_x/H'_x)\Tail(H')\le\tfrac1{24}3^{3-x}\cdot\tfrac{105}{48}\cdot
r^{x-3}\cdot\tfrac r{r-1}$, the stated bound.
\end{proof}

\begin{theorem}[The family near $r=1$]\label{thm:familyone}
Let $P_r=(x-r)(x-3)(x-5)(x-7)$ and $u=1$, let $\tau(\{x\})$ be the costs at
$P_r$ and $t(x)$ as in Lemma~\ref{lem:threenode}, and put
$\alpha=\log(5/3)/\log3=0.46497\ldots$.
\begin{enumerate}
\item[(a)] For $1<r<3$, $\inf_Qb_2(QP_r)>24$.
\item[(b)] For every $x\ge4$, $\tau(\{x\})\to t(x)$ as $r\to1^+$.
\item[(c)] For $1<r\le\tfrac{101}{100}$, with $\varepsilon=r-1$,
\[
  24+\tfrac32\,\varepsilon^\alpha\le\inf_Qb_2(QP_r)\le24+125\,\varepsilon^\alpha,
\]
and every worst placement $\{x\}$, that is $\tau(\{x\})=T$, satisfies
$\log_3\frac1\varepsilon-3<x<\log_3\frac1\varepsilon+6$.
\end{enumerate}
\end{theorem}

\begin{proof}
Let $y=\tfrac1r$.  The certificates $U^{(x)}$ and $W$ of
Lemma~\ref{lem:threenode}, with weight $0$ on $y$, are certificates on the
four nodes, so $\tau(\{x\})\le t(x)$ for $x\ge2$ and
$\tau(\{1\}),\tau(\{2\})\le\tfrac1{48}$.

(a) By Lemma~\ref{lem:threenode}, $\tfrac1{48}$, $t(3)=\tfrac{31}{1704}=\tau^*$
and $t(x)$ for $x\ge4$ are below $\tfrac1{24}$.  By
Theorem~\ref{thm:halfmass}(c), $T$ is the largest of $\tau^*$ and finitely
many $\tau(\{x\})$, so $T<\tfrac1{24}$, and Theorem~\ref{thm:value} gives
the claim.

\emph{The witness.}  We show: if $x\ge4$, $z_0\ge x+3$ is an integer, and
$1<r<3$ satisfies
\[
  \text{(W1)}\ \ y^x(1+y)-y^2\ge\tfrac{1-y}{24},\qquad
  \text{(W2)}\ \ y^2+2y^{z_0}-y^x-y^{x+1}\ge(1-y)\bigl(1+\tfrac{3^x}8\bigr),
\]
then $\tau(\{x\})\ge t(x)/(1+3^{6-z_0})$.  Let $\sigma$ be as in
Lemma~\ref{lem:threenode}(b) and $t=t(x)$.  For an integer $z\ge z_0$ and
$\varphi\in[0,1]$ let $\rho$ agree with $\sigma$ below $z$, with
$\rho_z=1-2\varphi$ and $\rho_d=-1$ for $d>z$.  As $\sigma_d=1$ for $d\ge z$,
$\sum_{d\ge1}\rho_da^d=t-\delta(a)$ at $a=\tfrac13,\tfrac15,\tfrac17$, where
$\delta(a)=2\varphi a^z+2a^{z+1}/(1-a)\in[0,2a^z/(1-a)]$.  Let
$c_1,c_2,c_3$ solve $\sum_{j\le3}c_ja^j=\delta(a)$ at the three nodes.  The
inverse of the matrix $(a^j)$, rows indexed by $j$, has rows
$(\tfrac{27}8,-\tfrac{125}4,\tfrac{343}8)$,
$(-\tfrac{81}2,\tfrac{625}2,-343)$,
$(\tfrac{945}8,-\tfrac{2625}4,\tfrac{5145}8)$, and as $z\ge7$ this gives
$|c_1|\le13\cdot3^{-z}$, $|c_2|\le146\cdot3^{-z}$, $|c_3|\le405\cdot3^{-z}$,
so $\sum_j|c_j|<\tfrac12$ and $|c_j|\le3^{6-z}$.  Put $\rho'_j=\rho_j+c_j$
for $j\le3$, positions below $x$, and $\rho'_d=\rho_d$ otherwise; then
$\sum_{d\ge1}\rho'_da^d=t$ at the three nodes.  Order the pairs
$(z,\varphi)$ by increasing $z$ and, for each $z$, decreasing $\varphi$;
$(z,0)$ and $(z+1,1)$ give the same $\rho$ and $c$, so along this path
$F=\sum_{d\ge1}\rho'_dy^d-t$ is continuous.  At $(z_0,1)$,
\[
  \sum_{d\ge1}\rho_dy^d=s_1y+s_xy^x-\frac{y^2+2y^{z_0}-y^x-y^{x+1}}{1-y},
\]
and $s_1\le\tfrac12$, $0<s_xy^x<3^x/8$, $\sum_j|c_j|<\tfrac12$ and (W2) give
$F<-t<0$.  As $z\to\infty$, $c\to0$ and
$F\to\sum_d\sigma_dy^d-t\ge(y^x+y^{x+1}-y^2)/(1-y)-t\ge\tfrac1{24}-t>0$ by
$s_1,s_x>0$ and (W1).  So $F=0$ at some $(z,\varphi)$ with $z\ge z_0$.  There
put $w_0=1$ and $w_d=-\rho'_d/t$ for $d\ge1$: $w$ is bounded and its series
$1-\tfrac1t\sum_d\rho'_da^d$ vanishes at all four nodes, which is all the
proof of Lemma~\ref{lem:exemptpair} uses.  Off the position $x$,
$|\rho'_1|\le\tfrac12+\tfrac12$, $|\rho'_2|,|\rho'_3|\le1+3^{6-z}$,
$|\rho'_z|\le1$ and $|\rho'_d|=1$ otherwise, so every certificate for
$\{x\}$ has tail at least $t/(1+3^{6-z})\ge t/(1+3^{6-z_0})$.

(b) Fix $x\ge4$ and $z_0\ge x+3$.  As $r\to1^+$ the left sides of (W1) and
(W2) tend to $1$ and the right sides to $0$, so
$t(x)/(1+3^{6-z_0})\le\tau(\{x\})\le t(x)$ for $r$ near $1$; and $z_0$ is
arbitrary.

(c) \emph{Upper bound.}  Let $x_r$ be the largest integer with
$3^{x_r}\le8/(r\varepsilon)-48$.  As $r\varepsilon\le\tfrac{101}{10000}<\tfrac8{777}$,
$x_r\ge6$; as $8/r-48\varepsilon\ge3$, $3^{x_r}>(8/(r\varepsilon)-48)/3\ge1/\varepsilon$,
so $(\tfrac35)^{x_r}=3^{-\alpha x_r}\le\varepsilon^\alpha$.  Take $x=x_r$ and
$z_0=x+3$.  The left side of (W2) is
$y^2-(1-y)y^x(1+2y+2y^2)\ge y^2-5(1-y)$, and
$y^2/(1-y)=1/(r\varepsilon)\ge6+3^x/8$, which gives (W2).  For (W1),
$x\log r\le\varepsilon\log_3(8/\varepsilon)\le\tfrac1{100}\log_3800<\tfrac1{14}$,
as $\varepsilon\log(8/\varepsilon)$ increases on $(0,\tfrac1{100}]$, so
$y^{x-2}\ge y^x>\tfrac{13}{14}$ and
$y^x(1+y)-y^2\ge y^2(\tfrac{13}{14}\cdot\tfrac{201}{101}-1)>\tfrac{1-y}{24}$.
So $\tau(\{x_r\})\ge t(x_r)/(1+3^{3-x_r})$.  By
Lemma~\ref{lem:threenode}(c), $24(\tfrac1{24}-t(x_r))\le4(\tfrac35)^{x_r}<\tfrac15$,
so $1/t(x_r)\le24+30\cdot4(\tfrac35)^{x_r}<30$, and
\[
  \inf_Qb_2(QP_r)\le\frac{1+3^{3-x_r}}{t(x_r)}
  \le24+120(\tfrac35)^{x_r}+810\cdot3^{-x_r}\le24+121(\tfrac35)^{x_r},
\]
as $810\cdot5^{-x_r}<1$; this is at most $24+125\varepsilon^\alpha$.

\emph{Lower bound.}  Let $X$ be the least integer $X\ge4$ with
$\tfrac{315}{128}(r/3)^{X+1}\le r^2\varepsilon/192$.  By
Lemma~\ref{lem:threenode}(d) every placement $x>X$ costs at most
$\tfrac1{48}+\tfrac1{192}=\tfrac5{192}$; the placements $x\le3$ cost at most
$\tfrac1{48}$ or $\tfrac{31}{1704}$, and $\tau^*=\tfrac{31}{1704}$; and
$4\le x\le X$ costs at most $t(x)\le\tfrac1{24}-\tfrac1{16}(\tfrac35)^X$.  As
$\tfrac1{16}(\tfrac35)^X\le\tfrac1{16}(\tfrac35)^4<\tfrac1{24}-\tfrac5{192}$,
$T\le\tfrac1{24}-\tfrac1{16}(\tfrac35)^X$, and
$\inf_Qb_2(QP_r)=1/T\ge24+36(\tfrac35)^X$.  If $X=4$, $(\tfrac35)^X>\tfrac1{23}$.
If $X>4$, minimality gives $(3/r)^X<\tfrac{945}{2r^2\varepsilon}\le\tfrac{945}{2\varepsilon}$,
so $X<\log\tfrac{945}{2\varepsilon}/\log\tfrac{300}{101}$; as
$\log r\le\varepsilon$, $\varepsilon\log\tfrac{945}{2\varepsilon}\le\tfrac1{100}\log47250<\tfrac{27}{250}$
(the left side increases on $(0,\tfrac1{100}]$) and $\log\tfrac{300}{101}>\tfrac{27}{25}$,
$X\log r<\tfrac1{10}$ and $r^X<e^{1/10}<\tfrac{111}{100}$, so $3^X<\tfrac{111}{100}\cdot\tfrac{945}{2\varepsilon}<\tfrac{525}\varepsilon$
and $(\tfrac35)^X=3^{-\alpha X}>(\varepsilon/525)^\alpha>\varepsilon^\alpha/23$,
as $\alpha<\tfrac12$ ($25<27$).  Either way $\inf_Qb_2(QP_r)\ge24+\tfrac{36}{23}\varepsilon^\alpha$.

\emph{The worst placement.}  As $r<r_c$, $T>\tau^*$
(Theorem~\ref{thm:family}), so $T=\tau(\{x\})$ for some $x$
(Proposition~\ref{prop:attainT}).  By the upper bound,
$T\ge\tau(\{x_r\})\ge(\tfrac1{24}-\tfrac16(\tfrac35)^6)/(1+\tfrac1{27})>\tfrac5{192}$,
so a worst placement has $4\le x\le X$; and $X<\log_3\tfrac{525}\varepsilon<\log_3\tfrac1\varepsilon+6$
when $X>4$, while $4<\log_3\tfrac1\varepsilon$.  Also
$T\ge1/(24+125\varepsilon^\alpha)\ge\tfrac1{24}-\tfrac{125}{576}\varepsilon^\alpha$
and $T=\tau(\{x\})\le t(x)\le\tfrac1{24}-\tfrac1{16}(\tfrac35)^x$, so
$3^{-\alpha x}=(\tfrac35)^x\le\tfrac{125}{36}\varepsilon^\alpha$ and
$x\ge\log_3\tfrac1\varepsilon-\log_3(\tfrac{125}{36})/\alpha>\log_3\tfrac1\varepsilon-3$,
as $\log_3\tfrac{125}{36}<\tfrac65$ and $\alpha>\tfrac25$
($(\tfrac53)^5>9$).
\end{proof}

The closed forms and inequalities of Lemma~\ref{lem:threenode}, the
constants $13,146,405$, the bound $(1+3^{3-x})/t(x)\le24+121(\tfrac35)^x$
for $6\le x\le60$, and the witness itself at $r=\tfrac{101}{100}$ and
$x=x_r=6$, where $F$ vanishes at $z=41$, are checked exactly in
\texttt{test\_family\_near\_one} (\texttt{tests/test\_attainment.py}).  The
witness needs $3^x$ of order at most $1/\varepsilon$: the entry $s_x$,
about $3^x/18$, must be balanced at the node $y$ by the flip, whose weight
there is at most $y^{x+1}/(1-y)$.  Beyond that the node $\tfrac1r$ makes
the placement cheap: at $r=\tfrac{101}{100}$ the certificate $\{1,2,9\}$
costs less than $t(9)/(1+3^{-6})$, which the same test checks as a control.

\begin{corollary}\label{cor:familyone}
Along $(r,3,5,7)$ with $u=1$, $\inf_Qb_2(QP_r)\to24$ as $r\to1^+$, and the
worst placement tends to $\infty$.  On no interval $(1,1+\delta)$ does
$T(r)$ agree at every $r$ with one of finitely many rational functions of
$r$.  In particular, as the tail of the full certificate with a given zero
set is a rational function of $r$ on $(1,3)$, infinitely many pieces as in
Theorem~\ref{thm:familylow} accumulate at $r=1$.
\end{corollary}

\begin{proof}
The first two claims are Theorem~\ref{thm:familyone}(c).  By it,
$\tfrac1{24}-T=(\inf_Qb_2-24)/(24\inf_Qb_2)$ lies between
$c\varepsilon^\alpha$ and $C\varepsilon^\alpha$ for some $c,C>0$ and all
$\varepsilon\le\tfrac1{100}$.  If $T$ took values among rational functions
$f_1,\ldots,f_k$ on $(1,1+\delta)$, one $f_j$ would equal $T$ at points
$r_n\to1^+$.  Then $f_j-\tfrac1{24}$ is a rational function, nonzero at the
$r_n$ and tending to $0$ along them, so it has a zero of some order $m\ge1$
at $1$ and is $O(r-1)$ there, against $\tfrac1{24}-T\ge c\varepsilon^\alpha$
with $\alpha<1$.  The tail is rational in $r$ because the weights solve a
linear system with entries $r^{-d}$, the zero bound fixes the sign pattern,
and the part past the largest zero is geometric.
\end{proof}

\paragraph{The constant near $r=1$.}  The constants $\tfrac32$ and $125$ of
Theorem~\ref{thm:familyone}(c) cannot be brought to a common value:
$(\inf_Qb_2(QP_r)-24)(r-1)^{-\alpha}$ oscillates log-periodically between
$80\cdot18^{-\alpha}$ and $\tfrac{400}3\cdot18^{-\alpha}$
(Theorem~\ref{thm:familyosc}).  The switch is at $3^x(r-1)=18$.  Below it
the flipped witness works with the true size of $s_x$, which is $3^x/18$ up
to a factor tending to $1$, and the placement $\{x\}$ costs $t(x)$ up to
$O(3^{-x})$; above it a small weight on $\tfrac1r$ added to a certificate
$U^{(c)}$ at $(3,5,7)$ makes the placement cheaper by an amount that does
not tend to $0$.  For $c\ge2$ let $\ell(c)$ and $\lambda_7(c)$ be the weights
$\lambda_3$ and $\lambda_7$ of the certificate $U^{(c)}$ of
Lemma~\ref{lem:threenode}(a), and put
\[
  K_c=\frac{\ell(c)}{\tfrac1{24}-t(c)},\qquad
  R(x)=\tfrac{36}5\bigl(\tfrac53\bigr)^x\bigl(\tfrac1{24}-t(x)\bigr).
\]

\begin{lemma}[Both sides of $3^x(r-1)=18$]\label{lem:eighteen}
Let $1<r<3$, $\varepsilon=r-1$, $y=\tfrac1r$, and let $\tau(\{x\})$ be the
costs at $P_r=(x-r)(x-3)(x-5)(x-7)$.
\begin{enumerate}
\item[(a)] Let $x\ge4$, and let $t=t(x)$, $s_1$ and $s_x$ be as in
Lemma~\ref{lem:threenode}.  If
\[
  y^x(1+y)-y^2\ge\frac{1-y}{24}\qquad\text{and}\qquad
  s_x+7\le\frac1{r\varepsilon},
\]
then for some integer $z\ge x+3$ the cost $\tau(\{x\})$ is the tail of the
full certificate with zero set $\{1,x,z\}$, and
$t(x)-14\cdot3^{-z}\le\tau(\{x\})\le t(x)$.
\item[(b)] For integers $c\ge2$ and $x>c$,
\[
  \tau(\{x\})\le t(c)+\frac{35}{24}\cdot
  \frac{\ell(c)3^{-x}+\lambda_7(c)7^{-x}}{h_{x-2}(\tfrac15,\tfrac17,1)}\cdot
  \frac{r^{x-1}}{\varepsilon},
\]
and the right side decreases as $x$ increases.
\item[(c)] As $x\to\infty$, $18\cdot3^{-x}s_x\to1$ and $R(x)\to1$; and
$R(x)\ge\tfrac9{20}$ for $x\ge4$.  For $c\ge2$, $K_c>18$, and $K_c\to18$ as
$c\to\infty$.
\end{enumerate}
\end{lemma}

\begin{proof}
Write $A=3^{-x}$, $B=5^{-x}$, $C=7^{-x}$ and $D=3A-10B+7C$, as in
Lemma~\ref{lem:threenode}.

(a) We run the witness of the proof of Theorem~\ref{thm:familyone} again,
correcting it at the positions $1$ and $x$ instead of $1,2,3$, so that no
entry off $x$ leaves $[-1,1]$.  Let $\sigma$ be as in
Lemma~\ref{lem:threenode}(b).  For an integer $z\ge x+3$ and
$\varphi\in[0,1]$ put $\delta(a)=2\varphi a^z+2a^{z+1}/(1-a)$, which lies in
$[0,2a^z/(1-a)]$, and let $p,q,\theta'$ solve
$pa+qa^x-\theta'=\delta(a)$ at $a=\tfrac13,\tfrac15,\tfrac17$.  Subtracting
the equations in pairs and solving,
\[
  q=\frac{3\delta(\tfrac13)-10\delta(\tfrac15)+7\delta(\tfrac17)}{D},\qquad
  p=-\frac{105}{2D}\Bigl(\delta(\tfrac13)(B-C)-\delta(\tfrac15)(A-C)
  +\delta(\tfrac17)(A-B)\Bigr),
\]
and $\theta'=\tfrac p3+qA-\delta(\tfrac13)$.  Here $D\ge\tfrac{17}{10}A$
(proof of Lemma~\ref{lem:threenode}(c)), $0\le B-C\le B$,
$0\le A-B,A-C\le A$ and $B\le(\tfrac35)^4A$, so, as $z\ge7$,
\[
  |p|\le\tfrac{525}{17}\,3^{-z}\bigl(3(\tfrac35)^4+\tfrac52(\tfrac35)^7
  +\tfrac73(\tfrac37)^7\bigr)\le15\cdot3^{-z};
\]
and $|q|\le\max\bigl(3\delta(\tfrac13)+7\delta(\tfrac17),10\delta(\tfrac15)\bigr)/D
\le\tfrac{10}{17}\,3^{x-z}\bigl(9+\tfrac{49}3(\tfrac37)^7\bigr)\le6\cdot3^{x-z}$,
as $25\cdot5^{-z}\le9\cdot3^{-z}$.  Hence $|\theta'|\le14\cdot3^{-z}$.  Let
$\rho$ agree with $\sigma$ except that $\rho_1=s_1+p$, $\rho_x=s_x+q$,
$\rho_z=1-2\varphi$ and $\rho_d=-1$ for $d>z$.  As $\sigma_d=1$ for $d>x$,
$\sum_{d\ge1}\rho_da^d=t-\delta(a)+pa+qa^x=t+\theta'$ at
$a=\tfrac13,\tfrac15,\tfrac17$.  Put
$F(z,\varphi)=\sum_{d\ge1}\rho_dy^d-t-\theta'$.  The pairs $(z,0)$ and
$(z+1,1)$ give the same $\delta$, hence the same $p,q,\theta'$ and $\rho$,
so $F$ is continuous along the path through the pairs ordered by increasing
$z$ and, for each $z$, decreasing $\varphi$.  At $(x+3,1)$,
\[
  \sum_{d\ge1}\rho_dy^d=(s_1+p)y+(s_x+q)y^x
  -\frac{y^2+2y^{x+3}-y^x-y^{x+1}}{1-y},
\]
where $y^2+2y^{x+3}-y^x-y^{x+1}=y^2-(1-y)y^x(2y^2+2y+1)\ge y^2-5(1-y)$ and
$y^2/(1-y)=1/(r\varepsilon)$.  With $0<s_1\le\tfrac12$ and $s_x>0$
(Lemma~\ref{lem:threenode}(c)), $|p|\le15\cdot3^{-7}$, $|q|\le6\cdot3^{-3}$
and $t+\theta'>0$, this gives $F(x+3,1)<s_x+7-1/(r\varepsilon)\le0$.  As
$z\to\infty$, $p,q,\theta'\to0$ and
$F\to\sum_d\sigma_dy^d-t\ge(y^x+y^{x+1}-y^2)/(1-y)-t\ge\tfrac1{24}-t>0$, as in
the proof of Theorem~\ref{thm:familyone}.  So $F=0$ at some $(z,\varphi)$
with $z\ge x+3$.  There $\theta=t+\theta'$ satisfies
$\sum_{d\ge1}\rho_da^d=\theta$ at all four nodes,
$|\rho_1|\le\tfrac12+15\cdot3^{-7}<1$, $|\rho_z|\le1$, and $|\rho_d|=1$ for
every other $d\ne x$.  For a certificate $v'$ for $\{x\}$, summing over the
nodes as in the proof of Lemma~\ref{lem:exemptpair} gives
$\theta=\sum_{d\ge1}\rho_dv'_d\le\Tail(v')$, so $\tau(\{x\})\ge\theta$.  The
full certificate $v$ with zero set $\{1,x,z\}$ is positive at $0$ and
changes sign exactly at $1$, $x$ and $z$, so $\rho_d=\sgn(v_d)$ off its zero
set, and the same sum gives $\theta=\Tail(v)$.  Hence
$\tau(\{x\})=\Tail(v)=\theta\ge t-14\cdot3^{-z}$, while $\tau(\{x\})\le t$
as $U^{(x)}$ is a certificate for $\{x\}$.

(b) By \eqref{eq:interr} at the nodes $\tfrac15,\tfrac17$ and $t=y$, the
sum $g_d=y^d-\tfrac{35y-5}2\,5^{-d}+\tfrac{35y-7}2\,7^{-d}$ equals
$(y-\tfrac15)(y-\tfrac17)h_{d-2}(\tfrac15,\tfrac17,y)$: it vanishes at
$d=0,1$, is positive for $d\ge2$, and
$\sum_dg_d=(y-\tfrac15)(y-\tfrac17)/\bigl((1-\tfrac15)(1-\tfrac17)(1-y)\bigr)$.
For $x>c$, $U^{(c)}-(U^{(c)}_x/g_x)g$ is a certificate for $\{x\}$ on the
four nodes, so
\[
  \tau(\{x\})\le t(c)+\frac{|U^{(c)}_x|}{g_x}\sum_dg_d
  =t(c)+\frac{35}{24}\cdot\frac{|U^{(c)}_x|}{(1-y)h_{x-2}(\tfrac15,\tfrac17,y)} .
\]
Here $h_k(\tfrac15,\tfrac17,y)=y^k\sum_{i+j\le k}(5y)^{-i}(7y)^{-j}\ge
y^kh_k(\tfrac15,\tfrac17,1)$ and $(1-y)y^{x-2}=\varepsilon r^{1-x}$;
and $0<U^{(c)}_x\le\ell(c)3^{-x}+\lambda_7(c)7^{-x}$, as $U^{(c)}$ is
positive past $c$ and its weight on $\tfrac15$ is negative by
Lemma~\ref{lem:threenode}(a).  The bound decreases in $x$, as $r<3$ and
$h_k(\tfrac15,\tfrac17,1)$ increases with $k$.

(c) Dividing the closed forms of Lemma~\ref{lem:threenode} by $A$,
\[
  18\cdot3^{-x}s_x=\frac{1-36A+90B-56C}{1-\tfrac{10}3(\tfrac35)^x+\tfrac73(\tfrac37)^x},
  \qquad
  R(x)=\frac{1-\tfrac75(\tfrac57)^x-9A+\tfrac{84}5A(\tfrac57)^x-7C}
  {1-\tfrac{10}3(\tfrac35)^x+\tfrac73(\tfrac37)^x},
\]
and both tend to $1$; $R(x)\ge\tfrac9{20}$ is
$\tfrac1{24}-t(x)\ge\tfrac1{16}(\tfrac35)^x$.  At $x=c$,
$K_c=18(5B-7C)/(5B-7C-45AB+84AC-35BC)$.  The denominator is
$12D(\tfrac1{24}-t(c))>0$ and falls short of $5B-7C$ by
$A(45B-84C)+35BC>0$, as $(\tfrac75)^c\ge\tfrac{49}{25}>\tfrac{84}{45}$; so
$K_c>18$.  And $K_c\to18$, since
$(45AB+84AC+35BC)/(5B-7C)\le(45A+84A(\tfrac57)^c+35C)/(5-7(\tfrac57)^c)\to0$.
\end{proof}

As $h_k(\tfrac15,\tfrac17,1)<\tfrac{35}{24}$, the bound in (b) exceeds
$t(c)+\ell(c)/(3^x\varepsilon)$, which exceeds $\tfrac1{24}$ while
$3^x\varepsilon<18$, as $K_c>18$: there only (a) is informative.  Once
$3^x\varepsilon\ge\kappa>18$ the bound tends, as $r\to1^+$, to at most
$t(c)+\ell(c)/\kappa$, below $\tfrac1{24}$ for $c$ with $K_c<\kappa$.

\begin{theorem}[The constant oscillates]\label{thm:familyosc}
Let $P_r=(x-r)(x-3)(x-5)(x-7)$, $u=1$, $\varepsilon=r-1$ and
$\alpha=\log(5/3)/\log3$, and put
\[
  \Delta(r)=\bigl(\inf_Qb_2(QP_r)-24\bigr)\varepsilon^{-\alpha},\qquad
  s=\log_3\frac{18}\varepsilon,\qquad x^*=\lceil s\rceil-1,\qquad
  \omega=s-x^*,
\]
so that $x^*$ is the largest integer with $3^{x^*}\varepsilon<18$ and
$0<\omega\le1$, and $\Phi(\omega)=80\cdot18^{-\alpha}(\tfrac53)^\omega$, so
that $\Phi(\omega)=80(\tfrac35)^{x^*}\varepsilon^{-\alpha}$.
\begin{enumerate}
\item[(a)] Let $0<\delta<\tfrac12$.  For $r>1$ close enough to $1$ with
$\delta\le\omega\le1-\delta$, the placement $\{x^*\}$ is the only worst
placement, $T=\tau(\{x^*\})$ is the tail of the full certificate with zero
set $\{1,x^*,z\}$ for some $z\ge x^*+3$, and
\[
  \frac1{t(x^*)}\le\inf_Qb_2(QP_r)\le\frac1{t(x^*)-14\cdot3^{-x^*-3}} .
\]
Moreover $\Delta(r)-\Phi(\omega)\to0$ as $r\to1^+$ with
$\delta\le\omega\le1-\delta$.
\item[(b)] $\liminf_{r\to1^+}\Delta(r)=\Phi(0)=80\cdot18^{-\alpha}=20.865\ldots$
and $\limsup_{r\to1^+}\Delta(r)=\Phi(1)=\tfrac{400}3\cdot18^{-\alpha}=34.775\ldots$,
and the limit points of $\Delta(r)$ as $r\to1^+$ fill $[\Phi(0),\Phi(1)]$.
\end{enumerate}
\end{theorem}

So $\Delta$ has no limit.  Along $r_n=1+18\cdot3^{-n-\omega_0}$ it tends to
$\Phi(\omega_0)$, a function of $\log_3\frac1{r-1}$ modulo $1$: the excess
$\inf_Qb_2-24$ is $(80+o(1))(\tfrac35)^{x^*}$ off the jumps, and it drops
by the factor $\tfrac35$ each time $3^x(r-1)$ crosses $18$ as $r$ decreases.

\begin{proof}
As $r\to1^+$, $\varepsilon\to0$, $s\to\infty$ and $s\varepsilon\to0$.  By
Theorem~\ref{thm:familyone} and Corollary~\ref{cor:familyone}, $T<\tfrac1{24}$,
$T\to\tfrac1{24}$ and $(\tfrac1{24}-T)\varepsilon^{-\alpha}$ is bounded, so
$\inf_Qb_2-24=24(\tfrac1{24}-T)/T$ gives
$\Delta(r)-576(\tfrac1{24}-T)\varepsilon^{-\alpha}\to0$.  As
$(\tfrac35)^x=3^{-\alpha x}$, $3^\alpha=\tfrac53$ and
$\varepsilon=18\cdot3^{-s}$,
\begin{equation}\label{eq:oscscale}
  80\,(\tfrac35)^x\varepsilon^{-\alpha}=\Phi(0)\,(\tfrac53)^{s-x}
  \qquad\text{for every }x .
\end{equation}
Recall from the proof of Theorem~\ref{thm:familyone} that
$\tau(\{1\}),\tau(\{2\})\le\tfrac1{48}$ and $\tau(\{x\})\le t(x)$ for
$x\ge3$, where $t(3)=\tau^*<\tfrac1{48}$ and
$\tfrac1{24}-t(x)=\tfrac5{36}(\tfrac35)^xR(x)$ for $x\ge4$.  Fix
$\delta\in(0,1)$.

\emph{(i) Far placements.}  There is $\eta>0$ with
$\tau(\{x\})\le\tfrac1{24}-\eta$ whenever $3^x\varepsilon\ge18\cdot3^\delta$,
for $r$ close to $1$.  By Lemma~\ref{lem:eighteen}(c) take $c\ge2$ with
$K_c<18\cdot3^{\delta/2}$.  Let $x_0$ be the least integer with
$3^{x_0}\varepsilon\ge18\cdot3^\delta$; as $r\to1^+$, $x_0\to\infty$ and
$x_0\varepsilon\to0$.  For $x\ge x_0>c$, Lemma~\ref{lem:eighteen}(b) and its
monotonicity bound $\tau(\{x\})$ by
\[
  t(c)+\frac{35}{24}\cdot
  \frac{\ell(c)+\lambda_7(c)(\tfrac37)^{x_0}}{h_{x_0-2}(\tfrac15,\tfrac17,1)}
  \cdot\frac{r^{x_0-1}}{18\cdot3^\delta},
\]
which tends to $t(c)+\ell(c)/(18\cdot3^\delta)
=\tfrac1{24}-(\tfrac1{24}-t(c))(1-K_c/(18\cdot3^\delta))$, since
$h_k(\tfrac15,\tfrac17,1)\to(1-\tfrac15)^{-1}(1-\tfrac17)^{-1}=\tfrac{35}{24}$
and $r^{x_0}\to1$.  As $K_c/(18\cdot3^\delta)<3^{-\delta/2}$, any
$\eta<(\tfrac1{24}-t(c))(1-3^{-\delta/2})$ will do.

\emph{(ii) Near placements.}  For $r$ close to $1$, every $x\ge4$ with
$3^x\varepsilon\le18\cdot3^{-\delta}$ satisfies the hypotheses of
Lemma~\ref{lem:eighteen}(a), so
$t(x)-14\cdot3^{-x-3}\le\tau(\{x\})\le t(x)$.  Such $x$ have $x\le s$, and
$y^x(1+y)-y^2\ge2y^{x+1}-1\ge2r^{-s-1}-1\to1$.  By
Lemma~\ref{lem:eighteen}(c) there is $x_1$ with
$18\cdot3^{-x}s_x\le3^{\delta/2}$ for $x\ge x_1$, and then
$s_x\le3^{-\delta/2}/\varepsilon$; for $4\le x<x_1$,
$s_x\le3^{x_1}/8$ by Lemma~\ref{lem:threenode}(c).  Either way
$s_x+7\le1/(r\varepsilon)$ once $\varepsilon$ is small.

(a) Let $\delta\le\omega\le1-\delta$.  Then
$3^{x^*}\varepsilon=18\cdot3^{-\omega}\le18\cdot3^{-\delta}$ and
$x^*\to\infty$, so by (ii) $\tau(\{x^*\})$ is the tail of a full certificate
with zero set $\{1,x^*,z\}$, $z\ge x^*+3$, and
\begin{equation}\label{eq:oscworst}
  \tfrac5{36}(\tfrac35)^{x^*}R(x^*)\le\tfrac1{24}-\tau(\{x^*\})
  \le\tfrac5{36}(\tfrac35)^{x^*}R(x^*)+\tfrac{14}{27}3^{-x^*},
\end{equation}
where $R(x^*)\to1$ and $3^{-x^*}=5^{-x^*}(\tfrac35)^{x^*}$.  Every other
placement costs less for $r$ close to $1$.  A placement $x\ge x^*+1$ has
$3^x\varepsilon\ge54\cdot3^{-\omega}\ge18\cdot3^\delta$, so costs at most
$\tfrac1{24}-\eta$ by (i).  Placements $x\le3$ cost at most
$\tfrac1{48}$.  For $4\le x\le x^*-2$,
$\tfrac1{24}-\tau(\{x\})\ge\tfrac5{36}(\tfrac35)^x\cdot\tfrac9{20}
\ge\tfrac54\cdot\tfrac5{36}(\tfrac35)^{x^*}$, and for $x=x^*-1$,
$\tfrac1{24}-\tau(\{x\})\ge\tfrac53\cdot\tfrac5{36}(\tfrac35)^{x^*}R(x^*-1)$.
All of these exceed the right side of \eqref{eq:oscworst} for $r$ close to
$1$.  So $T=\tau(\{x^*\})$, Theorem~\ref{thm:value} and (ii) give the bounds
on the infimum, and by \eqref{eq:oscworst} and \eqref{eq:oscscale},
$576(\tfrac1{24}-T)\varepsilon^{-\alpha}=\Phi(\omega)\bigl(R(x^*)+O(5^{-x^*})\bigr)$,
which differs from $\Phi(\omega)\le\Phi(1)$ by $o(1)$ uniformly in $\omega$.

(b) \emph{Lower bound.}  For $r$ close to $1$ every placement $x$ has
$\tfrac1{24}-\tau(\{x\})$ at least $\eta$ if $3^x\varepsilon\ge18\cdot3^\delta$
(by (i)), at least $\tfrac1{48}$ if $x\le3$, at least
$\tfrac54\cdot\tfrac5{36}(\tfrac35)^s$ if $4\le x\le s-2$, and at least
$\tfrac5{36}(\tfrac35)^{s+\delta}R(x)$ if $s-2<x<s+\delta$, where
$R(x)\to1$ uniformly as $x>s-2\to\infty$.  So
$\tfrac1{24}-T=\inf_x(\tfrac1{24}-\tau(\{x\}))$ and \eqref{eq:oscscale} give
$\liminf_{r\to1^+}\Delta(r)\ge\Phi(0)(\tfrac35)^\delta$.

\emph{Upper bound.}  Let $\hat x$ be the largest integer with
$3^{\hat x}\varepsilon\le18\cdot3^{-\delta}$, so $\hat x>s-1-\delta$ and
$\hat x\to\infty$.  By (ii),
$\tfrac1{24}-T\le\tfrac1{24}-\tau(\{\hat x\})\le\tfrac5{36}(\tfrac35)^{\hat x}R(\hat x)
+\tfrac{14}{27}3^{-\hat x}$, where
$80(\tfrac35)^{\hat x}\varepsilon^{-\alpha}<\Phi(0)(\tfrac53)^{1+\delta}
=\Phi(1)(\tfrac53)^\delta$ by \eqref{eq:oscscale} and
$3^{-\hat x}\varepsilon^{-\alpha}<3^{1+\delta}\varepsilon^{1-\alpha}/18\to0$.
So $\limsup_{r\to1^+}\Delta(r)\le\Phi(1)(\tfrac53)^\delta$.

As $\delta$ is arbitrary, $\Phi(0)\le\liminf\Delta\le\limsup\Delta\le\Phi(1)$.
For $0<\omega_0<1$ the points $r_n=1+18\cdot3^{-n-\omega_0}$ have $s=n+\omega_0$,
$x^*=n$ and $\omega=\omega_0$, so $\Delta(r_n)\to\Phi(\omega_0)$ by (a).
Hence every point of $[\Phi(0),\Phi(1)]$ is a limit point, and the
$\liminf$ and $\limsup$ are the endpoints.
\end{proof}

At $r=\tfrac{5701}{5700}$, far below $\rho_1$, everything is exact:
$3^{10}\varepsilon=10.35\ldots$ and $3^{11}\varepsilon=31.07\ldots$, so
$x^*=10$ and $\omega=0.50\ldots$.  The hypotheses of
Lemma~\ref{lem:eighteen}(a) hold at $x=10$, and the witness of the full
certificate with zero set $\{1,10,1335\}$ has $s_1=0.4525\ldots$ and
$s_{1335}=-0.3405\ldots$, so by \eqref{eq:witness} $\tau(\{10\})$ is its
tail, which lies within $10^{-600}$ below $t(10)$; the placements $4\le x\le9$
cost at most $t(x)<\tau(\{10\})$, those $x\le3$ at most $\tfrac1{48}$, and
Lemma~\ref{lem:eighteen}(b) with $c=4$ at $x=11$ bounds every $x\ge11$ by
$0.03593<\tau(\{10\})=0.04085\ldots$.  So $T=\tau(\{10\})$ and
$\inf_Qb_2(QP_r)=24.47884\ldots$, with $\Delta(r)=26.70\ldots$ against
$\Phi(\omega)=26.97\ldots$; this is checked exactly in
\texttt{test\_family\_oscillation} (\texttt{tests/test\_attainment.py}).
There $t(11)>t(10)$, so the placement $11$ is beaten only through the node
$\tfrac1r$, while Lemma~\ref{lem:eighteen}(b) at $x=10$ beats $t(10)$ for no
$c$, as $3^{10}\varepsilon<18$; the test checks both as controls, and that
the hypothesis $s_x+7\le1/(r\varepsilon)$ of Lemma~\ref{lem:eighteen}(a)
fails at $x=11$.
Theorem~\ref{thm:familyosc} leaves out neighbourhoods of the jumps, where
$3^{x^*}\varepsilon$ is just below $18$ or $3^{x^*+1}\varepsilon$ just above
it.  To cross them we first compute every cost along the family.

\paragraph{The zero set of one placement.}  The witness of
Lemma~\ref{lem:eighteen}(a) flips its entries far out.  Flipping entries
just below $x$ as well gives the cost of every placement along the family
in closed form, with its minimising zero set, at every $1<r<3$.  Fix
$1<r<3$, $y=\tfrac1r$ and an integer $x\ge2$, and let $U=U^{(x)}$,
$t=t(x)$, $s_1$, $s_x$ and $\sigma$ be as in Lemma~\ref{lem:threenode},
whose parts (a) and (b) hold for $x\ge2$; also $0<s_1\le\tfrac12$, by
part (c) for $x\ge4$, while $s_1=\tfrac7{24}$ and $\tfrac{859}{3408}$ at
$x=2,3$.  Let $V$ and $V'$ be the exponential sums on
$(\tfrac13,\tfrac15,\tfrac17)$ with $(V_0,V_1,V_x)=(0,1,0)$ and
$(V'_0,V'_1,V'_x)=(0,0,1)$.  With $A,B,C,D$ as in
Lemma~\ref{lem:threenode}, their weights are $-\tfrac{105}{2D}(B-C,C-A,A-B)$
and $\tfrac1D(3,-10,7)$, so
$V'_d=h_{d-2}(\tfrac13,\tfrac15,\tfrac17)/h_{x-2}(\tfrac13,\tfrac15,\tfrac17)$
by the proof of Lemma~\ref{lem:threenode}(a).  Put
\[
  H_d=y^d-U_d-yV_d-y^xV'_d,\qquad
  F=\sum_{d\ge1}\sigma_dy^d-t=s_1y+s_xy^x+\frac{y^x+y^{x+1}-y^2}{1-y}-t .
\]
Call the integers $d\ge2$, $d\ne x$, the \emph{positions}; unmarked sums
over $d$ run over them.  Order them as $\ldots,x+2,x+1,x-1,\ldots,2$, and for
a position $z$ let $\Pi(z)$ be the set of positions before $z$ in this
order: all $d>z$ if $z>x$, and all $d>x$ with all $z<d<x$ if $z<x$.

\begin{theorem}[One placement along the family]\label{thm:familyzero}
In this notation:
\begin{enumerate}
\item[(a)] $H$ is the exponential sum on $(y,\tfrac13,\tfrac15,\tfrac17)$
with weight $1$ on $y$ vanishing at $d=0,1,x$.  At every position $U_d$ and
$H_d$ have the sign of $d-x$ and $V_d$ the opposite sign, and
\[
  \sum_d|V_d|=s_1,\qquad\sum_d|H_d|=F>0 .
\]
\item[(b)] If $0\le\varphi_d\le1$ at the positions and
$2\sum_d\varphi_d|H_d|\ge F$, then $\tau(\{x\})\ge t-2\sum_d\varphi_d|U_d|$.
\item[(c)] Let $z$ be the position with
$\sum_{d\in\Pi(z)}|H_d|<\tfrac F2\le\sum_{d\in\Pi(z)}|H_d|+|H_z|$, and
$\varphi=(\tfrac F2-\sum_{d\in\Pi(z)}|H_d|)/|H_z|\in(0,1]$.  Then
\[
  \tau(\{x\})=t-2\sum_{d\in\Pi(z)}|U_d|-2\varphi|U_z|,
\]
the tail of the full certificate on $(y,\tfrac13,\tfrac15,\tfrac17)$ with
zero set $\{1,x,z\}$ if $z>x$ and $\{1,z,x\}$ if $z<x$.  If $\varphi<1$, it is
the only certificate achieving $\tau(\{x\})$.
\end{enumerate}
\end{theorem}

So the third zero is a half-mass point of $|H|$, flipped from the far end,
as $N$ is in Theorem~\ref{thm:halfmass}(c); it lies above $x$ exactly when
$\sum_{d>x}H_d\ge\tfrac F2$.  Parts (b) and (c) are the witness of
Lemma~\ref{lem:eighteen}(a) with the correction at the positions $1$ and
$x$ written through the kernel $H$.

\begin{proof}
Write $\mathbf a=(\tfrac13,\tfrac15,\tfrac17)$ and $a$ for one of its
nodes.  The $3\times3$ matrix $(a^d)$, $d\in\{0,1,x\}$, is nonsingular by the
zero bound, so for every integer $d$ there are unique $c_0,c_1,c_x$ with
$a^d=c_0+c_1a+c_xa^x$ at the three nodes.  Summing against the weights of
$U$, which satisfy $\sum_a\lambda_a=1$ and $\sum_a\lambda_aa=\sum_a\lambda_aa^x=0$,
gives $c_0=U_d$, and likewise $c_1=V_d$, $c_x=V'_d$:
\begin{equation}\label{eq:threebasis}
  a^d=U_d+V_da+V'_da^x\qquad\text{at }a=\tfrac13,\tfrac15,\tfrac17 .
\end{equation}

(a) So $H_d=y^d-\sum_a\mu_aa^d$ with $\mu_a$ the weight of $U+yV+y^xV'$ at
$a$, and $H_0=1-1=0$, $H_1=y-y=0$, $H_x=y^x-y^x=0$; also $\sum_a\mu_a=1$.
By the zero bound, $U$ is negative on $(1,x)$ and positive past $x$
(proof of Lemma~\ref{lem:threenode}(a)).  The sum $V$ is nonzero on three
nodes and vanishes at $0$ and $x$, so it has no other zeros and changes
sign at both; as $V_1=1$, it is positive on $(0,x)$ and negative past $x$.
The sum $H$ is nonzero on four nodes and vanishes at $0$, $1$ and $x$, so it
has no other zeros and changes sign at each; its weight on the largest node
$y>\tfrac13$ is $1$, so it is positive past $x$ and negative on $(1,x)$.  By
Lemma~\ref{lem:threenode}(b), $\sum_{d\ge1}\sigma_da^d=t$ at the three
nodes; summing against the weights of $V$, which add up to $V_0=0$, gives
$\sum_{d\ge1}\sigma_dV_d=0$, and against those of $H$,
$\sum_{d\ge1}\sigma_dH_d=\sum_{d\ge1}\sigma_dy^d-t\sum_a\mu_a=F$.  At
every position $\sigma_d$ is $-1$ below $x$ and $1$ above it, so
$\sigma_dV_d=-|V_d|$ and $\sigma_dH_d=|H_d|$; with $V_1=1$, $V_x=0$ and
$H_1=H_x=0$ the two sums read $s_1-\sum_d|V_d|=0$ and $\sum_d|H_d|=F$.

To see $F>0$, put $p(a)=(1-a)(\sum_{d\ge1}\sigma_da^d-t)$, the polynomial
\[
  p(a)=-t+(s_1+t)a-(s_1+1)a^2+(s_x+1)a^x+(1-s_x)a^{x+1}
\]
for $x\ge3$, the terms in $a^2$ combining for $x=2$.  It has at most
five terms, so by Descartes' rule at most four positive zeros counted with
multiplicity.  It vanishes at $\tfrac17,\tfrac15,\tfrac13$, while
$p(0)=-t<0<1=p(1)$, so the number of its zeros in $(0,1)$ is odd, hence
three: $\tfrac17,\tfrac15,\tfrac13$, all simple.  So $p>0$ on
$(\tfrac13,1]$, and $F=p(y)/(1-y)>0$.

(b) Put $\eta_d=-2\varphi_d\sgn(U_d)$ at the positions.  For $0\le s\le1$
let $\rho(s)$ agree with $\sigma+s\eta$ at the positions and have
$\rho_1=s_1-s\sum_d\eta_dV_d$ and $\rho_x=s_x-s\sum_d\eta_dV'_d$; the series
converge absolutely, as $V$ and $V'$ decay geometrically.  By
\eqref{eq:threebasis}, at each node $a$,
\[
  \sum_{d\ge1}\rho_da^d=t+s\sum_d\eta_d(a^d-V_da-V'_da^x)
  =t+s\sum_d\eta_dU_d=t-2s\sum_d\varphi_d|U_d|=:\theta(s),
\]
and at $y$, as $\sgn H_d=\sgn U_d$,
\[
  \sum_{d\ge1}\rho_dy^d-\theta(s)=F+s\sum_d\eta_d(y^d-U_d-V_dy-V'_dy^x)
  =F-2s\sum_d\varphi_d|H_d| .
\]
This is affine in $s$, positive at $s=0$ and at most $0$ at $s=1$, so it
vanishes at some $s\in(0,1]$.  There $\rho=\rho(s)$ satisfies
$\sum_{d\ge1}\rho_da^d=\theta$ at all four nodes, with
$\theta=\theta(s)\ge t-2\sum_d\varphi_d|U_d|$.  At a position
$\rho_d=\sigma_d(1-2s\varphi_d)\in[-1,1]$, and $\sgn(U_d)V_d=-|V_d|$ gives
$\rho_1=s_1-2s\sum_d\varphi_d|V_d|\in[-s_1,s_1]$ by (a).  So $|\rho_d|\le1$
for every $d\ne x$.  For a certificate $v$ for $\{x\}$ on the four nodes,
summing over them as in the proof of Lemma~\ref{lem:exemptpair} gives
$\theta=\sum_{d\ge1}\rho_dv_d=\sum_{d\ne x}\rho_dv_d\le\Tail(v)$, so
$\tau(\{x\})\ge\theta$.

(c) As $z$ runs through the positions in their order, the sums
$\sum_{\Pi(z)}|H_d|$ increase from $0$ (the $H_d$ are summable) by the
steps $|H_z|>0$ and end at $F-|H_2|$, so $z$ exists, is unique, and
$0<\varphi\le1$.  Take $\varphi_d=1$ on $\Pi(z)$, $\varphi_z=\varphi$ and
$\varphi_d=0$ otherwise.  Then $2\sum_d\varphi_d|H_d|=F$, so the proof of
(b) runs with $s=1$, and $\rho=\rho(1)$ has $\rho_d=-\sigma_d$ on $\Pi(z)$,
$\rho_z=\sigma_z(1-2\varphi)$ and $\rho_d=\sigma_d$ at the other positions:
if $z>x$, $\rho_d$ is $-1$ on $(1,x)$, $1$ on $(x,z)$ and $-1$ past $z$; if
$z<x$, it is $-1$ on $(1,z)$, $1$ on $(z,x)$ and $-1$ past $x$.  The full
certificate $v$ with zero set $\{1,x,z\}$ has $v_0=1$ and changes sign
exactly at its three zeros, so $\rho_d=\sgn(v_d)$ off its zero set, and
$\theta=\sum_{d\ge1}\rho_dv_d=\Tail(v)$.  Hence
$\Tail(v)=\theta\le\tau(\{x\})\le\Tail(v)$, and $\theta$ is the displayed
value.  If $\varphi<1$ and $v'$ achieves $\tau(\{x\})$, then
$\theta=\sum_{d\ne x}\rho_dv'_d\le\sum_{d\ne x}|v'_d|=\theta$ forces
$\rho_dv'_d=|v'_d|$ for $d\ne x$, so $v'_d=0$ where $|\rho_d|<1$: at $d=1$,
as $|\rho_1|\le s_1<1$, and at $d=z$, as $|1-2\varphi|<1$.  So $v'$ vanishes
on $\{1,x,z\}$ and is $v$.
\end{proof}

With Theorem~\ref{thm:halfmass}, this makes $T$ along the family a finite
computation in closed form at every $1<r<3$: $Z^*=\{1,3\}$, so
$T=\max(\tau^*,\tau(\{x\}):x=2\text{ or }4\le x<N)$, each $\tau(\{x\})$
given by (c) through finitely many powers of $r$ and geometric sums.  The
costs found in Theorems~\ref{thm:elevenvalue}, \ref{thm:family}
and~\ref{thm:familylow} are instances: at $r=\tfrac54$ the placement $4$ has
$z=5$, zero set $\{1,4,5\}$ and cost $\tfrac{7627}{314024}$, and at
$r=\tfrac{11}{10}$ the placement $5$ has $z=3$, zero set $\{1,3,5\}$ and
cost $\tfrac{96935}{3398808}$.  For a fixed placement $x$ the zero set can
change with $r$ only where $\tfrac F2=\sum_{\Pi(z)}|H_d|+|H_z|$ for a
position $z$, where $\{1,x,z\}$ (or $\{1,z,x\}$) and the certificate with
the next position in the order tie; each side of this equation is a
combination of powers of $r$ and of geometric sums in $r$, so these are
the boundaries of the pieces in closed form.  In the pattern after
Theorem~\ref{thm:familylow}, the entry $s_z$ of the witness is
$1-2\varphi$ for $z>x$, and passing from $-1$ to $1$ is $\varphi$ passing
from $1$ to $0$.

\paragraph{The limit profile.}  Theorem~\ref{thm:familyzero} settles what
the costs tend to on the scale $3^x(r-1)\to\kappa$.  Below, $g_d$ is the
sum $3^{-d}-\tfrac{10}3\,5^{-d}+\tfrac73\,7^{-d}=\tfrac8{315}h_{d-2}(\tfrac13,\tfrac15,\tfrac17)$
of the proof of Lemma~\ref{lem:threenode}(a), positive for $d\ge2$;
$U^\infty_d=\tfrac72\,7^{-d}-\tfrac52\,5^{-d}$ is the full certificate on
$(\tfrac15,\tfrac17)$ with zero set $\{1\}$, negative for $d\ge2$, with
tail $\tfrac1{24}$ (proof of Proposition~\ref{prop:cutoffsharp}); and
$M_c=2\sum_{d\ge c}g_d=3\cdot3^{-c}-\tfrac{25}3\,5^{-c}+\tfrac{49}9\,7^{-c}$.

\begin{theorem}[The limit profile]\label{thm:familylimit}
Let $\tau(\{x\})$ be the costs at $P_r=(x-r)(x-3)(x-5)(x-7)$, $u=1$, let
$t(c)$, $\ell(c)$ and $K_c$ be as in Lemmas~\ref{lem:threenode}
and~\ref{lem:eighteen}, and put $k_2=\infty$ and
\[
  k_c=\frac{18}{1-18\cdot3^{-c}+30\cdot5^{-c}-14\cdot7^{-c}}\qquad(c\ge3).
\]
\begin{enumerate}
\item[(a)] $k_3=\tfrac{66150}{1957}=33.80\ldots$, the $k_c$ decrease
strictly, and $k_c\to18$.  Let $G(\kappa)=\tfrac1{24}$ for
$0\le\kappa\le18$ and $G(\kappa)=t(c)+\ell(c)/\kappa$ for
$k_c\le\kappa\le k_{c-1}$, $c\ge3$.  Then $G$ is well defined, continuous and
nonincreasing, $G(\kappa)<\tfrac1{24}$ for $\kappa>18$, and
\[
  G(\kappa)=\min\Bigl(\frac1{24},\ \inf_{c\ge2}\Bigl(t(c)+\frac{\ell(c)}\kappa\Bigr)\Bigr)
  \qquad(\kappa>0).
\]
\item[(b)] If $x\to\infty$ and $r\to1^+$ with $3^x(r-1)\to\kappa\in[0,\infty)$,
then $\tau(\{x\})\to G(\kappa)$.
\item[(c)] In (b), if $\kappa<18$, the zero set of
Theorem~\ref{thm:familyzero}(c) is eventually $\{1,x,z\}$ with $z>x$ and
$(r-1)z\to\log\frac{36}{18+\kappa}$; if $k_c<\kappa<k_{c-1}$, it is eventually
$\{1,c,x\}$.
\item[(d)] There are constants $C_2>C_1>0$ with
$C_1\eta^{1+\alpha}\le\tfrac1{24}-G(18(1+\eta))\le C_2\eta^{1+\alpha}$ for
$0<\eta\le1$, where $\alpha=\log(5/3)/\log3$.
\item[(e)] As $\eta\to0^+$,
$(\tfrac1{24}-G(18(1+\eta)))(18/\eta)^{1+\alpha}$ has no limit: its limit
points fill $[\tfrac54,f(\theta^*)]=[1.25,1.3814\ldots]$, where
$f(\theta)=\tfrac54(2\theta-1)\theta^{-1-\alpha}$ and
$\theta^*=\tfrac{1+\alpha}{2\alpha}$.
\end{enumerate}
\end{theorem}

So $\tau(\{x\})$ does tend to $\min_c(t(c)+\ell(c)/\kappa)$, the limit of
the bounds of Lemma~\ref{lem:eighteen}(b), once $\kappa>18$; the minimum is
at the $c\ge3$ with $k_c\le\kappa\le k_{c-1}$ (so $c=3$ for
$\kappa\ge33.80\ldots$, $c=4$ for $21.95\ldots\le\kappa\le33.80\ldots$),
never at $c=2$, and $c\to\infty$ as $\kappa\downarrow18$.  At
$r=\tfrac{101}{100}$ the placement $7$ has $3^7\varepsilon=21.87$, just
below $k_4$, and zero set $\{1,4,7\}$ rather than the limiting $\{1,5,7\}$.

\begin{proof}
\emph{The profile.}  Write $A=3^{-c}$, $B=5^{-c}$, $C=7^{-c}$ and $D$ as in
Lemma~\ref{lem:threenode}.  Expanding the closed forms of
Lemma~\ref{lem:threenode}(a),
\begin{equation}\label{eq:profileid}
  |U^\infty_c|=\ell(c)\,g_c,\qquad
  \frac1{k_c}=\frac1{18}-M_{c+1},\qquad
  t(c)+\ell(c)\Bigl(\frac1{18}-M_c\Bigr)=\frac1{24}-2\sum_{d\ge c}|U^\infty_d| :
\end{equation}
the first is $\tfrac12(5B-7C)=\tfrac{3(5B-7C)}{2D}\cdot\tfrac D3$, the second is
$M_{c+1}=A-\tfrac53B+\tfrac79C$, and in the third both sides minus
$\tfrac1{24}$ equal $(98C-75B)/12$, the left one after the factor
$D=3A-10B+7C$ cancels.  As $g_d>0$, $M_c$ decreases strictly to $0$, with
$M_c-M_{c+1}=2g_c$, and $M_3=\tfrac{19}{315}>\tfrac1{18}>M_4$; so the $k_c$,
$c\ge3$, are positive, decrease strictly and tend to $18$.  Define $L(w)=0$
for $w\le0$ and
\[
  L(w)=2\sum_{d>c}|U^\infty_d|+\ell(c)(w-M_{c+1})\qquad
  (M_{c+1}\le w\le M_c,\ c\ge3).
\]
By the first identity the pieces agree at their common ends, where both
give $2\sum_{d\ge c}|U^\infty_d|$, so $L$ is continuous on $(-\infty,M_3]$
(at $0$ as $\sum_{d>c}|U^\infty_d|\to0$), nondecreasing, and positive on
$(0,M_3]$.  For $\kappa>0$ put $w(\kappa)=\tfrac1{18}-\tfrac1\kappa<M_3$.
If $\kappa\le18$ then $L(w(\kappa))=0$.  If $k_c\le\kappa\le k_{c-1}$ then
$M_{c+1}\le w(\kappa)\le M_c$ by the second identity, and
$\tfrac1{24}-L(w(\kappa))$ is affine in $1/\kappa$ with slope $\ell(c)$ and, by
the third, equal to $t(c)+\ell(c)/\kappa$ at $1/\kappa=\tfrac1{18}-M_c$;
so it is $t(c)+\ell(c)/\kappa$.  Hence $G(\kappa)=\tfrac1{24}-L(w(\kappa))$
for every $\kappa>0$, which is well defined, continuous, nonincreasing, and
below $\tfrac1{24}$ for $\kappa>18$; and $k_3=\tfrac{66150}{1957}$.

\emph{Asymptotics.}  Let $x\to\infty$ and $r\to1^+$ with
$3^x\varepsilon\to\kappa$, $\varepsilon=r-1$; then $x\varepsilon\to0$ and
$y^x\to1$.  Use the notation of Theorem~\ref{thm:familyzero} at $x$, and
write $A,B,C,D$ at $x$ now.  By Lemma~\ref{lem:eighteen}(c)
$3^{-x}s_x\to\tfrac1{18}$, and $1-y=\varepsilon/r$, so
\begin{equation}\label{eq:Fscale}
  (1-y)F=(1-y)(s_1y+s_xy^x-t)+y^x+y^{x+1}-y^2\to1+\tfrac\kappa{18}.
\end{equation}
Summing the geometric part of $\sum_{d>x}H_d$,
\[
  E:=F-2\sum_{d>x}H_d
  =s_1y+(s_x+1)y^x-\frac{y^2}{1-y}-t+2\sum_{d>x}(U_d+yV_d+y^xV'_d),
\]
where $0<\sum_{d>x}U_d\le t$, and $V_d<0$ for $d>x$ with
$\sum_{d>x}|V_d|\le s_1$, by Theorem~\ref{thm:familyzero}(a), and
$0<\sum_{d>x}V'_d=(\tfrac32A-\tfrac52B+\tfrac76C)/D<1$ for $x\ge3$.  As
$y^2/(1-y)=1/(r\varepsilon)$, this gives $E\le s_x+4-1/(r\varepsilon)$, and,
if $\kappa>0$, $3^{-x}E\to\tfrac1{18}-\tfrac1\kappa$.  Below $x$, $V_d>0$,
$-1<U_d<0$ and $y^d<1$, so for $2\le d<x$
\begin{equation}\label{eq:Hbelow}
  y^xV'_d-2\le|H_d|\le y^xV'_d+1,\qquad
  V'_d=\frac{3g_d}D=\frac{3^xg_d}{1-\tfrac{10}3(\tfrac35)^x+\tfrac73(\tfrac37)^x},
\end{equation}
hence $3^{-x}\sum_{c<d<x}|H_d|\to\sum_{d>c}g_d=\tfrac12M_{c+1}$ and
$3^{-x}|H_c|\to g_c$ for fixed $c\ge2$.  Finally the weights of $U^{(x)}$
tend to $(0,-\tfrac52,\tfrac72)$, so $\sum_{d\ge2}|U^{(x)}_d-U^\infty_d|\to0$.

(b) \emph{Lower bound.}  If $\kappa<18$, then
$r\varepsilon(s_x+4)\to\kappa/18<1$, so eventually $E<0$, that is,
$2\sum_{d>x}H_d>F$; Theorem~\ref{thm:familyzero}(b) with $\varphi_d=1$ for
$d>x$ gives $\tau(\{x\})\ge t(x)-2\sum_{d>x}|U^{(x)}_d|\to\tfrac1{24}$.  Now
let $\kappa>0$, $w=w(\kappa)$, and $\max(w,0)<w'\le M_3$; take $c\ge3$ with
$M_{c+1}<w'\le M_c$ and $\varphi=(w'-M_{c+1})/(2g_c)\in(0,1]$, and put
$\varphi_d=1$ for $d>x$ and for $c<d<x$, and $\varphi_c=\varphi$.  Then
$3^{-x}(2\sum_d\varphi_d|H_d|-2\sum_{d>x}H_d)\to M_{c+1}+2\varphi g_c=w'$,
which exceeds $w=\lim3^{-x}E$, so eventually $2\sum_d\varphi_d|H_d|\ge F$,
and Theorem~\ref{thm:familyzero}(b) gives
\[
  \tau(\{x\})\ge t(x)-2\sum_{d>x}|U_d|-2\sum_{c<d<x}|U_d|-2\varphi|U_c|
  \to\frac1{24}-2\sum_{d>c}|U^\infty_d|-2\varphi|U^\infty_c|=\frac1{24}-L(w'),
\]
by the first identity of \eqref{eq:profileid}.  Letting $w'\downarrow\max(w,0)$,
$\liminf\tau(\{x\})\ge\tfrac1{24}-L(\max(w,0))=G(\kappa)$ in every case.

\emph{Upper bound.}  $\tau(\{x\})\le t(x)\to\tfrac1{24}$, and for $c\ge2$
and $x>c$ the bound of Lemma~\ref{lem:eighteen}(b) is
$t(c)+\tfrac{35}{24}(\ell(c)+\lambda_7(c)(\tfrac37)^x)r^{x-1}/(3^x\varepsilon\,h_{x-2}(\tfrac15,\tfrac17,1))$,
which tends to $t(c)+\ell(c)/\kappa$ if $\kappa>0$, as
$h_{x-2}(\tfrac15,\tfrac17,1)\to\tfrac{35}{24}$.  So $\limsup\tau(\{x\})$ is at
most $\min(\tfrac1{24},\inf_c(t(c)+\ell(c)/\kappa))\le G(\kappa)$, the last as
$G(\kappa)$ is $\tfrac1{24}$ or one of the $t(c)+\ell(c)/\kappa$.  So
$\tau(\{x\})\to G(\kappa)$, and, as every $\kappa>0$ is such a limit (take
$r=1+\kappa3^{-x}$), $G(\kappa)$ is that minimum; this proves (b) and the
rest of (a).

(c) If $\kappa<18$, eventually $E<0$, that is, $\sum_{d>x}H_d>\tfrac F2$, so
$z>x$ and $\sum_{d>z}H_d<\tfrac F2\le\sum_{d\ge z}H_d$.  For $z>x$,
$(1-y)\sum_{d\ge z}H_d=y^z-(1-y)\sum_{d\ge z}(U_d+yV_d+y^xV'_d)=y^z+O(\varepsilon)$
uniformly, by the bounds above, so by \eqref{eq:Fscale}
$y^{z+1}+o(1)<\tfrac12(1+\tfrac\kappa{18})\le y^z+o(1)$; hence
$y^z\to\tfrac{18+\kappa}{36}$, and $z\log r\to\log\tfrac{36}{18+\kappa}$.
If $k_c<\kappa<k_{c-1}$, then $M_{c+1}<w<M_c$, while $3^{-x}E\to w$,
$3^{-x}\cdot2\sum_{c<d<x}|H_d|\to M_{c+1}$ and
$3^{-x}\cdot2\sum_{c\le d<x}|H_d|\to M_c$; as $E=F-2\sum_{d>x}H_d$, eventually
$\sum_{\Pi(c)}|H_d|<\tfrac F2\le\sum_{\Pi(c)}|H_d|+|H_c|$, so $z=c$.

(d) Let $0<\eta\le1$, $\kappa=18(1+\eta)$ and $w=w(\kappa)=\eta/(18(1+\eta))$.
\emph{Upper.}  Let $c\ge3$ with $M_{c+1}<w\le M_c$.  On $[0,w]$ the function
$L$ is affine on pieces with slopes $\ell(c')$, $c'\ge c$, and
$\ell(c')\le5(\tfrac35)^{c'}$ ($\ell(3)=\tfrac{81}{142}$, and
$\ell(c')\le\tfrac{15B}{2D}\le\tfrac{75}{17}(\tfrac35)^{c'}$ for $c'\ge4$,
as $D\ge\tfrac{17}{10}A$ there by the proof of
Lemma~\ref{lem:threenode}(c)).  Also
$w>M_{c+1}\ge3^{-c}(1-\tfrac53(\tfrac35)^c)\ge\tfrac58\,3^{-c}$, so
$(\tfrac35)^c=(3^{-c})^\alpha\le(\tfrac85w)^\alpha$, and
$\tfrac1{24}-G(\kappa)=L(w)\le5(\tfrac85w)^\alpha w\le8(\eta/18)^{1+\alpha}$.
\emph{Lower.}  By (a), $\tfrac1{24}-G(\kappa)\ge\tfrac1{24}-t(c)-\ell(c)/\kappa
=(\tfrac1{24}-t(c))(1-K_c/\kappa)$ for every $c\ge2$.  For $c\ge4$, by the
proof of Lemma~\ref{lem:eighteen}(c), $K_c=18/(1-q)$ with
$q=(45AB-84AC+35BC)/(5B-7C)\le(45A+35C)/(5-7(\tfrac57)^c)\le14.6A\le\tfrac{14.6}{81}$,
so $K_c\le18(1+18\cdot3^{-c})$.  Take the least $c\ge4$ with
$18\cdot3^{-c}\le\eta/2$; then $3^{-c}>\eta/108$, $1-K_c/\kappa\ge
1-(1+\tfrac\eta2)/(1+\eta)\ge\tfrac\eta4$, and
$\tfrac1{24}-t(c)\ge\tfrac1{16}(\tfrac35)^c>\tfrac1{16}(\eta/108)^\alpha$ by
Lemma~\ref{lem:threenode}(c), so $\tfrac1{24}-G(\kappa)\ge\tfrac\eta{64}(\eta/108)^\alpha$.

(e) Let $w\to0^+$, let $c$ be the piece with $M_{c+1}<w\le M_c$, so
$c\to\infty$, and put $\theta=3^cw$; as $3^cM_{c+1}\to1$ and
$3^cM_c\to3$, the limit points of $\theta$ lie in $[1,3]$.  Here
$2\sum_{d>c}|U^\infty_d|=\tfrac54\,5^{-c}-\tfrac76\,7^{-c}$ and
$\ell(c)=\tfrac52(\tfrac35)^c(1+o(1))$, so
$L(w)=5^{-c}(\tfrac54+\tfrac52(\theta-1)+o(1))$, while
$w^{1+\alpha}=\theta^{1+\alpha}5^{-c}$ as $3^{1+\alpha}=5$.  Hence
$L(w)/w^{1+\alpha}-f(\theta)\to0$.  Along $w=\theta_03^{-c}$, $c\to\infty$,
with $\theta_0\in(1,3)$ fixed, $\theta\to\theta_0$; and $f(1)=f(3)=\tfrac54$,
while on $[1,3]$ the function $f$ has its maximum at $\theta^*$, where
$f'$ vanishes.  So the limit points of $L(w)/w^{1+\alpha}$ fill
$[\tfrac54,f(\theta^*)]$.  With $w=\eta/(18(1+\eta))$ the quantity in (e)
is $L(w)w^{-1-\alpha}(1+\eta)^{-1-\alpha}$, with the same limit points.
\end{proof}

\paragraph{Across the jumps.}  Near a jump the placements $x^*$ and $x^*+1$
compete on the scale $(\tfrac35)^{x^*}$, and the profile $G$ decides
between them.

\begin{theorem}[Across the jumps]\label{thm:familyjump}
In the notation of Theorems~\ref{thm:familyosc} and~\ref{thm:familylimit},
extend $\Phi(\omega)=80\cdot18^{-\alpha}(\tfrac53)^\omega$ to all real $\omega$, and put
\[
  \Gamma(r)=576\,\varepsilon^{-\alpha}\Bigl(\frac1{24}-G\bigl(3^{x^*+1}\varepsilon\bigr)\Bigr),
  \qquad\nu(r)=(1-\omega)\,\varepsilon^{-\alpha/(1+\alpha)} .
\]
\begin{enumerate}
\item[(a)] For $r$ close to $1$ every worst placement is $\{x^*\}$ or
$\{x^*+1\}$, and, with no restriction on $\omega$,
\[
  \Delta(r)-\min\bigl(\Phi(\omega),\,\Phi(\omega-1)+\Gamma(r)\bigr)\to0
  \qquad(r\to1^+).
\]
\item[(b)] There are constants $a_2>a_1>0$ such that, for $r$ close to $1$,
$\{x^*\}$ is the only worst placement if $\nu(r)\ge a_2$, and $\{x^*+1\}$ is
the only one if $\nu(r)\le a_1$.
\item[(c)] As $r\to1^+$: $\Delta(r)-\Phi(\omega)\to0$ if $\nu(r)\to\infty$,
in particular uniformly on $0<\omega\le1-\delta$ for each $\delta>0$; and
$\Delta(r)-\Phi(\omega-1)\to0$ if $\nu(r)\to0$.
\end{enumerate}
\end{theorem}

Here $\Phi(\omega-1)=\tfrac35\Phi(\omega)$, and $\Gamma(r)$ is of order
$\nu(r)^{1+\alpha}$ when $1-\omega\le\tfrac12$ (proof of (b)), with a ratio
that oscillates as $1-\omega\to0$ by Theorem~\ref{thm:familylimit}(e).  So as $r$
decreases through a jump, $\Delta$ does not jump from $\Phi(1)$ to $\Phi(0)$:
it passes from $\Phi(\omega)$ to $\Phi(\omega-1)$ while $1-\omega$ runs
through a window of order $(r-1)^{\alpha/(1+\alpha)}$,
$\alpha/(1+\alpha)=\log(5/3)/\log5=0.317\ldots$, which shrinks to the jump,
the worst placement switching from $x^*$ to $x^*+1$ inside it; the costs
there are those of Theorem~\ref{thm:familyzero}(c), with zero sets
$\{1,x,z\}$ or $\{1,z,x\}$.

\begin{proof}
Let $r\to1^+$; then $x^*\to\infty$, $s\varepsilon\to0$ and
$\varepsilon^{-\alpha}(\tfrac35)^{x^*}=\Phi(\omega)/80\le\Phi(1)/80$.  Write
$\beta_x=(\tfrac35)^x$ and $\gamma_x=(\tfrac37)^x$, and let $o(\beta)$ stand
for a quantity whose ratio to $\beta_{x^*}$ tends to $0$.  For
$x\in\{x^*,x^*+1\}$ put $\kappa=3^x\varepsilon\in[6,54]$ and
$w=\tfrac1{18}-\tfrac1\kappa$; so $x\varepsilon=o(\beta)$, $3^{-x}=o(\beta)$,
$\gamma_x=o(\beta)$, and $R(x)\to1$.  We show
\begin{align}
  \tfrac1{24}-\tau(\{x\})&\ge\tfrac5{36}\beta_x+\min\bigl(\tfrac1{24}-G(\kappa),\beta_x\bigr)-o(\beta),
  \label{eq:jumplow}\\
  \tfrac1{24}-\tau(\{x\})&\le\tfrac5{36}\beta_x+\bigl(\tfrac1{24}-G(\kappa)\bigr)+o(\beta)
  \quad\text{if }\tfrac1{24}-G(\kappa)\le\beta_x .\label{eq:jumpup}
\end{align}

\emph{Proof of \eqref{eq:jumplow}.}  As $\tau(\{x\})\le t(x)=\tfrac1{24}-\tfrac5{36}\beta_xR(x)$,
it holds when $\tfrac1{24}-G(\kappa)=o(\beta)$, in particular when
$\kappa\le18$.  Let $\kappa>18$ and $c\ge3$ with $k_c\le\kappa\le k_{c-1}$,
so $G(\kappa)=t(c)+\ell(c)/\kappa$.  If $c\ge x$, then
$0<w\le\tfrac1{18}-\tfrac1{k_{x-1}}=M_x$ and, the slopes of $L$ being at
most $1$ (proof of Theorem~\ref{thm:familylimit}(d)),
$\tfrac1{24}-G(\kappa)=L(w)\le M_x\le4\cdot3^{-x}=o(\beta)$.  Let $c<x$.  The
proof of Lemma~\ref{lem:eighteen}(b) gives
$\tau(\{x\})\le t(c)+J\cdot3^xU^{(c)}_x/\kappa$ with
$J=\tfrac{35}{24}r^{x-1}/h_{x-2}(\tfrac15,\tfrac17,1)$, where
$3^xU^{(c)}_x=\ell(c)+\lambda_5(c)\beta_x+\lambda_7(c)\gamma_x>0$,
$\lambda_5(c)$ being the weight of $U^{(c)}$ on $\tfrac15$.  Here $J\ge1$ and
$J-1=O(x\varepsilon+x5^{-x})=o(\beta)$, as
$\tfrac{35}{24}-h_k(\tfrac15,\tfrac17,1)=\sum_{i+j>k}5^{-i}7^{-j}\le(k+2)5^{-k}$;
$\lambda_5(c)\le-\tfrac52$ by Lemma~\ref{lem:threenode}(a), as
$10\cdot5^{-c}\ge14\cdot7^{-c}$; and
$\ell(c)\le1$, $0<\lambda_7(c)\le7$ for $c\ge3$.  So
$\tau(\{x\})\le G(\kappa)-\tfrac5{2\kappa}\beta_x+o(\beta)$.  Fix
$\zeta\in(0,1)$.  If $\kappa\le18(1+\zeta)$, then
$\tfrac5{2\kappa}\ge\tfrac5{36}(1-\zeta)$; otherwise
$\tfrac1{24}-G(\kappa)\ge\tfrac1{24}-G(18(1+\zeta))>0$, a fixed quantity,
which exceeds $\tfrac{41}{36}\beta_x$ for $r$ close to $1$.  Either way
$\tfrac1{24}-\tau(\{x\})\ge\tfrac5{36}\beta_x+\min(\tfrac1{24}-G(\kappa),\beta_x)
-\tfrac5{36}\zeta\beta_x-o(\beta)$, and $\zeta$ is arbitrary.

\emph{Proof of \eqref{eq:jumpup}.}  Now $L(\max(w,0))=\tfrac1{24}-G(\kappa)\le\beta_x$,
so by Theorem~\ref{thm:familylimit}(d) $\max(w,0)\to0$.  By the bound
$E\le s_x+4-1/(r\varepsilon)$ of the proof of
Theorem~\ref{thm:familylimit}, with
$18\cdot3^{-x}s_x=(1-36A+90B-56C)/(1-\tfrac{10}3\beta_x+\tfrac73\gamma_x)\le1+4\beta_x$
and $3^{-x}/(r\varepsilon)=1/(r\kappa)\ge(1-\varepsilon)/\kappa$, we get
$3^{-x}E\le w+\tfrac14\beta_x$ for $r$ close to $1$.  Let
$w'=\max(w,0)+\tfrac12\beta_x$, take $c'\ge3$ with $M_{c'+1}<w'\le M_{c'}$
and $\varphi=(w'-M_{c'+1})/(2g_{c'})$, and put $\varphi_d=1$ for $d>x$ and
for $c'<d<x$, and $\varphi_{c'}=\varphi$.  As $w'\to0$, $c'\to\infty$; as
$M_{c'}\ge w'\ge\tfrac12\beta_x$ and $M_{c'}\le4\cdot3^{-c'}$,
$3^{c'}\le8/\beta_x$, so $c'<x$.  By \eqref{eq:Hbelow},
$|H_d|\ge3^xg_d(1-\xi)-2$ for $2\le d<x$, with
$\xi=1-y^x/(1+\tfrac73\gamma_x)=o(\beta)$, so
\[
  2\sum_d\varphi_d|H_d|-2\sum_{d>x}H_d
  \ge3^x(1-\xi)(w'-M_x)-4x\ge3^x(w'-\tfrac14\beta_x)\ge E
\]
for $r$ close to $1$, that is, $2\sum_d\varphi_d|H_d|\ge F$.  By
Theorem~\ref{thm:familyzero}(b),
$\tau(\{x\})\ge t(x)-2\sum_{d>c'}|U_d|-2\varphi|U_{c'}|$, with
$U=U^{(x)}$.  Its weights differ from $(0,-\tfrac52,\tfrac72)$ by at most
$15\beta_x$ each ($D\ge\tfrac{17}{10}A$), so
$|U_d|\le|U^\infty_d|+45\beta_x3^{-d}$ and the loss is at most
$L(w')+135\beta_x3^{-c'}=L(w')+o(\beta)$.  Finally
$L(w')-L(\max(w,0))\le\tfrac12\beta_x\cdot5(\tfrac35)^{c'}=o(\beta)$, the
slopes of $L$ on $[\max(w,0),w']$ being $\ell(c'')\le5(\tfrac35)^{c''}$ with
$c''\ge c'$.  So
$\tau(\{x\})\ge t(x)-(\tfrac1{24}-G(\kappa))-o(\beta)$, which is
\eqref{eq:jumpup}.

\emph{The other placements.}  Placements $x\le3$ cost at most
$\tfrac1{48}$, and $4\le x<x^*$ at most $t(x)$, with
$\tfrac1{24}-t(x)=\tfrac5{36}\beta_xR(x)$, at least
$\tfrac5{36}\cdot\tfrac53R(x^*-1)\beta_{x^*}$ for $x=x^*-1$ and
$\tfrac5{36}\cdot\tfrac{25}9\cdot\tfrac9{20}\beta_{x^*}$ below; both exceed
$\tfrac5{36}\cdot\tfrac98\beta_{x^*}$ for $r$ close to $1$.  Placements
$x\ge x^*+2$ have $3^x\varepsilon\ge54$ and cost at most
$\tfrac1{24}-\eta$ for a fixed $\eta>0$ by step (i) of the proof of
Theorem~\ref{thm:familyosc}.

(a) At $x^*$, $\kappa<18$ and $G(\kappa)=\tfrac1{24}$, so by
\eqref{eq:jumplow} and \eqref{eq:jumpup}
$\tfrac1{24}-\tau(\{x^*\})=\tfrac5{36}\beta_{x^*}+o(\beta)$; every placement
$x\ne x^*,x^*+1$ falls short of $\tau(\{x^*\})$ by a multiple of
$\beta_{x^*}$, so every worst placement is $\{x^*\}$ or $\{x^*+1\}$ and
$\tfrac1{24}-T=\min(\tfrac1{24}-\tau(\{x^*\}),\tfrac1{24}-\tau(\{x^*+1\}))$.
At $x^*+1$ put $g=\tfrac1{24}-G(3^{x^*+1}\varepsilon)$.  If
$g\le\beta_{x^*+1}$, \eqref{eq:jumplow} and \eqref{eq:jumpup} give
$\tfrac1{24}-\tau(\{x^*+1\})=\tfrac5{36}\beta_{x^*+1}+g+o(\beta)$.  If
$g>\beta_{x^*+1}$, \eqref{eq:jumplow} gives
$\tfrac1{24}-\tau(\{x^*+1\})\ge\tfrac{41}{36}\cdot\tfrac35\beta_{x^*}-o(\beta)$,
more than $\tfrac5{36}\beta_{x^*}$ by a multiple of $\beta_{x^*}$, as is
$\tfrac5{36}\beta_{x^*+1}+g$.  Either way
\[
  \tfrac1{24}-T=\min\bigl(\tfrac5{36}\beta_{x^*},\ \tfrac5{36}\beta_{x^*+1}+g\bigr)+o(\beta).
\]
Multiply by $576\varepsilon^{-\alpha}$: by \eqref{eq:oscscale}
$80\beta_{x^*}\varepsilon^{-\alpha}=\Phi(\omega)$ and
$80\beta_{x^*+1}\varepsilon^{-\alpha}=\Phi(\omega-1)$, and
$576\varepsilon^{-\alpha}g=\Gamma(r)$.  As in the proof of
Theorem~\ref{thm:familyosc}, $\Delta(r)-576(\tfrac1{24}-T)\varepsilon^{-\alpha}\to0$,
which gives (a).

(b) Let $v=1-\omega$ and $\eta=3^v-1$, so $3^{x^*+1}\varepsilon=18(1+\eta)$.
If $v\le\tfrac12$, then $v\le\eta\le2v\le1$, and
Theorem~\ref{thm:familylimit}(d) gives
$576C_1\nu^{1+\alpha}\le\Gamma(r)\le2304C_2\nu^{1+\alpha}$, as
$\varepsilon^{-\alpha}v^{1+\alpha}=\nu^{1+\alpha}$; if $v>\tfrac12$, then
$g\ge\tfrac1{24}-G(18\sqrt3)>0$ and $\Gamma(r)\to\infty$.  Choose $a_2$ with
$576C_1a_2^{1+\alpha}\ge\Phi(1)+1$ and $a_1$ with
$2304C_2a_1^{1+\alpha}\le\tfrac1{10}\Phi(0)$.  If $\nu\ge a_2$ (and $r$ is close
to $1$), then $g>\beta_{x^*+1}$ or $\Gamma\ge\Phi(1)+1$, and in both cases
the estimates above give $\tfrac1{24}-\tau(\{x^*+1\})$ at least
$\tfrac1{24}-\tau(\{x^*\})$ plus a fixed multiple of
$\varepsilon^\alpha$, since $576\varepsilon^{-\alpha}\beta_{x^*+1}=\tfrac{36}5\Phi(\omega-1)\ge90>\Phi(1)+1$;
so $\{x^*\}$ is the only worst placement.  If $\nu\le a_1$, then
$g\le\tfrac1{5760}\Phi(0)\varepsilon^\alpha<\beta_{x^*+1}$, and
$576\varepsilon^{-\alpha}(\tfrac1{24}-\tau(\{x^*+1\}))\le\Phi(\omega-1)+\tfrac1{10}\Phi(0)+o(1)
\le\Phi(\omega)-\tfrac3{10}\Phi(0)+o(1)$, while
$576\varepsilon^{-\alpha}(\tfrac1{24}-\tau(\{x^*\}))=\Phi(\omega)+o(1)$; so
$\{x^*+1\}$ is the only worst placement.

(c) If $\nu\to\infty$ then $\Gamma\to\infty$ by the bounds in (b), and if
$\nu\to0$ then $\Gamma\to0$; (a) gives the two limits, as
$\Phi(\omega-1)<\Phi(\omega)$.  On $0<\omega\le1-\delta$,
$\nu\ge\delta\varepsilon^{-\alpha/(1+\alpha)}\to\infty$.
\end{proof}

Two exact examples show the switch.  At $r=\tfrac{101}{100}$,
$3^6\varepsilon=7.29$ and $3^7\varepsilon=21.87$, so $x^*=6$, $\omega=0.82\ldots$
and $\nu(r)=0.76\ldots$.  By Theorem~\ref{thm:familyzero}(c) the placement
$6$ costs the tail $0.035595\ldots$ of $\{1,6,41\}$, and the placement $7$
the tail $0.035944\ldots$ of $\{1,4,7\}$; the placements $4,5$ cost at most
$t(x)<t(6)$, those $x\le3$ at most $\tfrac1{48}$, and
Lemma~\ref{lem:eighteen}(b) with $c=3$ at $x=8$ bounds every $x\ge8$ by
$0.02761$.  So $T=\tau(\{7\})$ and $\inf_Qb_2(QP_r)=27.8209\ldots$: the
worst placement is $x^*+1$, as the floating-point search after
Theorem~\ref{thm:familylow} suggested.  At
$r=\tfrac{9842}{9841}$, $3^{10}\varepsilon=6.0003\ldots$ and
$3^{11}\varepsilon=18.0009\ldots$, so $x^*=10$ and $1-\omega=4.6\ldots\cdot10^{-5}$,
deep in the window.  The placement $11$ costs the tail $0.0411505\ldots$ of
$\{1,7,11\}$, more than $t(10)=0.0408515\ldots$ and than $t(x)$ for
$4\le x\le9$, while Lemma~\ref{lem:eighteen}(b) with $c=3$ at $x=12$ bounds
every $x\ge12$ by $0.0288$.  So $T=\tau(\{11\})$,
$\inf_Qb_2(QP_r)=24.3009\ldots$, and $\Delta(r)=21.63\ldots$, near
$\Phi(\omega-1)=20.86\ldots$ rather than $\Phi(\omega)=34.77\ldots$.  Both are
checked exactly in \texttt{test\_family\_jumps}.

Theorem~\ref{thm:familyosc}(a) holds on an explicit range.

\begin{proposition}[An explicit range]\label{prop:familyrdelta}
Let $0<\delta\le\tfrac12$.  If $1<r\le1+(\delta/10)^4$ and
$\delta\le\omega\le1-\delta$, in the notation of
Theorem~\ref{thm:familyosc}, then $\{x^*\}$ is the only worst placement,
$T=\tau(\{x^*\})$ is the tail of the full certificate with zero set
$\{1,x^*,z\}$ for some $z\ge x^*+3$, and
$1/t(x^*)\le\inf_Qb_2(QP_r)\le1/(t(x^*)-14\cdot3^{-x^*-3})$.
\end{proposition}

\begin{proof}
Write $x=x^*$, $\kappa=3^x\varepsilon=18\cdot3^{-\omega}$ and
$\beta=(\tfrac35)^x$, and let $\varepsilon\le(\delta/10)^4\le20^{-4}$.  We
use $0.4649<\alpha<0.465$, so $1/\alpha<2.151$, $\log_3\tfrac73>0.77$ and
$\tfrac12\log_35>0.73$.  Each bound below of a power of $\varepsilon$ by a
multiple of $\delta$ is used at $\varepsilon=(\delta/10)^4$, where its two
sides have a ratio increasing in $\delta$ (a positive power of $\delta$,
times a logarithm for the bound on $x\varepsilon$); so it suffices to check
it at $\delta=\tfrac12$.  As
$3^{x+1}\varepsilon\ge18$, $3^x\ge6/\varepsilon\ge960000$, so $x\ge13$ and
$3^{-x}\le\varepsilon/6$; hence $\beta=(3^{-x})^\alpha\le(\varepsilon/6)^\alpha
\le\delta/250$, $(\tfrac37)^x\le(\varepsilon/6)^{0.77}$ and
$5^{-x/2}\le(\varepsilon/6)^{0.73}$.  As $3^x<18/\varepsilon$,
$x\varepsilon<\varepsilon\log_3(18/\varepsilon)\le\delta/2000$, the middle
term increasing in $\varepsilon$ and at most $\delta/2000$ at
$\varepsilon=(\delta/10)^4$.

\emph{The placement $x$.}  With $y^x\ge1-x\varepsilon$,
$y^x(1+y)-y^2\ge(1-x\varepsilon)(2-\varepsilon)-1>\varepsilon>\tfrac{1-y}{24}$.
By the proof of Lemma~\ref{lem:eighteen}(c),
$18\cdot3^{-x}s_x\le(1+90\cdot5^{-x})/(1-\tfrac{10}3\beta)\le1+4\beta$ for
$x\ge13$, so $s_x\le3^{-\delta}(1+4\beta)/\varepsilon$ as
$\kappa\le18\cdot3^{-\delta}$; with $3^{-\delta}\le1-\tfrac56\delta$ on
$[0,\tfrac12]$ (convexity, $3^{-1/2}<\tfrac7{12}$) and
$4\beta+9\varepsilon<\tfrac56\delta$, this gives
$(s_x+7)r\varepsilon\le(1-\tfrac56\delta+4\beta+7\varepsilon)(1+\varepsilon)<1$.
So Lemma~\ref{lem:eighteen}(a) applies: $\tau(\{x\})$ is the tail of
$\{1,x,z\}$ with $z\ge x+3$, and $t(x)-14\cdot3^{-x-3}\le\tau(\{x\})\le t(x)$.
For $x\ge11$ the closed form of $R(x)$ in Lemma~\ref{lem:eighteen}(c) has
numerator in $[1-\tfrac75(\tfrac57)^{11}-9\cdot3^{-11}-7^{-10},1+\tfrac{84}5\,3^{-11}]$
and denominator in $[1-\tfrac{10}3(\tfrac35)^{11},1+\tfrac73(\tfrac37)^{11}]$,
so $0.96\le R(x)\le1.02$, and
$\tfrac1{24}-\tau(\{x\})\le\tfrac5{36}\cdot1.02\,\beta+14\cdot3^{-x-3}\le0.143\beta$.

\emph{Placements below $x$.}  Those $x'\le3$ cost at most $\tfrac1{48}$, and
$4\le x'<x$ at most $t(x')$, with
$\tfrac1{24}-t(x-1)=\tfrac5{36}\cdot\tfrac53R(x-1)\beta\ge0.22\beta$ and
$\tfrac1{24}-t(x')\ge\tfrac1{16}(\tfrac35)^{x'}\ge\tfrac1{16}\cdot\tfrac{25}9\beta>0.17\beta$
for $x'\le x-2$ (Lemma~\ref{lem:threenode}(c)).  All are below
$\tau(\{x\})$.

\emph{Placements above $x$.}  Let $c$ be the least integer $c\ge4$ with
$18\cdot3^{-c}\le\delta/2$, so $3^c<108/\delta<3^x$, and put
$\kappa'=3^{x+1}\varepsilon=18\cdot3^{1-\omega}\ge18\cdot3^\delta\ge18(1+\tfrac{21}{20}\delta)$.
By Lemma~\ref{lem:eighteen}(b) and its monotonicity every $x'>x$ costs at
most $t(c)+\ell(c)\Pi/\kappa'$, where
\[
  \Pi=\Bigl(1+\frac{\lambda_7(c)}{\ell(c)}\bigl(\tfrac37\bigr)^{x+1}\Bigr)\,
  r^x\,\frac{35}{24\,h_{x-1}(\tfrac15,\tfrac17,1)} .
\]
Here, with $A,B,C$ at $c$, $\lambda_7(c)/\ell(c)=\tfrac{7(3A-5B)}{3(5B-7C)}\le\tfrac{7A}{5B-7C}\le2.3(\tfrac53)^c
<2.3(108/\delta)^\alpha$, and $(\tfrac37)^x\le(\varepsilon/6)^{0.77}$,
so the first factor is at most $1+\delta/100$; $r^x\le e^{x\varepsilon}\le1+\delta/1000$;
and $h_{x-1}(\tfrac15,\tfrac17,1)\ge\tfrac{35}{24}(1-5^{-x/2})(1-7^{-x/2})$,
the terms with both exponents at most $\tfrac{x-1}2$, so the last factor is
at most $1+3\cdot5^{-x/2}\le1+\delta/100$.  So $\Pi\le1+\delta/30$.  By the
proof of Theorem~\ref{thm:familylimit}(d), $K_c\le18(1+18\cdot3^{-c})\le18(1+\tfrac\delta2)$,
so
\[
  \frac1{24}-t(c)-\frac{\ell(c)\Pi}{\kappa'}
  =\Bigl(\frac1{24}-t(c)\Bigr)\Bigl(1-\frac{K_c\Pi}{\kappa'}\Bigr)
  \ge\frac1{16}\Bigl(\frac35\Bigr)^c\Bigl(1-\frac{(1+\tfrac\delta2)(1+\tfrac\delta{30})}{1+\tfrac{21}{20}\delta}\Bigr)
  \ge0.02\,\delta\Bigl(\frac35\Bigr)^c .
\]
This exceeds $0.143\beta\ge\tfrac1{24}-\tau(\{x\})$ once
$(\tfrac35)^{x-c}\le\delta/7.2$, that is, $x-c\ge\log_{5/3}(7.2/\delta)$,
which holds as $c<\log_3(108/\delta)$, $x\ge\log_3(6/\varepsilon)$ and
$(\delta/10)^4\le\tfrac\delta{18}(\tfrac\delta{7.2})^{1/\alpha}$ for
$0<\delta\le\tfrac12$ ($1/\alpha=2.150\ldots$).  So every placement above
$x$ costs less than $\tau(\{x\})$.

Hence $\{x\}$ is the only worst placement, $T=\tau(\{x\})$, and
Theorem~\ref{thm:value} gives the bounds on the infimum.
\end{proof}

So the escaping placement stays extremal at some roots below $2$, here down
to $r_c$, and the failure below $r_c$ starts at the placement $\{4\}$ just
above $Z^*$, where the conclusion of Lemma~\ref{lem:delete} fails.  The cutoff cannot be
lowered uniformly, however.

\begin{proposition}[The cutoff is sharp]\label{prop:cutoffsharp}
Let $1<r<2$, $P_r=(x-r)(x-5)(x-7)$ and $u=1$.  Then $\tau^*=\tfrac1{24}$
and
\[
  \inf_Qb_2(QP_r)\le\max(r+11,\,24(r-1))<24=\frac1{\tau^*},
\]
so $T>\tau^*$.  At $r=2$ the infimum is $24$.  Hence for every
$r_1\in(1,2)$ the conclusion of Theorem~\ref{thm:extremal} fails at some
root set with smallest root $r_1$.
\end{proposition}

\begin{proof}
The full certificate on $(\tfrac15,\tfrac17)$ with zero set $\{1\}$ has
tail $\tfrac1{(5-1)(7-1)}=\tfrac1{24}$ by~\cite[Lemma~2.1]{coefficient-mass}, and $(x-5)(x-7)$ has
partial sums $1,-11,24$, so \eqref{eq:dualpair} with $M=1$ gives
$\tau^*=\tfrac1{24}$.  The partial sums of $P_r$ are
$1,-(r+11),11r+24,-24(r-1)$, constant from $j=3$ on, so
Lemma~\ref{lem:exemptpair} with $X=\{2\}$ and $M=1$ gives the bound, which
is below $24$ for $r<2$.  At $r=2$, Theorem~\ref{thm:extremal} gives
$1/\tau^*=24$.
\end{proof}

\section{Scope and open questions}\label{sec:scope}

\paragraph{What is left open here.}  For roots at least $2$ the infimum is
$1/\tau^*$ (Theorem~\ref{thm:extremal}); below $2$ it is $1/T$ with
$T\ge\tau^*$, and four questions about $T$ remain open.
\begin{itemize}
\item An a priori bound on the zero set minimising a single $\tau(\{x\})$
at general roots, which would make the case $u=1$ a finite computation
outright.  Along $(r,3,5,7)$ Theorem~\ref{thm:familyzero} gives the zero
set exactly, $\{1,x,z\}$ or $\{1,z,x\}$ with $z$ a half-mass point; its
proof uses that the minimiser at the three fast nodes is $\{1,x\}$ with a
witness whose corrections at $1$ and $x$ keep their signs.
\item An analogue of the threshold $N$ of Theorem~\ref{thm:halfmass} for
$u\ge2$, where only the qualitative Proposition~\ref{prop:attainT} is
proved.
\item A description of the root sets below $2$ at which the escaping
placement stays extremal, which depends on more than $r_1$
(Theorem~\ref{thm:family} and Proposition~\ref{prop:cutoffsharp}).
\item The infimum along $(r,3,5,7)$ for $1<r<\rho_1=1.0745\ldots$ as a
function of $r$.  Theorem~\ref{thm:familylow} gives it on
$[\rho_1,\tfrac{13}{10}]$ in seven pieces; below, every cost is explicit
at each $r$ (Theorem~\ref{thm:familyzero}), infinitely many pieces
accumulate at $r=1$ (Corollary~\ref{cor:familyone}), and near $r=1$ the
infimum is described uniformly: $(\inf_Qb_2(QP_r)-24)(r-1)^{-\alpha}$ is
$\min(\Phi(\omega),\Phi(\omega-1)+\Gamma(r))+o(1)$, with $\Gamma$ given
by the limit profile $G$ (Theorems~\ref{thm:familylimit}
and~\ref{thm:familyjump}), and the conclusions of
Theorem~\ref{thm:familyosc}(a) hold for $r\le1+(\delta/10)^4$
(Proposition~\ref{prop:familyrdelta}).  Open: the order of the pieces on
$(1,\rho_1)$, that is, whether the worst placement is nonincreasing in $r$
and whether each boundary equation after Theorem~\ref{thm:familyzero}
has a single root in $r$; and explicit constants $a_1,a_2$ and a rate of
convergence in Theorem~\ref{thm:familyjump}.  Since
$(\tfrac1{24}-G(18(1+\eta)))\eta^{-1-\alpha}$ oscillates
(Theorem~\ref{thm:familylimit}(e)), the window itself is log-periodic:
whether $\Delta$ is monotone in $r$ across each window, and where exactly
the worst placement switches, in the scale $\nu$, are open.
\end{itemize}

\paragraph{Relation to the literature.}  The nearest work is the reduced
height of Dubickas and Jankauskas~\cite{dj}, which is our infimum at $u=0$;
its relatives, Schinzel's reduced length~\cite{schinzel} and the minimal
Euclidean norm of Filaseta, Robinson and Wheeler~\cite{frw}; and the height
of polynomials with split roots, studied by Hajdu, Tijdeman and
Varga~\cite{htv} and by Tromp~\cite{tromp}.  The bounds here concern order
statistics of the coefficient magnitudes of multiples, which at $u=0$ reduce
to the height; they are uniform over the degree
of $F$ and over all additional real or complex roots.  In the opposite
direction, the degree of an integer polynomial of height $H\ge2$ all of whose
roots are nonzero rationals is at most a constant times $\log H$~\cite{htv},
and the minimal height of a degree-$n$ polynomial over a fixed number
field $\mathbb K$ with all roots in $\mathbb K^\times$ grows exponentially in $n$~\cite{tromp}.  At
$u=0$ the infimum of $b_1(QP)$ over monic $Q$ is the reduced height of $P$,
the $\ell^\infty$ member of a family
of extremal problems over normalized multiples: reduced length minimizes
$\ell^1$~\cite{schinzel}, and the integer-polynomial $\ell^2$ analogue, the
minimal Euclidean norm of an algebraic number, is effectively
computable~\cite{frw}.  For $u\ge1$ the infimum of $b_{u+1}$, which ignores
the $u$ largest nonleading coefficient magnitudes, is an order-statistic
variant of the reduced height.  Dubickas and Jankauskas evaluate the tail
of the certificate with a consecutive zero set under a sign
condition~\cite[p.~333]{dj}, which
\cite[Lemma~2.1]{coefficient-mass} removes, and the case $u=0$ of the
construction of Section~\ref{sec:criterion} is their
Corollary~7~\cite{dj}.

\paragraph{Data availability.}
No datasets were generated or analyzed in this study.  The exact
computations quoted in the paper are in \texttt{tests/test\_attainment.py},
with helpers in \texttt{tests/sweep.py}, of the repository
\href{https://github.com/bangyen/coefficient-mass}{bangyen/coefficient-mass},
and run with \texttt{uv run pytest tests/test\_attainment.py}.

\paragraph{Computer assistance.}
Six computations are computer-assisted: steps of the proofs of
Theorems~\ref{thm:elevenvalue}, \ref{thm:family} and~\ref{thm:familylow},
the example at $r=\tfrac{5701}{5700}$ after Theorem~\ref{thm:familyosc},
the two examples at $r=\tfrac{101}{100}$ and $\tfrac{9842}{9841}$ after
Theorem~\ref{thm:familyjump}, and the failure of
\eqref{eq:runchord} above the cutoff.  In
Theorem~\ref{thm:elevenvalue}, the sum $H$, its tail, the threshold $N=11$
and the tails of the seven full certificates for the placements
$x=2,4,6,7,8,9,10$, with their comparisons with $\theta$, are computed over
$\mathbb Q$ in \texttt{test\_value\_eleven\_tenths}.  In
Theorem~\ref{thm:family}, the tails of the full certificates with zero sets
$\{1,3,4\}$, $\{1,2,5\}$ and $\{1,3,5\}$ and the closed forms of
$2\sum_{d\le n}|H_d|-\Tail(H)$ for $n=5,6$ are identities in $\mathbb Q(r)$,
verified with $r$ an indeterminate in \texttt{test\_family\_threshold} and
\texttt{test\_family\_window};
Appendix~\ref{sec:familydata} lists the solutions of the linear systems, from
which they can also be checked by hand.  In Theorem~\ref{thm:familylow}, the
tails $\theta_z$, the witnesses and their identities, the ties between
pieces, the closed forms of $2\sum_{d\le n}|H_d|-\Tail(H)$ for
$n=6,\ldots,13$ and the tails of the full certificates $\{1,3,x\}$
($6\le x\le12$), $\{1,2,13\}$ and $\{1,2,10\}$ are computed over
$\mathbb Q(r)$ in \texttt{test\_family\_below}, and every sign of a rational
function on an interval that its proof uses is decided by Sturm sequences
over $\mathbb Q$ between rational endpoints; Appendix~\ref{sec:familydata}
lists the tails, the breakpoint polynomials and one witness.  At
$r=\tfrac{5701}{5700}$, \texttt{test\_family\_oscillation} solves the
$4\times4$ system for the witness of the zero set $\{1,10,1335\}$ and the
certificate itself over $\mathbb Q$, forms its tail from the sign pattern
as a combination of geometric sums, and compares it with $t(x)$ for
$4\le x\le9$ and with the bound of Lemma~\ref{lem:eighteen}(b) for
$11\le x\le20$.  At $r=\tfrac{101}{100}$ and $\tfrac{9842}{9841}$,
\texttt{test\_family\_jumps} finds the half-mass points of
Theorem~\ref{thm:familyzero}(c) over $\mathbb Q$, checks that the tails of
the resulting full certificates $\{1,6,41\}$, $\{1,4,7\}$ and
$\{1,7,11\}$ equal the costs given there, and compares them with $t(x)$
and with the bound of Lemma~\ref{lem:eighteen}(b), decreasing in $x$, at
$x=8$ and $x=12$.  The statement after
Proposition~\ref{prop:runchord} that \eqref{eq:runchord} can fail above the
cutoff rests on the two rationals computed in
\texttt{test\_run\_chord\_control\_above\_cutoff}.  All arithmetic is exact:
rationals are Python \texttt{fractions.Fraction}, and $\mathbb Q(r)$ is
represented by quotients of polynomials with \texttt{Fraction}
coefficients, a small class defined in the test file; no floating point and
no computer algebra system is used, and every decimal in the paper is a
truncation of an exact rational, of a zero of an explicit polynomial, or
of an explicit constant such as $\alpha$ or $80\cdot18^{-\alpha}$.  The
floating-point searches reported after Theorems~\ref{thm:familylow}
and~\ref{thm:familyosc} and in Section~\ref{sec:scope} are exploration,
used in no proof and not part of the tests.  The other computations only check steps
proved by hand: the example after Theorem~\ref{thm:extremal}
(\texttt{test\_extremal\_example\_2357}); the multiple $F$ of
Proposition~\ref{prop:belowtwo} and the certificate with zero set
$\{1,3,5\}$ after it (\texttt{test\_below\_two}); the partial sums of $PM$
in Theorem~\ref{thm:elevenvalue} and the comparisons in the remark after it
(\texttt{test\_value\_eleven\_tenths}); the partial sums, quadratics and
cubics of the proof of Theorem~\ref{thm:family}
(\texttt{test\_family\_threshold} and \texttt{test\_family\_window});
Proposition~\ref{prop:cutoffsharp} (\texttt{test\_cutoff\_sharp}); and,
in \texttt{test\_family\_near\_one}, Lemma~\ref{lem:threenode} and
Theorem~\ref{thm:familyone}, where the closed forms are checked as
identities of rational functions of $A,B,C$ (multiplied by $D$, each side
is a polynomial of degree at most $2$ in each variable, so agreement on a
$4\times4\times4$ grid proves it), together with the rational inequalities
behind Lemma~\ref{lem:threenode}(c) and the witness at
$r=\tfrac{101}{100}$.  The estimates with logarithms in the proof of
Theorem~\ref{thm:familyone}, and those with powers $\alpha$ and
$1/\alpha$ in the proofs of Theorem~\ref{thm:familyjump}(b) and
Proposition~\ref{prop:familyrdelta}, reduce to comparisons of integer
powers, with $e$ bounded by its series and $\alpha$ bracketed as
$0.4649<\alpha<0.465$; \texttt{test\_family\_hand\_estimates} checks them.  In
\texttt{test\_family\_oscillation} the solution for $p,q,\theta'$, the
constants $15,6,14$, the identity at the start of the path in
Lemma~\ref{lem:eighteen}(a), the closed forms in (c) with $K_c>18$ for
$c\le60$, and the sum $g$ of (b) are checked in the same way; the limits
in the proof of Theorem~\ref{thm:familyosc} are by hand.  In
\texttt{test\_family\_zero\_sets}, the weights of $V$ and $V'$, the
identity \eqref{eq:threebasis}, the identities of
Theorem~\ref{thm:familyzero}(a) and the polynomial $p$ of its proof are
checked for $x\le15$, and at several $r$ and $x$ the witness of part (c)
is built and checked to vanish at the four nodes with entries in
$[-1,1]$, against all zero sets in $\{1,\ldots,16\}$; the costs at
$r=\tfrac54$ and $\tfrac{11}{10}$ quoted after the theorem and the zero set
$\{1,10,1335\}$ at $r=\tfrac{5701}{5700}$ are recovered.  In
\texttt{test\_family\_jumps} the identities \eqref{eq:profileid} are
checked on a grid as above, the $k_c$ and $M_c$ for $c\le40$, and the
rational constants in the proofs of Theorems~\ref{thm:familylimit}
and~\ref{thm:familyjump} and Proposition~\ref{prop:familyrdelta} for
$c\le60$ or $x\le60$; the limits are by hand.  The file also holds tests used in no
proof: seeded random exact checks of
Proposition~\ref{prop:runchord} and of the identity in
Theorem~\ref{thm:halfmass}(b), and exact linear programs over monic
multiples of bounded degree, which confirm the value $\tfrac{1704}{31}$ at
$(2,3,5,7)$ and serve as controls that no multiple found undercuts a stated
infimum.

\paragraph{Use of AI tools.}
Large language model assistants (Claude, Anthropic) were used in drafting
and revising the text, the proofs and the verification code.  The author has checked the mathematics and is responsible
for the content.

\paragraph{Funding and competing interests.}
The author received no specific funding for this work and has no competing
interests to declare.

\appendix

\section{Data for the family \texorpdfstring{$(r,3,5,7)$}{(r,3,5,7)}}\label{sec:familydata}

This appendix lists the solutions of the linear systems in the proof of
Theorem~\ref{thm:family}, and the data of Theorem~\ref{thm:familylow}, as
rational functions of $r$, written with
\[
  p(r)=(3-r)(5-r)(7-r),\qquad q(r)=3466r^2+7455r+11025 ,
\]
so that the tails $1/\beta(r)$,
$t_2(r)=(4267r^2+5070r-5871)/\bigl(48(r-1)q(r)\bigr)$ and $t_5(r)$ there can be checked
by hand or by any computer algebra system.  Let $1<r<3$, so that
$\tfrac1r$ is the largest of the nodes
$\tfrac1r,\tfrac13,\tfrac15,\tfrac17$; then $p(r)$ and $q(r)$ are positive.  Each sum below has the form
$v_d=\lambda_1r^{-d}+\lambda_33^{-d}+\lambda_55^{-d}+\lambda_77^{-d}$, and
past its largest prescribed zero $z$ its values sum to
\[
  \sum_{d>z}v_d=\frac{\lambda_1}{r^{z}(r-1)}+\frac{\lambda_3}{2\cdot3^{z}}
  +\frac{\lambda_5}{4\cdot5^{z}}+\frac{\lambda_7}{6\cdot7^{z}} .
\]
Each weight vector is verified by substituting it into the defining
conditions; the values listed are then direct evaluations.  The signs are
those of the zero bound: a full certificate is positive at $0$ and changes
sign exactly at its zeros, and $H$ vanishes exactly at $0,1,3$, with
$H_d>0$ for $d>3$.  The escaping minimiser $V$, with zero set $Z^*=\{1,3\}$,
is displayed after Theorem~\ref{thm:extremal}.

\paragraph{The certificate with zero set $\{1,3,4\}$.}  Its weights are
\[
  \lambda_1=-\frac{15r^4}{p(r)(15r+71)},\quad
  \lambda_3=\frac{81(r+12)}{8(3-r)(15r+71)},\quad
  \lambda_5=-\frac{625(r+10)}{4(5-r)(15r+71)},\quad
  \lambda_7=\frac{2401(r+8)}{8(7-r)(15r+71)},
\]
and $v_0=1$, $v_1=v_3=v_4=0$,
\[
  v_2=-\frac1{15r+71},\qquad
  \sum_{d\ge5}v_d=-\frac{r+14}{48(r-1)(15r+71)},
\]
with $v_2<0$ and $v_d<0$ for $d\ge5$, so
\[
  \Tail(v)=|v_2|-\sum_{d\ge5}v_d=\frac{49r-34}{48(r-1)(15r+71)}=\frac1{\beta(r)}.
\]

\paragraph{The certificate with zero set $\{1,2,5\}$.}  Its weights are
\begin{gather*}
  \lambda_1=-\frac{3466r^5}{p(r)q(r)},\qquad
  \lambda_3=\frac{243(109r^2+420r+1225)}{8(3-r)q(r)},\\
  \lambda_5=-\frac{3125(79r^2+210r+441)}{4(5-r)q(r)},\qquad
  \lambda_7=\frac{16807(49r^2+120r+225)}{8(7-r)q(r)},
\end{gather*}
and $v_0=1$, $v_1=v_2=v_5=0$,
\[
  v_3=\frac{71r+105}{q(r)},\qquad
  v_4=\frac{15r+71}{q(r)},\qquad
  \sum_{d\ge6}v_d=-\frac{139r^2+750r+2577}{48(r-1)q(r)},
\]
with $v_3,v_4>0$ and $v_d<0$ for $d\ge6$, so
\[
  \Tail(v)=v_3+v_4-\sum_{d\ge6}v_d
  =\frac{4267r^2+5070r-5871}{48(r-1)q(r)}=t_2(r).
\]

\paragraph{The certificate with zero set $\{1,3,5\}$.}  With
$q_5(r)=960r^2+5041r+7455$, its weights are
\begin{gather*}
  \lambda_1=-\frac{960r^5}{p(r)q_5(r)},\qquad
  \lambda_3=\frac{729(r+5)(r+7)}{2(3-r)q_5(r)},\\
  \lambda_5=-\frac{15625(r+3)(r+7)}{2(5-r)q_5(r)},\qquad
  \lambda_7=\frac{16807(r+3)(r+5)}{(7-r)q_5(r)},
\end{gather*}
and $v_0=1$, $v_1=v_3=v_5=0$,
\[
  v_2=-\frac{71r+105}{q_5(r)},\qquad
  v_4=\frac{r+15}{q_5(r)},\qquad
  \sum_{d\ge6}v_d=-\frac{7r^2+98r+375}{24(r-1)q_5(r)},
\]
with $v_2<0$, $v_4>0$ and $v_d<0$ for $d\ge6$, so
\[
  \Tail(v)=-v_2+v_4-\sum_{d\ge6}v_d
  =\frac{5(347r^2+250r-501)}{24(r-1)q_5(r)}=t_5(r).
\]

\paragraph{The sum $H$ of Theorem~\ref{thm:halfmass}.}  With weight $1$ on
$\tfrac1r$, $H_0=0$ and $H|_{Z^*}=0$, its weights are
\[
  \lambda_1=1,\quad
  \lambda_3=-\frac{27(5-r)(7-r)(12r+35)}{568r^3},\quad
  \lambda_5=\frac{125(3-r)(7-r)(10r+21)}{284r^3},\quad
  \lambda_7=-\frac{343(3-r)(5-r)(8r+15)}{568r^3},
\]
and $H_0=H_1=H_3=0$,
\begin{gather*}
  H_2=-\frac{p(r)}{71r^3},\qquad
  H_4=\frac{p(r)(15r+71)}{7455r^4},\qquad
  H_5=\frac{p(r)(960r^2+5041r+7455)}{782775r^5},\\
  \sum_{d\ge4}H_d=\frac{p(r)(14r+57)}{3408r^3(r-1)},\qquad
  \sum_{d\ge6}H_d=\frac{p(r)(10890r^3+50099r^2+80656r+119280)}{12524400r^5(r-1)} .
\end{gather*}
So $\Tail(H)=-H_2+\sum_{d\ge4}H_d$ and $\sum_{d\le5}|H_d|=-H_2+H_4+H_5$,
which give the identity for $2\sum_{d\le5}|H_d|-\Tail(H)$ in the proof of
Theorem~\ref{thm:family}.  For the identity with $d\le6$ there,
\[
  H_6=\frac{p(r)(44535r^3+246086r^2+529305r+782775)}{82191375r^6},\qquad
  \sum_{d\ge7}H_d=\frac{p(r)\,e_7(r)}{1315062000r^6(r-1)},
\]
with $e_7(r)=430890r^4+2035579r^3+3937376r^2+8468880r+12524400$, and
$\sum_{d\le6}|H_d|=-H_2+H_4+H_5+H_6$.

\paragraph{The certificates of Theorem~\ref{thm:familylow}.}  With
\begin{align*}
  Q_5&=77r^2+480r+1733,\qquad Q_6=9359r^3+60705r^2+246086r+363930,\\
  Q_7&=421939r^4+2798055r^3+12013156r^2+25839030r+38212650,\\
  Q_8&=8457547r^5+56817765r^4+251780638r^3+630690690r^2\\
  &\qquad+1356549075r+2006164125,
\end{align*}
the tails of the full certificates with zero sets $\{1,4,z\}$ are
\[
  \theta_5=\frac{n_5}{96(r-1)Q_5},\qquad
  \theta_6=\frac{n_6}{48(r-1)Q_6},\qquad
  \theta_7=\frac{n_7}{48(r-1)Q_7},\qquad
  \theta_8=\frac{n_8}{96(r-1)Q_8},
\]
where
\begin{align*}
  n_5&=769r^2+3374r-3989,\qquad n_6=49550r^3+244468r^2+148970r-433629,\\
  n_7&=2308315r^4+12063338r^3+17059705r^2+15420090r-46429509,\\
  n_8&=94332044r^5+508724104r^4+933286436r^3+1788992289r^2\\
  &\qquad+1607184810r-4915604589 .
\end{align*}
For $\{1,4,5\}$ the weights are
\begin{gather*}
  \lambda_1=-\frac{77r^5}{p(r)Q_5},\qquad
  \lambda_3=\frac{243(r^2+12r+109)}{16(3-r)Q_5},\\
  \lambda_5=-\frac{3125(r^2+10r+79)}{8(5-r)Q_5},\qquad
  \lambda_7=\frac{16807(r^2+8r+49)}{16(7-r)Q_5},
\end{gather*}
and $v_0=1$, $v_1=v_4=v_5=0$,
\[
  v_2=-\frac{15r+71}{2Q_5},\qquad v_3=-\frac{r+15}{2Q_5},\qquad
  \sum_{d\ge6}v_d=-\frac{r^2+14r+139}{96(r-1)Q_5},
\]
all negative, so $\Tail(v)=-v_2-v_3-\sum_{d\ge6}v_d=\theta_5(r)$.  Its witness
in \eqref{eq:witness} has
\begin{gather*}
  s_1=\frac{769r^3+10766r^2+42091r-52276}{96(r-1)Q_5},\qquad
  s_4=\frac{211252r^3+380157r^2+254562r-991257}{96(r-1)Q_5},\\
  s_5=-\frac{142079r^3+137846r^2-345779r-55456}{32(r-1)Q_5},
\end{gather*}
and $s_d=-1$ for $d=2,3$ and $d\ge6$.  The other witnesses have larger
coefficients and are only computed.  The polynomials of the breakpoints are
\begin{align*}
  \gamma_1&=218644r^3+418845r^2+374850r-1157625,\\
  \gamma_2&=203860r^3+341469r^2+134274r-824889,\\
  \gamma'&=172587443623r^6+203696380262r^5-245023502843r^4-16113960832r^3\\
  &\qquad-40364204160r^2-86819140800r-128394504000,\\
  \gamma_3&=31863386060r^7+527299949544r^6+2216976743064r^5+5148085143336r^4\\
  &\qquad+7473795889089r^3+2333772047565r^2-9236198868099r-16544067678495,\\
  \gamma_4&=1638534863r^5+1910488822r^4-2452391683r^3-384420992r^2\\
  &\qquad-826848960r-1222804800,\\
  \gamma_5&=15476503r^4+17470982r^3-26164523r^2-7874752r-11645760,\\
  \gamma_6&=144543r^3+150742r^2-305683r-110912,\qquad
  \gamma_7=3989r^2+3570r-11025,\\
  \gamma_+&=111700r^3-50307r^2-97470r-109209,
\end{align*}
each with exactly one zero in $(1,\tfrac32)$, where it changes sign from
$-$ to $+$.  Besides the ties after Theorem~\ref{thm:familylow},
\begin{gather*}
  \theta_5-\theta_6=-\frac{(r^2+15r+154)\gamma_6}{32(r-1)Q_5Q_6},\qquad
  \theta_6-\theta_7=-\frac{(15r^3+225r^2+2310r+9359)\gamma_5}{16(r-1)Q_6Q_7},\\
  \theta_7-\theta_8=-\frac{7(22r^4+330r^3+3388r^2+17745r+60277)\gamma_4}{32(r-1)Q_7Q_8} .
\end{gather*}

\begin{thebibliography}{9}
\small
\bibitem{karlin-studden}
S.~Karlin and W.~J. Studden,
\newblock \emph{Tchebycheff Systems: With Applications in Analysis and Statistics},
\newblock Interscience, New York, 1966.

\bibitem{polya-szego}
G.~P\'olya and G.~Szeg\H{o},
\newblock \emph{Problems and Theorems in Analysis II},
\newblock Classics in Mathematics, Springer, 1998 (reprint of the 1976
  edition), Part Five, Chapter 1,
\newblock \href{https://doi.org/10.1007/978-3-642-61905-2}
{doi:10.1007/978-3-642-61905-2}.

\bibitem{schinzel}
A.~Schinzel,
\newblock On the reduced length of a polynomial with real coefficients,
\newblock \emph{Functiones et Approximatio Commentarii Mathematici} 35 (2006), 271--306,
\newblock \href{https://doi.org/10.7169/facm/1229442629}
{doi:10.7169/facm/1229442629}.

\bibitem{dj}
A.~Dubickas and J.~Jankauskas,
\newblock On the reduced height of a polynomial,
\newblock \emph{Publicationes Mathematicae Debrecen} 71 (2007), 325--348,
\newblock \href{https://doi.org/10.5486/PMD.2007.3728}
{doi:10.5486/PMD.2007.3728}.

\bibitem{frw}
M.~Filaseta, M.~L. Robinson and F.~S. Wheeler,
\newblock The minimal Euclidean norm of an algebraic number is effectively
computable,
\newblock \emph{Journal of Algorithms} 16 (1994), 309--333,
\newblock \href{https://doi.org/10.1006/jagm.1994.1015}
{doi:10.1006/jagm.1994.1015}.

\bibitem{htv}
L.~Hajdu, R.~Tijdeman and N.~Varga,
\newblock On polynomials with only rational roots,
\newblock \emph{Mathematika} 69 (2023), 867--878,
\newblock \href{https://doi.org/10.1112/mtk.12209}
{doi:10.1112/mtk.12209}.

\bibitem{tromp}
T.~Tromp,
\newblock On the height of polynomials that split completely over a fixed
number field,
\newblock arXiv:2509.12064 (2025),
\newblock \href{https://doi.org/10.48550/arXiv.2509.12064}
{doi:10.48550/arXiv.2509.12064}.

\bibitem{coefficient-mass}
B.~Pham,
\newblock Displaced-zero certificates and the coefficient mass of polynomial
multiples,
\newblock preprint, 2026, \url{https://github.com/bangyen/coefficient-mass};
Lemma~2.1 (Consecutive zero set), Theorem~2.2 (Displaced-zero tail bound),
Lemma~2.3 (Ratio bound for the correction sum),
Lemma~3.1 (Certificate with excluded positions), Theorem~3.2 (Coefficient order statistics),
Proposition~4.4 (Sharpness of the row $k=2$ at $(2,3)$).
\end{thebibliography}

\end{document}
